% Oscillations électriques libres (circuits LC et RLC) — Physique Tle D — Cours signé OnBuch+
% Programme officiel MINESEC. Compiler avec : tectonic P11-oscillations-electriques-libres.tex
\newcommand{\DOCMATIERE}{Physique}
\newcommand{\DOCNIVEAU}{Tle D}
\newcommand{\DOCMODULE}{Module 3 : Électricité — évolutions temporelles des circuits électriques et électroniques}
\newcommand{\DOCLECON}{P11}
\newcommand{\DOCTITRE}{Oscillations électriques libres (circuits LC et RLC)}
% OnBuch+ — préambule « cours signés » ENRICHI (compilé avec tectonic / XeLaTeX)
% Système visuel « Pop » d'OnBuch+ : fond crème (jamais blanc), bordures encre
% épaisses, ombres « dures » décalées, titres Archivo Black en capitales.
% Version MATHÉMATIQUES : graphes (pgfplots/tikz), boîtes pédagogiques
% (définition, propriété, méthode, attention, le savais-tu, exercice résolu,
% exercice, TP, bilan), page de garde et signature OnBuch+ ancrée en bas.
%
% Macros attendues dans main.tex AVANT \input{preamble} :
%   \DOCMATIERE  \DOCNIVEAU  \DOCMODULE  \DOCLECON (numéro)  \DOCTITRE
\documentclass[11pt]{article}

\usepackage{fontspec}
\usepackage[french]{babel}
\usepackage[a4paper,margin=2cm,top=3.4cm,bottom=2.8cm,headheight=1.4cm,footskip=1.3cm]{geometry}
\usepackage{amsmath,amssymb,amsfonts}
\usepackage{mathtools} % \xmapsto, \coloneqq... souvent utilisés par les modèles
\usepackage{cancel} % simplifications barrées (bilans de forces)
% \abs{z} (module, valeur absolue) et \norm{u} (norme), utilisés par les modèles
\providecommand{\abs}{}\let\abs\relax\DeclarePairedDelimiter{\abs}{\lvert}{\rvert}
\providecommand{\norm}{}\let\norm\relax\DeclarePairedDelimiter{\norm}{\lVert}{\rVert}
\usepackage[table]{xcolor} % [table] : fournit aussi \rowcolors (alternance de lignes)
\usepackage{pagecolor}
\usepackage{tikz}
\usetikzlibrary{arrows.meta,positioning,calc,shapes.geometric,shapes.misc,shapes.arrows,decorations.pathmorphing,decorations.markings,patterns,shadows,mindmap,trees,fit,backgrounds,matrix,intersections,angles,quotes}
\usepackage{pgfplots}
\usepackage[version=4]{mhchem} % équations nucléaires \ce{^{235}_{92}U}
\usepackage[european,siunitx]{circuitikz} % schémas électriques
\pgfplotsset{compat=1.18}
\usepgfplotslibrary{fillbetween,groupplots} % aires entre courbes (calcul intégral)
\usepackage{enumitem}
\usepackage{fancyhdr}
% [most] charge toutes les bibliothèques tcolorbox (listings, minted, raster,
% documentation...) et gonfle inutilement la mémoire TeX sur un document long
% (~300 boîtes) : on ne charge que ce qui est réellement utilisé.
\usepackage[skins,breakable]{tcolorbox}
\usepackage{graphicx}
\usepackage{colortbl}
\usepackage{booktabs}
\usepackage{tabularx}
\usepackage{diagbox} % case d'en-tête diagonale des tableaux à double entrée
% Colonnes extensibles centrée / gauche / droite (C, L, R), souvent utilisées par les modèles
\newcolumntype{C}{>{\centering\arraybackslash}X}
\newcolumntype{L}{>{\raggedright\arraybackslash}X}
\newcolumntype{R}{>{\raggedleft\arraybackslash}X}
\usepackage{array}
\usepackage{multicol}
\usepackage{siunitx}
% parse-numbers=false : accepte aussi \SI{2,5 \times 10^{-3}}{...}
\sisetup{locale=FR,output-decimal-marker={,},group-minimum-digits=4,parse-numbers=false}
\DeclareSIUnit\ohm{\ensuremath{\symup{\Omega}}} % U+2126 absent de la police
\DeclareSIUnit\division{div} % réglages d'oscilloscope (ms/div, V/div)
\usepackage{qrcode}
% \degree : commande fréquemment utilisée par les modèles pour un angle ou une
% température en dehors de \SI{}{} (non fournie par les paquets chargés).
\providecommand{\degree}{^{\circ}}
% \qed (fin de démonstration), utilisé par les modèles sans amsthm
\providecommand{\qed}{\hfill$\square$}
% \barre : double barre des valeurs interdites dans un tableau de variations (tkz-tab)
\providecommand{\barre}{\ensuremath{\|}}
% environnement proof (amsthm non chargé) utilisé par les modèles
\ifdefined\proof\else\newenvironment{proof}[1][Démonstration]{\par\noindent\textit{#1.}\ }{\qed\par}\fi
\usepackage{lastpage}
\usepackage{hyperref}
\hypersetup{hidelinks}

% ============================================================
% Palette « Pop » OnBuch+ — mêmes hex que la classe P (plus_ui.dart)
% ============================================================
\definecolor{poporange}{HTML}{F59321}
\definecolor{popdark}{HTML}{E07A0C}
\definecolor{popcream}{HTML}{FFF6E8}
\definecolor{popink}{HTML}{1C1714}
\definecolor{poppurple}{HTML}{7A5AE0}
\definecolor{popgreen}{HTML}{2E9E6A}
\definecolor{poppink}{HTML}{E0457B}
\definecolor{popblue}{HTML}{2F7DE1}
\definecolor{popgold}{HTML}{C98A12}
\definecolor{popmuted}{HTML}{8A7F76}
\definecolor{popline}{HTML}{EADFD2}
\definecolor{popgray}{HTML}{8A7F76}
\colorlet{popgrey}{popgray}
% Alias tolérants (le modèle nomme parfois les couleurs différemment)
\colorlet{popwhite}{white}
\colorlet{popblack}{popink}
\colorlet{popred}{poppink}
\colorlet{popyellow}{popgold}
% Teintes claires pour fonds de tableaux / zones
\colorlet{popblueL}{popblue!10!white}
\colorlet{poporangeL}{poporange!14!white}
\colorlet{popgreenL}{popgreen!12!white}
\colorlet{poppurpleL}{poppurple!12!white}
\colorlet{poppinkL}{poppink!12!white}
\colorlet{popgoldL}{popgold!14!white}

\pagecolor{popcream}

% ============================================================
% Typographie
% ============================================================
\defaultfontfeatures{Ligatures=TeX}
\setmainfont{PlusJakartaSans-Medium.ttf}[Path=../../../fonts/,BoldFont=PlusJakartaSans-Bold.ttf]
% Maths en sans-serif (Fira Math) avec chiffres et lettres droites de Plus
% Jakarta, pour que les formules se fondent dans le texte.
\usepackage[math-style=ISO,bold-style=ISO]{unicode-math}
\setmathfont{FiraMath-Regular.otf}[Path=../../../fonts/]
\setmathfont{PlusJakartaSans-Medium.ttf}[Path=../../../fonts/,range={up/num,up/{Latin,latin},\mathup}]
\usepackage{icomma} % virgule décimale sans espace en mode math
% Caractères mathématiques Unicode absents des polices de texte (Plus Jakarta,
% Archivo Black) : rendus par leur équivalent mathématique, en texte comme en
% formule — sinon ils s'affichent en carré vide (ex. « Arithmétique dans ℤ »).
\usepackage{newunicodechar}
\newunicodechar{ℝ}{\ensuremath{\mathbb{R}}}
\newunicodechar{ℤ}{\ensuremath{\mathbb{Z}}}
\newunicodechar{ℂ}{\ensuremath{\mathbb{C}}}
\newunicodechar{ℕ}{\ensuremath{\mathbb{N}}}
\newunicodechar{ℚ}{\ensuremath{\mathbb{Q}}}
\newunicodechar{𝔻}{\ensuremath{\mathbb{D}}}
\newunicodechar{α}{\ensuremath{\alpha}}
\newunicodechar{β}{\ensuremath{\beta}}
\newunicodechar{γ}{\ensuremath{\gamma}}
\newunicodechar{δ}{\ensuremath{\delta}}
\newunicodechar{ε}{\ensuremath{\varepsilon}}
\newunicodechar{θ}{\ensuremath{\theta}}
\newunicodechar{λ}{\ensuremath{\lambda}}
\newunicodechar{μ}{\ensuremath{\mu}}
\newunicodechar{σ}{\ensuremath{\sigma}}
\newunicodechar{τ}{\ensuremath{\tau}}
\newunicodechar{φ}{\ensuremath{\varphi}}
\newunicodechar{ψ}{\ensuremath{\psi}}
\newunicodechar{ω}{\ensuremath{\omega}}
\newunicodechar{Σ}{\ensuremath{\Sigma}}
% majuscules produites par \MakeUppercase dans les titres (α-aminés -> Α)
\newunicodechar{Α}{\ensuremath{\alpha}}
\newunicodechar{Β}{\ensuremath{\beta}}
\newunicodechar{Γ}{\ensuremath{\Gamma}}
\newunicodechar{Δ}{\ensuremath{\Delta}}
\newunicodechar{Ε}{\ensuremath{\varepsilon}}
\newunicodechar{Θ}{\ensuremath{\Theta}}
\newunicodechar{Λ}{\ensuremath{\Lambda}}
\newunicodechar{Μ}{\ensuremath{\mu}}
\newunicodechar{Τ}{\ensuremath{\tau}}
\newunicodechar{Φ}{\ensuremath{\Phi}}
\newunicodechar{Ψ}{\ensuremath{\Psi}}
\newunicodechar{Ω}{\ensuremath{\Omega}}
\newunicodechar{↦}{\ensuremath{\mapsto}}
\newunicodechar{∈}{\ensuremath{\in}}
\newunicodechar{∉}{\ensuremath{\notin}}
\newunicodechar{∧}{\ensuremath{\wedge}}
\newunicodechar{⇔}{\ensuremath{\Leftrightarrow}}
\newunicodechar{⟺}{\ensuremath{\Longleftrightarrow}}
\newunicodechar{⇒}{\ensuremath{\Rightarrow}}
\newunicodechar{⟹}{\ensuremath{\Longrightarrow}}
\newunicodechar{∘}{\ensuremath{\circ}}
\newunicodechar{∪}{\ensuremath{\cup}}
\newunicodechar{∩}{\ensuremath{\cap}}
\newunicodechar{≡}{\ensuremath{\equiv}}
\newunicodechar{⊂}{\ensuremath{\subset}}
\newunicodechar{⊥}{\ensuremath{\perp}}
\newunicodechar{∃}{\ensuremath{\exists}}
\newunicodechar{∀}{\ensuremath{\forall}}
\newunicodechar{∛}{\ensuremath{\sqrt[3]{\,}}}
\newunicodechar{∜}{\ensuremath{\sqrt[4]{\,}}}
\newunicodechar{⃗}{\ensuremath{{}^{\to}}}
\newunicodechar{ⁿ}{\ifmmode^{n}\else\textsuperscript{\ensuremath{n}}\fi}
\newunicodechar{ˣ}{\ifmmode^{x}\else\textsuperscript{\ensuremath{x}}\fi}
\newunicodechar{ᵇ}{\ifmmode^{b}\else\textsuperscript{\ensuremath{b}}\fi}
\newunicodechar{ʳ}{\ifmmode^{r}\else\textsuperscript{\ensuremath{r}}\fi}
\newunicodechar{ᵘ}{\ifmmode^{u}\else\textsuperscript{\ensuremath{u}}\fi}
\newunicodechar{ʸ}{\ifmmode^{y}\else\textsuperscript{\ensuremath{y}}\fi}
\newunicodechar{ᵃ}{\ifmmode^{a}\else\textsuperscript{\ensuremath{a}}\fi}
\newunicodechar{ᵉ}{\ifmmode^{e}\else\textsuperscript{\ensuremath{e}}\fi}
\newunicodechar{ᵏ}{\ifmmode^{k}\else\textsuperscript{\ensuremath{k}}\fi}
\newunicodechar{ᵖ}{\ifmmode^{p}\else\textsuperscript{\ensuremath{p}}\fi}
\newunicodechar{ᴺ}{\ifmmode^{N}\else\textsuperscript{\ensuremath{N}}\fi}
\newunicodechar{ᶜ}{\ifmmode^{c}\else\textsuperscript{\ensuremath{c}}\fi}
\newunicodechar{ʰ}{\ifmmode^{h}\else\textsuperscript{\ensuremath{h}}\fi}
\newunicodechar{ᵠ}{\ifmmode^{q}\else\textsuperscript{\ensuremath{q}}\fi}
\newunicodechar{ₙ}{\ifmmode_{n}\else\textsubscript{\ensuremath{n}}\fi}
\newunicodechar{ₐ}{\ifmmode_{a}\else\textsubscript{\ensuremath{a}}\fi}
\newunicodechar{ᵢ}{\ifmmode_{i}\else\textsubscript{\ensuremath{i}}\fi}
\newunicodechar{ᵣ}{\ifmmode_{r}\else\textsubscript{\ensuremath{r}}\fi}
\newunicodechar{ₖ}{\ifmmode_{k}\else\textsubscript{\ensuremath{k}}\fi}
\newunicodechar{ᵦ}{\ifmmode_{\beta}\else\textsubscript{\ensuremath{\beta}}\fi}
\newunicodechar{ₓ}{\ifmmode_{x}\else\textsubscript{\ensuremath{x}}\fi}
\newfontfamily\popdisplay{ArchivoBlack-Regular.ttf}[Path=../../../fonts/]
\newfontfamily\poplogo{SpaceGrotesk-Bold.ttf}[Path=../../../fonts/]
\newfontfamily\popsemi{PlusJakartaSans-SemiBold.ttf}[Path=../../../fonts/]
\newfontfamily\popxbold{PlusJakartaSans-ExtraBold.ttf}[Path=../../../fonts/]
\newfontfamily\popxboldls{PlusJakartaSans-ExtraBold.ttf}[Path=../../../fonts/,LetterSpace=3.0]

\setlist[enumerate]{leftmargin=1.7em,itemsep=5pt,topsep=4pt}
\setlist[itemize]{leftmargin=1.7em,itemsep=4pt,topsep=4pt,label={\color{poporange}\textbullet}}
\setlength{\parindent}{0pt}
\setlength{\parskip}{5pt}
\renewcommand{\arraystretch}{1.25}

% pgfplots : style Pop par défaut
\pgfplotsset{
  every axis/.append style={axis line style={popink,line width=1pt},tick style={popink},
    grid style={popline,line width=0.5pt},label style={font=\small},tick label style={font=\footnotesize}},
  popcurve/.style={poporange,line width=1.8pt,no markers,smooth},
}
% Même style disponible dans un \draw TikZ simple (hors pgfplots)
\tikzset{popcurve/.style={poporange,line width=1.8pt}}
\tikzset{popgrid/.style={popline,very thin}}
\colorlet{popcurve}{poporange} % le nom du style est parfois utilisé comme couleur

% ============================================================
% Pill / squiggle
% ============================================================
\newcommand{\poppill}[3]{%
  \tikz[baseline=-0.55ex]\node[draw=popink,line width=1.2pt,rounded corners=2.6pt,%
    fill=#2,inner xsep=6.5pt,inner ysep=3pt]{%
    \popxboldls\fontsize{7}{7}\selectfont\color{#3}\MakeUppercase{#1}};%
}
\newcommand{\popsquiggle}[1]{%
  \tikz[baseline=-0.4ex]{
    \draw[#1,line width=2.6pt,line cap=round]
      plot [smooth] coordinates {(0,0.09) (0.34,0.01) (0.68,0.21) (1.02,0.01) (1.36,0.10)};
  }%
}
% Mot-clé surligné (marqueur orange sous le texte)
\newcommand{\cle}[1]{{\popxbold\color{popdark}#1}}

% ============================================================
% En-tête / pied
% ============================================================
\pagestyle{fancy}
\fancyhf{}
\renewcommand{\headrulewidth}{0pt}
\renewcommand{\footrulewidth}{0pt}
\lhead{\raisebox{-0.15ex}{\poplogo\fontsize{13}{13}\selectfont\color{popink}OnBuch\color{poporange}+}}
\rhead{\poppill{\DOCMATIERE\ \textbullet\ \DOCNIVEAU\ \textbullet\ Leçon \DOCLECON}{white}{popink}}
\lfoot{\popxbold\fontsize{7.6}{7.6}\selectfont\color{popmuted}\MakeUppercase{Cours signé OnBuch+}\ \textbullet\ {\popsemi\DOCTITRE}}
\rfoot{\tikz[baseline=-0.6ex]\node[rounded corners=8.5pt,fill=popink,minimum height=17pt,minimum width=17pt,inner xsep=5pt,inner sep=0]{\popxbold\fontsize{8}{8}\selectfont\color{popcream}\thepage\,/\,\pageref*{LastPage}};}

% ============================================================
% Titres
% ============================================================
\newcommand{\chaptitre}[2]{%
  \begin{tcolorbox}[enhanced,colback=poporange,colframe=popink,boxrule=2.6pt,arc=11pt,
      drop shadow=popink,left=15pt,right=15pt,top=12pt,bottom=11pt,before skip=3pt,after skip=18pt]
    \poppill{#2}{popink}{popcream}\par\vspace{9pt}
    {\raggedright\hyphenpenalty=10000\exhyphenpenalty=10000\popdisplay\fontsize{22}{24}\selectfont\color{popcream}\MakeUppercase{#1}\par}
  \end{tcolorbox}
}
% Section numérotée : pastille orange avec le numéro + titre Archivo
\newcounter{popsec}
\newcommand{\coursec}[1]{%
  \stepcounter{popsec}\setcounter{popsub}{0}%
  \par\vspace{16pt}\noindent
  \begin{minipage}{\linewidth}
  \tikz[baseline=-0.7ex]\node[draw=popink,line width=1.6pt,fill=poporange,rounded corners=4pt,minimum size=22pt,inner sep=0]{\popdisplay\fontsize{12}{12}\selectfont\color{popcream}\thepopsec};%
  \hspace{8pt}{\popdisplay\fontsize{15}{17}\selectfont\color{popink}\MakeUppercase{#1}}\par
  \vspace{4pt}\noindent{\color{poporange}\rule{36pt}{3.6pt}}
  \end{minipage}\par\nopagebreak\vspace{8pt}
}
\newcounter{popsub}
\newcommand{\courssub}[1]{%
  \stepcounter{popsub}%
  \par\vspace{9pt}\noindent{\popxbold\fontsize{11.5}{13}\selectfont\color{popink}{\color{poporange}\thepopsec.\thepopsub}\hspace{6pt}#1}\par\nopagebreak\vspace{4pt}
}

% ============================================================
% Boîtes pédagogiques — même recette : fond blanc, bordure épaisse colorée,
% ombre dure encre, badge plein. Titre optionnel [#1].
% ============================================================
\tcbset{popbase/.style={breakable,enhanced,colback=white,boxrule=2.3pt,arc=9pt,
  shadow={3pt}{-3pt}{0pt}{popink},left=14pt,right=14pt,top=11pt,bottom=11pt,before skip=11pt,after skip=15pt,
  fontupper=\popsemi\fontsize{10.6}{14.5}\selectfont\color{popink},
  fontlower=\popsemi\fontsize{10.6}{14.5}\selectfont\color{popink}}}
% Coche dessinée (le glyphe ✓ n'existe pas dans les polices du document)
\newcommand{\popcheck}[1]{\tikz[baseline=-0.5ex]\draw[#1,line width=1.6pt,line cap=round,line join=round](-0.13,0.0)--(-0.04,-0.09)--(0.13,0.1);}
\providecommand{\coche}{\popcheck{popgreen}}
\providecommand{\resultat}[1]{\par\smallskip\noindent\fcolorbox{popgreen}{popgreenL}{\parbox{\dimexpr\linewidth-2\fboxsep-2\fboxrule\relax}{\textbf{Résultat.} #1}}\par\smallskip}
\newcommand{\popboxhead}[3]{\poppill{#1}{#2}{white}\ifx\relax#3\relax\else\hspace{6pt}{\popxbold\fontsize{10.6}{12}\selectfont\color{popink}#3}\fi\par\vspace{7pt}}

\newtcolorbox{aretenir}[1][]{popbase,colframe=popgreen,before upper={\popboxhead{À retenir}{popgreen}{#1}}}
\newtcolorbox{exemplebox}[1][]{popbase,colframe=poporange,before upper={\popboxhead{Exemple}{poporange}{#1}}}
\newtcolorbox{definition}[1][]{popbase,colframe=popblue,before upper={\popboxhead{Définition}{popblue}{#1}}}
\newtcolorbox{propriete}[1][]{popbase,colframe=poppurple,before upper={\popboxhead{Propriété}{poppurple}{#1}}}
\newtcolorbox{methode}[1][]{popbase,colframe=popgold,before upper={\popboxhead{Méthode}{popgold}{#1}}}
\newtcolorbox{attention}[1][]{popbase,colframe=poppink,before upper={\popboxhead{Attention}{poppink}{#1}}}
\newtcolorbox{experience}[1][]{popbase,colframe=popblue,colback=popblueL,before upper={\popboxhead{Expérience}{popblue}{#1}}}
% « Le savais-tu ? » : carte violette pleine (contexte, culture, Cameroun)
\newtcolorbox{savaistu}[1][]{popbase,colback=poppurple,colframe=popink,
  fontupper=\popsemi\fontsize{10.4}{14.2}\selectfont\color{white},
  before upper={\poppill{Le savais-tu\,?}{white}{poppurple}\ifx\relax#1\relax\else\hspace{6pt}{\popxbold\color{white}#1}\fi\par\vspace{7pt}}}
% Exercice résolu : énoncé en haut, \tcblower puis la solution sur fond crème
\newcounter{popexo}\newcounter{popexr}
\newtcolorbox{exoresolu}[1][]{popbase,colframe=poporange,colbacklower=popcream,
  segmentation style={popink,line width=1.2pt,dashed},
  before upper={\stepcounter{popexr}\popboxhead{Exercice résolu \thepopexr}{poporange}{#1}},
  before lower={\poppill{Solution}{popink}{popcream}\par\vspace{6pt}}}
% Exercice (non résolu en place) — la correction est dans la partie Corrigés
\newtcolorbox{exercice}[1][]{popbase,colframe=popink,
  before upper={\stepcounter{popexo}\popboxhead{Exercice \thepopexo}{popink}{#1}}}
% Corrigé
\newtcolorbox{corrige}[1][]{popbase,colframe=popgreen,colback=popgreenL,
  before upper={\popboxhead{Corrigé}{popgreen}{#1}}}
% Objectifs / prérequis / bilan
\newtcolorbox{objectifs}{popbase,colframe=popink,colback=poporangeL,
  before upper={\popboxhead{Objectifs de la leçon}{popink}{}}}
\newtcolorbox{prerequis}{popbase,colframe=popink,colback=popblueL,
  before upper={\popboxhead{Prérequis}{popblue}{}}}
\newtcolorbox{bilan}{popbase,colframe=popink,colback=white,boxrule=2.8pt,
  before upper={\popboxhead{Fiche bilan}{poporange}{L'essentiel à connaître pour l'examen}}}
% Figure encadrée avec légende
\newtcolorbox{popfigure}[1][]{enhanced,breakable=false,colback=white,colframe=popline,boxrule=1.4pt,arc=8pt,
  left=10pt,right=10pt,top=10pt,bottom=8pt,before skip=10pt,after skip=12pt,halign=center,
  lower separated=false,halign lower=center,
  fontlower=\popsemi\fontsize{9}{11.5}\selectfont\color{popmuted},
  #1}
\newcommand{\legende}[1]{\par\vspace{6pt}{\popsemi\fontsize{9}{11.5}\selectfont\color{popmuted}#1}}

% ============================================================
% Signature OnBuch+
% ============================================================
\newcommand{\poplogoinline}[1]{{\poplogo\fontsize{#1}{#1}\selectfont\color{popink}OnBuch\color{poporange}+}}
\newcommand{\signatureonbuch}{%
  \begin{tcolorbox}[enhanced,colback=popink,colframe=popink,boxrule=2.2pt,arc=10pt,
      drop shadow=poporange,left=14pt,right=14pt,top=10pt,bottom=10pt,before skip=0pt,after skip=0pt]
    \tikz[baseline=-0.7ex]\node[circle,fill=popgreen,draw=popcream,line width=1.3pt,minimum size=22pt,inner sep=0]{\popcheck{white}};%
    \hspace{9pt}\begin{minipage}[c]{0.62\linewidth}
      {\popxbold\fontsize{12}{13}\selectfont\color{popcream}Signé\ \textbullet\ {\poplogo OnBuch\color{poporange}+}}\par\vspace{2pt}
      {\raggedright\popsemi\fontsize{8}{10.5}\selectfont\color{popline}\MakeUppercase{Équipe pédagogique OnBuch+}\\\MakeUppercase{Conforme au programme officiel MINESEC}\par}
    \end{minipage}\hfill
    {\poplogo\fontsize{22}{22}\selectfont\color{popcream}OnBuch\color{poporange}+}
  \end{tcolorbox}%
}

% ============================================================
% Page de garde
% #1 titre · #2 matière · #3 niveau · #4 module · #5 n° leçon · #6 durée ·
% #7 description / objectifs (texte ou itemize)
% ============================================================
\newcommand{\pagegarde}[7]{%
  \thispagestyle{empty}
  \newgeometry{a4paper,margin=1.7cm,top=1.6cm,bottom=1.4cm}%
  \begin{tikzpicture}[remember picture,overlay]
    \fill[poporange!18] ([xshift=-3.2cm,yshift=-3.6cm]current page.north east) circle (4.6cm);
    \fill[poppurple!14] ([xshift=2.4cm,yshift=7.6cm]current page.south west) circle (3.8cm);
  \end{tikzpicture}%
  {\poplogo\fontsize{30}{30}\selectfont\color{popink}OnBuch\color{poporange}+}\hfill
  \raisebox{0.6ex}{\poppill{Cours signé \textbullet\ Premium}{popink}{popcream}}\par
  \vspace{1pt}\popsquiggle{poppurple}\par
  \vspace{8pt}
  \begin{tcolorbox}[enhanced,colback=poporange,colframe=popink,boxrule=2.8pt,arc=13pt,
      drop shadow=popink,left=18pt,right=18pt,top=12pt,bottom=13pt,after skip=10pt]
    \poppill{#2 \textbullet\ #3}{popink}{popcream}\hspace{5pt}\poppill{#4}{white}{popink}\par\vspace{8pt}
    {\popdisplay\fontsize{14}{14}\selectfont\color{popink}LEÇON #5}\par\vspace{4pt}
    {\raggedright\hyphenpenalty=10000\exhyphenpenalty=10000\popdisplay\fontsize{25}{27}\selectfont\color{popcream}\MakeUppercase{#1}\par}
    \vspace{8pt}
    {\popxbold\fontsize{9.5}{9.5}\selectfont\color{popink}\MakeUppercase{Durée conseillée : #6}}
  \end{tcolorbox}
  \begin{tcolorbox}[enhanced,colback=white,colframe=popink,boxrule=2.2pt,arc=10pt,
      drop shadow=popink,left=16pt,right=16pt,top=10pt,bottom=10pt,after skip=8pt]
    \poppill{Dans cette leçon}{poppurple}{white}\par\vspace{5pt}
    \popsemi\fontsize{9.3}{12.6}\selectfont\color{popink}#7
  \end{tcolorbox}
  \noindent\begin{minipage}[t]{0.487\linewidth}
    \begin{tcolorbox}[enhanced,colback=popgreen,colframe=popink,boxrule=2.2pt,arc=10pt,
        drop shadow=popink,left=11pt,right=11pt,top=10pt,bottom=10pt,height=3.1cm,valign=center]
      \tcbox[colback=white,colframe=popink,boxrule=1.3pt,arc=5pt,boxsep=0pt,nobeforeafter,
        left=3pt,right=3pt,top=3pt,bottom=3pt,valign=center]{\qrcode[height=1.9cm]{https://play.google.com/store/apps/details?id=com.luvvix.onbuch&hl=fr}}%
      \hspace{8pt}\begin{minipage}[c]{0.5\linewidth}
        {\popdisplay\fontsize{11}{12.5}\selectfont\color{white}\MakeUppercase{Télécharge OnBuch}}\par\vspace{4pt}
        {\raggedright\popsemi\fontsize{8.4}{11}\selectfont\color{white}Gratuit \textbullet\ Google Play\\Tuteur IA, annales, concours, résultats\par}
      \end{minipage}
    \end{tcolorbox}
  \end{minipage}\hfill
  \begin{minipage}[t]{0.487\linewidth}
    \begin{tcolorbox}[enhanced,colback=poppurple,colframe=popink,boxrule=2.2pt,arc=10pt,
        drop shadow=popink,left=11pt,right=11pt,top=10pt,bottom=10pt,height=3.1cm,valign=center]
      \tikz[baseline=-0.7ex]\node[circle,fill=white,draw=popink,line width=1.3pt,minimum size=1.6cm,inner sep=0]{\popdisplay\fontsize{24}{24}\selectfont\color{poppurple}+};%
      \hspace{8pt}\begin{minipage}[c]{0.6\linewidth}
        {\popdisplay\fontsize{11}{12.5}\selectfont\color{white}\MakeUppercase{Passe à OnBuch+}}\par\vspace{4pt}
        {\raggedright\popsemi\fontsize{8.4}{11}\selectfont\color{white}Tous les cours signés, Tuteur IA illimité, fiches PDF — onglet OnBuch+ de l'app\par}
      \end{minipage}
    \end{tcolorbox}
  \end{minipage}
  \vspace{8pt}
  \popcontenu
  \vfill
  \signatureonbuch
  \restoregeometry
  \newpage
}

% Rangée « Ce cours contient » (page de garde)
\newcommand{\poptile}[3]{% #1 couleur #2 gros texte #3 légende
  \begin{tcolorbox}[enhanced,colback=#1,colframe=popink,boxrule=2pt,arc=9pt,shadow={2.5pt}{-2.5pt}{0pt}{popink},
      left=6pt,right=6pt,top=9pt,bottom=9pt,halign=center,valign=center,height=1.75cm,width=\linewidth,nobeforeafter]
    {\popdisplay\fontsize{12.5}{13}\selectfont\color{white}\MakeUppercase{#2}}\par\vspace{3pt}
    {\centering\popsemi\fontsize{7.8}{9.5}\selectfont\color{white}#3\par}
  \end{tcolorbox}}
\newcommand{\popcontenu}{%
  \par\noindent{\popxboldls\fontsize{7.5}{7.5}\selectfont\color{popmuted}CE COURS SIGNÉ CONTIENT}\par\vspace{7pt}
  \noindent\begin{minipage}[t]{0.235\linewidth}\poptile{popblue}{Cours}{complet, illustré, pas à pas}\end{minipage}\hfill
  \begin{minipage}[t]{0.235\linewidth}\poptile{poporange}{Méthodes}{et exercices résolus}\end{minipage}\hfill
  \begin{minipage}[t]{0.235\linewidth}\poptile{poppink}{Exercices}{type examen + corrigés}\end{minipage}\hfill
  \begin{minipage}[t]{0.235\linewidth}\poptile{popgreen}{Fiche bilan}{l'essentiel à retenir}\end{minipage}\par}

% Sommaire
\newtcolorbox{sommaire}{enhanced,colback=white,colframe=popink,boxrule=2.2pt,arc=10pt,
  drop shadow=popink,left=18pt,right=18pt,top=13pt,bottom=13pt,before skip=6pt,after skip=20pt,
  before upper={\poppill{Sommaire}{poppurple}{white}\par\vspace{9pt}\popsemi\fontsize{10.6}{14.5}\selectfont\color{popink}}}

% Signature de fin de cours — poussée tout en bas de la dernière page
\newcommand{\findecours}{%
  \par\vfill\nopagebreak
  \begin{center}\popsquiggle{poporange}\end{center}\vspace{4pt}
  \signatureonbuch
}
\AtBeginDocument{\def\llbracket{\mathopen{[\mkern-3.5mu[}}\def\rrbracket{\mathclose{]\mkern-3.5mu]}}} % intervalles d'entiers (glyphe ⟦ absent de la police)
% Glyphes absents de Fira Math : redéfinis à partir de glyphes disponibles
\AtBeginDocument{%
  \def\setminus{\mathbin{\backslash}}\let\smallsetminus\setminus
  \def\vdots{\mathord{\vbox{\baselineskip3.2pt\lineskiplimit0pt\kern4pt\hbox{.}\hbox{.}\hbox{.}}}}%
  \def\ddots{\mathinner{\mkern1mu\raise6.5pt\hbox{.}\mkern2mu\raise3.6pt\hbox{.}\mkern2mu\raise0.7pt\hbox{.}\mkern1mu}}%
  \def\triangle{\mathord{\scalebox{0.9}{$\symup{\Delta}$}}}%
  \def\nabla{\mathord{\raisebox{\depth}{\scalebox{1}[-1]{$\symup{\Delta}$}}}}%
  \def\checkmark{\popcheck{popdark}}%
  \def\star{\mathbin{\ast}}\def\bigstar{\mathord{\ast}}%
  \def\diamond{\mathbin{\scalebox{0.6}{$\Diamond$}}}%
}

%% FIGFIX-BEGIN v5 — correctifs globaux des figures (docs/onbuch-generation/tools/figfix)
\usepackage{adjustbox}\usepackage{etoolbox}\tcbuselibrary{raster}
% toute figure TikZ/pgfplots plus large que la ligne est réduite à la largeur disponible
\BeforeBeginEnvironment{tikzpicture}{\begin{adjustbox}{max width=\linewidth}}
\AfterEndEnvironment{tikzpicture}{\end{adjustbox}}
% légendes pgfplots lisibles par-dessus les courbes
\pgfplotsset{every axis/.append style={legend style={fill=white,fill opacity=0.92,text opacity=1,draw=black!15,inner sep=3pt}}}
% étiquettes d'arêtes (syntaxe edge["..."]) avec fond blanc
\tikzset{every edge quotes/.append style={fill=white,inner sep=1.2pt}}
%% FIGFIX-END
\begin{document}

\pagegarde{Oscillations électriques libres (circuits LC et RLC)}{Physique}{Tle D}{Module 3 : Électricité — évolutions temporelles des circuits électriques et électroniques}{P11}{4 h}{Cette leçon étudie les oscillations électriques libres : d'abord dans le circuit LC idéal où l'énergie se conserve en s'échangeant entre condensateur et bobine, puis dans le circuit RLC réel où l'amortissement fait apparaître les régimes pseudo-périodique, critique et apériodique. Elle fournit les outils d'analyse (équations différentielles, solutions, période propre, bilan énergétique) indispensables pour comprendre la résonance et la réception radio.
\vspace{4pt}
\begin{multicols}{2}\small
\begin{itemize}
\item À la fin de cette leçon, je sais écrire la tension aux bornes d'une bobine réelle u = ri + L di/dt et l'énergie magnétique qu'elle stocke
\item À la fin de cette leçon, je sais établir l'équation différentielle des oscillations libres d'un circuit LC à partir de la loi des mailles
\item À la fin de cette leçon, je sais exprimer la charge q(t), l'intensité i(t) et vérifier qu'elles sont solutions de l'équation différentielle
\item À la fin de cette leçon, je sais calculer la pulsation propre ω0, la période propre T0 = 2π√(LC) et la fréquence propre f0 d'un circuit LC
\item À la fin de cette leçon, je sais exprimer l'énergie électromagnétique totale d'un circuit LC et démontrer sa conservation
\item À la fin de cette leçon, je sais établir l'équation différentielle des oscillations amorties d'un circuit RLC série
\end{itemize}\end{multicols}}

\begin{sommaire}
\begin{multicols}{2}
\begin{enumerate}
\item La bobine : inductance, tension et énergie magnétique
\item Le circuit LC idéal : équation différentielle des oscillations libres
\item Période propre, fréquence et applications
\item Étude énergétique du circuit LC : conservation de l'énergie électromagnétique
\item Le circuit RLC série : oscillations amorties
\item Analogies électromécaniques et synthèse
\item Activité d'intégration
\item Méthodes et exercices résolus
\item Exercices
\item Corrigés des exercices
\item Fiche bilan
\end{enumerate}
\end{multicols}
\end{sommaire}

\begin{objectifs}
\begin{itemize}
\item À la fin de cette leçon, je sais écrire la tension aux bornes d'une bobine réelle u = ri + L di/dt et l'énergie magnétique qu'elle stocke
\item À la fin de cette leçon, je sais établir l'équation différentielle des oscillations libres d'un circuit LC à partir de la loi des mailles
\item À la fin de cette leçon, je sais exprimer la charge q(t), l'intensité i(t) et vérifier qu'elles sont solutions de l'équation différentielle
\item À la fin de cette leçon, je sais calculer la pulsation propre ω0, la période propre T0 = 2π√(LC) et la fréquence propre f0 d'un circuit LC
\item À la fin de cette leçon, je sais exprimer l'énergie électromagnétique totale d'un circuit LC et démontrer sa conservation
\item À la fin de cette leçon, je sais établir l'équation différentielle des oscillations amorties d'un circuit RLC série
\item À la fin de cette leçon, je sais distinguer et reconnaître graphiquement les régimes pseudo-périodique, critique et apériodique
\item À la fin de cette leçon, je sais exploiter l'analogie électromécanique (ressort amorti ↔ circuit RLC) pour interpréter les grandeurs électriques
\end{itemize}
\end{objectifs}

\begin{prerequis}
\begin{itemize}
\item Condensateur : q = C·u, énergie électrostatique ½ q²/C, charge et décharge du circuit RC (Première et début de Tle)
\item Dérivation des fonctions sinusoïdales et équations différentielles du type x'' + ω²x = 0 (mathématiques, Tle C)
\item Oscillations mécaniques libres : pendule élastique, équation x'' + (k/m)x = 0, période propre (leçon de mécanique du Module 2)
\item Loi des mailles et loi d'Ohm en régime variable ; relation i = dq/dt
\item Notion d'inductance vue en Première : phénomène d'auto-induction et loi de Lenz
\end{itemize}
\end{prerequis}

\newpage

\coursec{La bobine : inductance, tension et énergie magnétique}

Au quartier, quand le groupe électrogène s'arrête brutalement, l'ampoule alimentée par un condensateur de flash d'appareil photo ne s'éteint pas instantanément : elle clignote faiblement avant de s'éteindre. De même, le poste radio capte la fréquence \SI{94,5}{MHz} de CRTV grâce à un minuscule circuit bobine-condensateur qui « résonne » exactement à cette fréquence. Comment une bobine et un condensateur, sans aucune pile, peuvent-ils entretenir un échange d'énergie rythmé, comme un pendule ? Et pourquoi cet échange finit-il toujours par s'éteindre dans un circuit réel ? C'est toute la physique des oscillations électriques libres.

Avant d'assembler le circuit oscillant, il faut comprendre son acteur principal : la \cle{bobine}. Tu connais déjà le condensateur, capable d'emmagasiner de l'énergie électrique. La bobine, elle, emmagasine de l'énergie \emph{magnétique} et s'oppose à toute variation du courant qui la traverse. C'est cette « inertie électrique » qui va jouer, dans le circuit LC, le rôle de la masse du pendule.

\courssub{Le phénomène d'auto-induction : une bobine « têtue »}

\textbf{Observation expérimentale.}\par

Branchons en série un générateur, un interrupteur, une lampe et une bobine de nombreuses spires. À la fermeture de l'interrupteur, la lampe s'allume \emph{progressivement}, avec un retard perceptible par rapport à un circuit identique sans bobine. À l'ouverture, on observe parfois une petite étincelle entre les lames de l'interrupteur. La bobine semble « retarder » l'établissement du courant et « protester » quand on veut le supprimer.

\textbf{Interprétation.}\par

Une bobine parcourue par un courant d'intensité $i$ crée un champ magnétique $\vec{B}$ proportionnel à $i$ dans son espace intérieur. Le flux de ce champ à travers ses propres spires, appelé \cle{flux propre} $\Phi$, est donc proportionnel à $i$. Or, d'après la loi de Faraday, toute variation du flux à travers un circuit y induit une force électromotrice. Ici, c'est le circuit \emph{lui-même} qui crée le flux variable : on parle d'\cle{auto-induction}. La f.é.m. induite $e = -\dfrac{d\Phi}{dt}$ s'oppose, d'après la \cle{loi de Lenz}, à la variation de courant qui lui a donné naissance :

\begin{itemize}
\item si $i$ augmente, la bobine crée une f.é.m. qui freine cette augmentation ;
\item si $i$ diminue, la bobine crée une f.é.m. qui s'oppose à cette diminution.
\end{itemize}

\begin{attention}[La loi de Lenz n'interdit pas la variation]
La bobine ne \emph{supprime} pas la variation du courant : elle la \emph{retarde}. Si elle l'empêchait totalement, il n'y aurait plus de variation de flux, donc plus de f.é.m. induite : le raisonnement serait contradictoire. La loi de Lenz exprime une opposition, jamais une annulation.
\end{attention}

\courssub{Inductance propre $L$ : définition et unité}

La proportionnalité entre le flux propre et l'intensité conduit à la définition centrale de cette leçon.

\begin{definition}[Inductance propre d'une bobine]
L'\cle{inductance propre} $L$ d'une bobine est le coefficient de proportionnalité entre le flux propre $\Phi$ à travers la bobine et l'intensité $i$ du courant qui la parcourt :
\[ \Phi = L\,i \]
$L$ s'exprime en \cle{henry} (symbole H). D'après $\Phi = Li$, on a $\SI{1}{H} = \SI{1}{Wb.A^{-1}}$.
\end{definition}

L'inductance ne dépend que de la \emph{géométrie} de la bobine (nombre de spires $N$, longueur $\ell$, section $S$) et du milieu qui l'entoure. Pour une bobine longue (solénoïde) de $N$ spires, de section $S$ et de longueur $\ell$, dans l'air, on montre en classe de Première que :
\[ L = \mu_0\,\frac{N^2 S}{\ell} \]
où $\mu_0 = 4\pi \times 10^{-7}\ \mathrm{H.m^{-1}}$ est la perméabilité magnétique du vide. Insérer un noyau de fer doux multiplie $L$ par un facteur pouvant dépasser $10^3$ : c'est le principe des bobines des transformateurs d'ENEO.

\begin{aretenir}[Ordres de grandeur]
En électronique, le henry est une grande unité. On utilise couramment le \cle{millihenry} ($\SI{1}{mH} = 10^{-3}\ \mathrm{H}$) et le \cle{microhenry} ($\SI{1}{\micro H} = 10^{-6}\ \mathrm{H}$). Une bobine de laboratoire : $L \approx 0,1$ à $\SI{1}{H}$ ; une bobine de circuit radio : quelques µH.
\end{aretenir}

\begin{savaistu}[Des bobines minuscules dans ta radio]
Les bobines des postes radio portables ont des inductances de quelques µH seulement, réalisées par quelques spires de fil de cuivre bobiné en l'air, parfois autour d'un petit bâtonnet de ferrite. Associée à un condensateur, une bobine de $\SI{0,3}{\micro H}$ permet de « syntoniser » la fréquence $\SI{94,5}{MHz}$ de CRTV : c'est le circuit LC que nous étudierons dans la section suivante. Heinrich Hertz, qui a donné son nom à l'unité de fréquence, et Joseph Henry, physicien américain qui découvrit l'auto-induction en 1832 (en même temps que Faraday), sont les deux « héros » cachés de ta radio.
\end{savaistu}

\courssub{Tension aux bornes d'une bobine réelle}

\textbf{Modélisation.}\par

Une bobine réelle est un enroulement de fil conducteur : elle possède donc une \cle{résistance interne} $r$ (celle du fil, d'autant plus grande que le fil est long et fin). On la modélise par l'association en série d'un résistor de résistance $r$ et d'une \emph{bobine idéale} d'inductance $L$. En \cle{convention récepteur} (flèche de la tension $u$ opposée au sens du courant $i$), les tensions s'ajoutent :

\begin{propriete}[Tension aux bornes d'une bobine réelle]
En convention récepteur, la tension aux bornes d'une bobine réelle d'inductance $L$ et de résistance interne $r$, parcourue par un courant $i$, est :
\[ u = r\,i + L\,\dfrac{di}{dt} \]
\begin{itemize}
\item le terme $r\,i$ est la chute de tension ohmique dans le fil ;
\item le terme $L\,\dfrac{di}{dt}$ est la tension d'auto-induction, qui ne s'annule que si le courant est constant.
\end{itemize}
\end{propriete}

\begin{experience}[Établissement guidé de la relation $u = ri + L\,\dfrac{di}{dt}$]
\begin{enumerate}
\item La bobine réelle est équivalente à un résistor $r$ en série avec une bobine idéale $L$. En convention récepteur, la tension totale est la somme : $u = u_r + u_L$.
\item Aux bornes du résistor : $u_r = r\,i$ (loi d'Ohm).
\item Aux bornes de la bobine idéale, la tension compense la f.é.m. d'auto-induction $e = -\dfrac{d\Phi}{dt}$. Avec $\Phi = Li$ et $L$ constante : $u_L = -e = \dfrac{d\Phi}{dt} = L\,\dfrac{di}{dt}$.
\item Conclusion : $u = r\,i + L\,\dfrac{di}{dt}$.
\end{enumerate}
Vérification dimensionnelle : $[L\,di/dt] = \mathrm{H.A.s^{-1}} = \mathrm{Wb.s^{-1}} = \mathrm{V}$, homogène à une tension, comme $[r\,i] = \Omega\cdot\mathrm{A} = \mathrm{V}$.
\end{experience}

\begin{popfigure}
\begin{center}
\begin{tikzpicture}[european, scale=1.1]
\draw[popblue, thick] (0,0) to[R, l_={$r$}, v=$r\,i$, i=$i$] (3,0)
to[L, l_={$L$}, v=$L\,\dfrac{di}{dt}$] (6,0);
\draw[popdark, thick] (0,0) -- (-0.8,0) node[left]{A};
\draw[popdark, thick] (6,0) -- (6.8,0) node[right]{B};
\draw[->, poporange, very thick] (0.2,1.6) -- (5.8,1.6) node[midway, above, poporange]{$u = r\,i + L\,\dfrac{di}{dt}$};
\node[popmuted] at (3,-1.1) {Convention récepteur : flèche de $u$ opposée à $i$};
\end{tikzpicture}
\end{center}
\legende{Modèle série de la bobine réelle : résistance interne $r$ en série avec l'inductance idéale $L$.}
\end{popfigure}

\textbf{Cas limites à connaître par cœur.}\par

\begin{itemize}
\item \textbf{Bobine idéale} ($r = 0$) : $u = L\,\dfrac{di}{dt}$. Si le courant est \emph{constant} ($di/dt = 0$), alors $u = 0$ : la bobine idéale se comporte comme un simple \emph{fil}. Elle ne s'oppose pas au courant constant, seulement à ses \emph{variations}.
\item \textbf{Régime permanent} ($i = I_0$ constant) dans une bobine réelle : $u = r\,I_0$. La bobine se comporte comme un résistor de résistance $r$.
\item \textbf{Régime variable} : les deux termes coexistent, et le terme $L\,di/dt$ peut devenir très grand si le courant varie vite.
\end{itemize}

\begin{attention}[Les deux erreurs classiques]
\begin{enumerate}
\item \textbf{Oublier le terme $r\,i$} pour une bobine réelle. Au Bac, l'énoncé précise souvent « bobine d'inductance $L$ et de résistance $r$ » : la tension est alors $u = ri + L\dfrac{di}{dt}$, et non $L\dfrac{di}{dt}$ seul. Cette omission fausse l'équation différentielle du circuit RLC.
\item \textbf{Croire que la bobine idéale s'oppose au courant constant}. Faux : en régime permanent, $u = 0$ aux bornes d'une bobine idéale. La bobine est une « passoire » pour le courant continu établi et un « mur » pour les variations brutales.
\end{enumerate}
\end{attention}

\courssub{Continuité de l'intensité dans une bobine}

\begin{propriete}[Continuité du courant dans une bobine]
L'intensité $i$ du courant traversant une bobine est une fonction \cle{continue} du temps : elle ne peut pas subir de discontinuité.
\end{propriete}

\textbf{Justification énergétique.}\par

Supposons que $i$ passe brutalement d'une valeur $i_1$ à une valeur $i_2 \neq i_1$ en un intervalle de temps $\Delta t \to 0$. Alors $\dfrac{di}{dt} \approx \dfrac{i_2 - i_1}{\Delta t} \to \infty$, donc la tension $L\,\dfrac{di}{dt}$ et la puissance $p = u\,i$ deviendraient infinies. Une puissance infinie est physiquement impossible : aucun dispositif réel ne peut la fournir. L'intensité dans une bobine varie donc toujours \emph{progressivement}.

C'est l'analogue électrique d'un résultat mécanique bien connu : la vitesse d'un solide ne peut pas changer instantanément (cela exigerait une force infinie). La bobine confère au courant une \emph{inertie}.

\begin{attention}[Conséquence pratique : l'étincelle de rupture]
Quand on ouvre brutalement un interrupteur dans un circuit contenant une grosse bobine, le courant « veut » continuer à circuler : la tension $L\,di/dt$ devient si grande qu'elle ionise l'air entre les lames de l'interrupteur, créant une étincelle. C'est pourquoi on place une diode de « roue libre » aux bornes des bobines des relais automobiles, et pourquoi on ne débranche jamais violemment un gros électroaimant sous tension.
\end{attention}

\courssub{Énergie magnétique emmagasinée : $E_L = \dfrac{1}{2}Li^2$}

\textbf{Intuition.}\par

Pour établir un courant $i$ dans une bobine, le générateur doit « lutter » contre la f.é.m. d'auto-induction : il fournit un travail. Ce travail n'est pas perdu (hors effet Joule) : il est stocké dans la bobine sous forme d'\cle{énergie magnétique}, restituable lorsque le courant diminue. C'est exactement comme lancer une masse : le travail fourni devient énergie cinétique $\frac{1}{2}mv^2$, récupérable au freinage.

\begin{propriete}[Énergie magnétique d'une bobine — démonstration exigible]
Une bobine idéale d'inductance $L$, parcourue par un courant d'intensité $i$, emmagasine l'énergie magnétique :
\[ E_L = \frac{1}{2}\,L\,i^2 \]
$E_L$ en joules (J), $L$ en henrys (H), $i$ en ampères (A).
\end{propriete}

\begin{experience}[Démonstration guidée par intégration de la puissance]
Considérons une bobine idéale ($r = 0$, donc aucune perte Joule) dont le courant passe progressivement de $0$ à la valeur $i$.
\begin{enumerate}
\item La puissance électrique reçue par la bobine en convention récepteur est :
\[ p = u\,i = L\,\dfrac{di}{dt}\,i \]
\item L'énergie élémentaire reçue pendant $dt$ est :
\[ dE_L = p\,dt = L\,i\,\dfrac{di}{dt}\,dt = L\,i\,di \]
\item On intègre de l'état initial ($i = 0$, $E_L = 0$) à l'état final (courant $i$) :
\[ E_L = \int_0^{i} L\,i'\,di' = L\left[\dfrac{i'^2}{2}\right]_0^{i} = \frac{1}{2}L\,i^2 \]
\end{enumerate}
Vérification dimensionnelle : $[L i^2] = \mathrm{H.A^2} = \mathrm{V.s.A^{-1}.A^2} = \mathrm{V.A.s} = \mathrm{W.s} = \mathrm{J}$. La formule est homogène.
\end{experience}

\begin{popfigure}
\begin{center}
\begin{tikzpicture}
\begin{axis}[
width=0.85\linewidth, height=6cm,
xlabel={$t$ (s)}, ylabel={$i$ (A)},
xmin=0, xmax=5, ymin=0, ymax=2.6,
xtick={5}, xticklabels={},
ytick={2}, yticklabels={$i$},
grid=both, grid style={popline!40},
axis lines=left]
\addplot[popcurve, domain=0:5, samples=200] {2*(1-exp(-x/1.2))};
\addplot[popblue!25, draw=none, domain=0:5, samples=200] {2*(1-exp(-x/1.2))} \closedcycle;
\addplot[popmuted, dashed, domain=0:5] {2};
\node[popdark] at (axis cs:3.4,0.55) {$E_L = \displaystyle\int_0^i L\,i'\,di' = \frac{1}{2}Li^2$};
\node[popmuted] at (axis cs:4.2,2.15) {régime permanent};
\end{axis}
\end{tikzpicture}
\end{center}
\legende{Établissement progressif du courant dans une bobine : l'intensité croît continûment de $0$ à $i$ ; l'énergie magnétique stockée s'obtient en intégrant la puissance $p = Li\,\dfrac{di}{dt}$ pendant toute la montée du courant.}
\end{popfigure}

\textbf{Analogie formelle avec l'énergie cinétique.}\par

\begin{center}
\begin{tabularx}{\linewidth}{|l|X|X|}
\hline
\rowcolor{popblueL} & \textbf{Mécanique} & \textbf{Électricité} \\
\hline
Grandeur d'« inertie » & masse $m$ (kg) & inductance $L$ (H) \\
\hline
Grandeur de « mouvement » & vitesse $v$ (m/s) & intensité $i$ (A) \\
\hline
Énergie stockée & $E_c = \dfrac{1}{2}mv^2$ & $E_L = \dfrac{1}{2}Li^2$ \\
\hline
Continuité & $v$ continue (force finie) & $i$ continue (tension finie) \\
\hline
\end{tabularx}
\end{center}

L'inductance $L$ joue le rôle d'une \cle{inertie électrique} : de même qu'une masse lourde est difficile à accélérer ou à arrêter, une bobine de forte inductance s'oppose vigoureusement à toute variation du courant. Cette analogie, que nous complèterons en fin de leçon, est la clé de voûte pour comprendre les oscillations du circuit LC : l'énergie « cinétique » $\frac{1}{2}Li^2$ de la bobine et l'énergie « potentielle » $\frac{q^2}{2C}$ du condensateur vont s'échanger indéfiniment, comme dans un pendule.

\begin{exemplebox}[Énergie stockée dans une bobine de laboratoire]
Une bobine d'inductance $L = \SI{0,20}{H}$ est parcourue par un courant $i = \SI{2,0}{A}$. Son énergie magnétique vaut :
\[ E_L = \frac{1}{2}Li^2 = \frac{1}{2}\times 0,20 \times (2,0)^2 = \SI{0,40}{J} \]
Comparons avec un condensateur de capacité $C = \SI{100}{\micro F}$ chargé sous $U = \SI{90}{V}$ :
\[ E_C = \frac{1}{2}CU^2 = \frac{1}{2}\times 100\times 10^{-6}\times 90^2 = \SI{0,405}{J} \approx \SI{0,40}{J} \]
Les deux réservoirs stockent la même énergie : s'ils sont associés dans un circuit, ils pourront se l'échanger intégralement — c'est le cœur de l'oscillation électrique libre.
\end{exemplebox}

\begin{exoresolu}[Bobine réelle en régime variable]
Une bobine réelle d'inductance $L = \SI{50}{mH}$ et de résistance interne $r = \SI{10}{\Omega}$ est parcourue par un courant $i(t) = 3{,}0\,t$ (en A, avec $t$ en s) pendant la phase de démarrage d'un montage.
\begin{enumerate}
\item Calculer la tension $u$ aux bornes de la bobine à l'instant $t = \SI{0,50}{s}$.
\item Calculer l'énergie magnétique emmagasinée à cet instant.
\item Que vaut $u$ lorsque le courant atteint sa valeur permanente $I_0 = \SI{1,5}{A}$ ?
\end{enumerate}
\tcblower
\textbf{1.} On applique la loi de la bobine réelle en convention récepteur : $u = ri + L\dfrac{di}{dt}$.
À $t = \SI{0,50}{s}$ : $i = 3{,}0 \times 0{,}50 = \SI{1,5}{A}$ et $\dfrac{di}{dt} = \SI{3,0}{A.s^{-1}}$.
\begin{align*}
u &= r\,i + L\,\dfrac{di}{dt} = 10 \times 1{,}5 + 50\times 10^{-3} \times 3{,}0 \\
&= 15 + 0{,}15 = \SI{15,15}{V} \approx \SI{15,2}{V}
\end{align*}
Remarque : le terme inductif ($\SI{0,15}{V}$) est ici négligeable devant le terme ohmique ($\SI{15}{V}$), car le courant varie lentement.

\textbf{2.} Énergie magnétique à $t = \SI{0,50}{s}$ :
\[ E_L = \frac{1}{2}Li^2 = \frac{1}{2}\times 50\times 10^{-3}\times (1{,}5)^2 = \SI{56,3}{mJ} \]

\textbf{3.} En régime permanent, $i = I_0$ constant donc $\dfrac{di}{dt} = 0$ :
\[ u = r\,I_0 = 10 \times 1{,}5 = \SI{15}{V} \]
La bobine se comporte alors comme un simple résistor de résistance $r = \SI{10}{\Omega}$ : toute la puissance reçue est dissipée par effet Joule, et l'énergie magnétique reste figée à $\frac{1}{2}LI_0^2$.
\end{exoresolu}

\begin{aretenir}[L'essentiel de la section]
\begin{itemize}
\item Une bobine parcourue par un courant variable est le siège d'une f.é.m. d'\cle{auto-induction} qui s'oppose à la variation du courant (loi de Lenz).
\item Flux propre : $\Phi = Li$ ; $L$ en henrys (H), sous-multiples usuels mH et µH.
\item Bobine réelle en convention récepteur : $u = ri + L\dfrac{di}{dt}$ ; bobine idéale : $u = L\dfrac{di}{dt}$ ; en régime permanent, une bobine idéale équivaut à un fil.
\item L'intensité $i$ dans une bobine est toujours \cle{continue} (sinon puissance infinie).
\item Énergie magnétique : $E_L = \dfrac{1}{2}Li^2$, obtenue en intégrant $p = Li\dfrac{di}{dt}$ de $0$ à $i$.
\item $L$ est l'analogue électrique de la masse $m$ : $\frac{1}{2}Li^2 \leftrightarrow \frac{1}{2}mv^2$.
\end{itemize}
\end{aretenir}

Nous disposons maintenant des deux « réservoirs » d'énergie : le condensateur ($E_C = \frac{q^2}{2C}$, énergie électrique) et la bobine ($E_L = \frac{1}{2}Li^2$, énergie magnétique). Dans la section suivante, nous les associerons dans le circuit LC et nous établirons l'équation différentielle qui gouverne leur échange rythmé d'énergie : les oscillations électriques libres.

\coursec{Le circuit LC idéal : équation différentielle des oscillations libres}

Dans la section précédente, nous avons découvert la bobine : un dipôle capable de s'opposer aux variations de l'intensité qui la traverse, caractérisée par son inductance $L$, avec une tension $u_L = ri + L\dfrac{di}{dt}$ et une énergie magnétique emmagasinée $\mathcal{E}_m = \dfrac{1}{2}Li^2$. Nous savons aussi, depuis la classe de Première, qu'un condensateur chargé stocke de l'énergie électrique $\mathcal{E}_e = \dfrac{q^2}{2C}$. Que se passe-t-il lorsqu'on réunit ces deux réservoirs d'énergie dans un même circuit ? C'est toute la magie des \cle{oscillations électriques libres} : l'énergie va voyager sans fin du condensateur vers la bobine, puis de la bobine vers le condensateur, exactement comme l'énergie oscille entre forme potentielle et forme cinétique dans un pendule. Cette section établit rigoureusement ce phénomène, qui est au cœur du programme de Terminale C.

\courssub{Description du dispositif expérimental}

Considérons le montage suivant : un générateur de tension continue de force électromotrice $E$, un interrupteur $K$ à deux positions, un condensateur de capacité $C$ et une bobine \emph{idéale} d'inductance $L$ (c'est-à-dire de résistance interne $r$ négligeable, $r = 0$).

\begin{experience}[Charge du condensateur puis association avec la bobine]
\textbf{Matériel :} générateur de tension continue ($E = \SI{12}{V}$), condensateur ($C = \SI{10}{\micro F}$), bobine idéale ($L = \SI{25}{mH}$), interrupteur à deux positions, oscilloscope à mémoire.

\textbf{Protocole :}
\begin{enumerate}
\item \textbf{Phase 1 — Charge.} On bascule $K$ en position 1 : le condensateur se charge (presque instantanément) sous la tension $E$. Il porte alors la charge $Q_0 = CE$ sur son armature reliée au pôle positif du générateur.
\item \textbf{Phase 2 — Oscillations.} À l'instant $t = 0$, on bascule $K$ en position 2 : le condensateur chargé est connecté à la bobine seule. On observe la tension $u_C$ à l'oscilloscope.
\end{enumerate}

\textbf{Observation :} la tension $u_C$ ne décroît pas vers zéro comme dans un circuit $RC$ ! Elle \emph{oscille} : elle devient alternativement positive, nulle, négative, nulle, positive\dots{} avec une période bien définie. Le condensateur se décharge, se recharge en sens inverse, se décharge à nouveau, et ainsi de suite.

\textbf{Interprétation qualitative :} le condensateur chargé fait circuler un courant dans la bobine ; celle-ci, par auto-induction, s'oppose à l'arrêt du courant et recharge le condensateur en sens inverse. Le phénomène se répète indéfiniment si rien ne dissipe l'énergie : ce sont des \cle{oscillations électriques libres non amorties}.
\end{experience}

\begin{popfigure}
\begin{center}
\begin{circuitikz}[european, scale=1.0, transform shape]
% Générateur et interrupteur
\draw (0,0) to[battery1, l_={$E$}] (0,3);
\draw (0,3) -- (1.2,3);
\draw (1.2,3) node[above left]{\small 1};
\fill (1.5,3) circle (1.5pt);
\draw (1.5,3.6) node[above]{\small 2};
\fill (1.5,3.6) circle (1.5pt);
% levier de l'interrupteur (position 2)
\draw[thick, poporange] (1.5,3.6) -- (2.6,3.15);
\fill (2.7,3.1) circle (1.5pt);
\draw (2.7,3.1) -- (3.5,3.1);
% condensateur
\draw (3.5,3.1) to[C, l_={$C$}, v=$u_C$] (3.5,0);
% armatures et charges
\node at (4.15,2.55) {$+q$};
\node at (4.15,0.55) {$-q$};
% bobine
\draw (3.5,0) -- (6.5,0) to[L, l_={$L$}, v=$u_L$] (6.5,3.1) -- (3.5,3.1);
% retour vers interrupteur
\draw (0,0) -- (3.5,0);
% sens du courant
\draw[-{Stealth[length=3mm]}, popblue, thick] (4.6,3.35) -- (5.6,3.35) node[midway, above, popblue]{$i$};
\node[popdark] at (3.3,-0.7) {\small Position 1 : charge du condensateur \quad Position 2 : oscillations $LC$};
\end{circuitikz}
\end{center}
\legende{Circuit d'étude des oscillations libres : l'interrupteur $K$ en position 1 charge le condensateur ; en position 2, le condensateur chargé est branché aux bornes de la bobine idéale. Les flèches $u_C$ et $u_L$ indiquent les tensions en convention récepteur, le courant $i$ est orienté vers l'armature portant la charge $+q$.}
\end{popfigure}

\courssub{Choix des conventions d'orientation}

Avant tout calcul, il faut \emph{orienter le circuit} : c'est l'étape que les élèves négligent trop souvent, et qui est pourtant la clé de la cohérence des signes.

\begin{definition}[Conventions adoptées pour le circuit LC]
On choisit, pour la maille formée par le condensateur et la bobine :
\begin{itemize}
\item un \cle{sens positif du courant} $i$, arbitraire mais fixé une fois pour toutes (flèche sur le schéma) ;
\item la charge $q$ étudiée est celle de \textbf{l'armature vers laquelle arrive le courant} dans le sens choisi. On a alors la relation fondamentale :
\[ i = \dfrac{dq}{dt} \]
\item les tensions $u_C$ (aux bornes du condensateur) et $u_L$ (aux bornes de la bobine) sont fléchées en \cle{convention récepteur} : la flèche tension est opposée à la flèche courant pour chaque dipôle.
\end{itemize}
\end{definition}

Avec ces conventions, les relations tension-courant des deux dipôles s'écrivent :
\[ u_C = \dfrac{q}{C} \qquad \text{et} \qquad u_L = L\,\dfrac{di}{dt} \quad (\text{bobine idéale : } r = 0). \]

\begin{attention}[L'erreur de signe la plus fréquente]
La relation $i = \dfrac{dq}{dt}$ n'est vraie que si $q$ désigne la charge de l'armature \emph{vers laquelle arrive} le courant orienté. Si tu choisis l'autre armature, ou si tu flèches le courant dans l'autre sens après avoir défini $q$, tu dois écrire $i = -\dfrac{dq}{dt}$. Un mélange des deux conventions conduit à l'équation fausse $\ddot{q} - \dfrac{q}{LC} = 0$, dont les solutions sont exponentielles (et non sinusoïdales) : le circuit « exploserait », ce qui est physiquement absurde. \textbf{Réflexe de survie :} dessine le circuit, flèche $i$, puis écris $+q$ sur l'armature d'arrivée de $i$ \emph{avant} d'écrire la moindre équation.
\end{attention}

\courssub{Établissement de l'équation différentielle}

\begin{methode}[Établir l'équation différentielle d'un circuit LC]
\begin{enumerate}
\item \textbf{Orienter} le circuit : choisir le sens de $i$, repérer l'armature dont on suit la charge $q$ (armature d'arrivée de $i$), flécher $u_C$ et $u_L$ en convention récepteur.
\item \textbf{Écrire la loi des mailles} (loi d'additivité des tensions) dans la maille $LC$ : $u_C + u_L = 0$.
\item \textbf{Remplacer} chaque tension par son expression : $u_C = \dfrac{q}{C}$ et $u_L = L\dfrac{di}{dt}$.
\item \textbf{Utiliser} $i = \dfrac{dq}{dt}$, donc $\dfrac{di}{dt} = \dfrac{d^2q}{dt^2}$, pour obtenir une équation en $q$ seul.
\item \textbf{Mettre sous forme canonique} en divisant par $L$ : $\ddot{q} + \dfrac{1}{LC}\,q = 0$.
\end{enumerate}
\end{methode}

Appliquons cette méthode. Dans la maille formée par le condensateur et la bobine (interrupteur en position 2), la loi des mailles s'écrit :
\[ u_C + u_L = 0. \]
En remplaçant les tensions par leurs expressions :
\[ \dfrac{q}{C} + L\,\dfrac{di}{dt} = 0. \]
Comme $i = \dfrac{dq}{dt}$, on a $\dfrac{di}{dt} = \ddot{q}$, d'où :
\[ L\,\ddot{q} + \dfrac{q}{C} = 0. \]

\begin{propriete}[Équation différentielle des oscillations libres — démonstration exigible]
Un circuit $LC$ idéal (bobine sans résistance) est le siège d'\cle{oscillations électriques sinusoïdales libres non amorties}. La charge $q(t)$ du condensateur obéit à l'équation différentielle :
\[ \ddot{q} + \dfrac{1}{LC}\,q = 0 \]
C'est l'équation d'un \cle{oscillateur harmonique} de \textbf{pulsation propre} :
\[ \omega_0 = \dfrac{1}{\sqrt{LC}} \qquad (\text{en } \mathrm{rad\cdot s^{-1}}). \]

\emph{Démonstration (à savoir refaire au Bac) :} voir les étapes ci-dessus — loi des mailles $u_C + u_L = 0$, substitution $u_C = q/C$ et $u_L = L\,di/dt$, puis $i = \dot{q}$. \emph{Vérification dimensionnelle :} $[LC] = [L][C]$. Or $u_L = L\,di/dt$ donne $[L] = \mathrm{V\cdot s\cdot A^{-1}}$ et $q = Cu$ donne $[C] = \mathrm{A\cdot s\cdot V^{-1}}$. Donc $[LC] = \mathrm{s^2}$, et $1/\sqrt{LC}$ s'exprime bien en $\mathrm{s^{-1}}$, homogène à une pulsation. $\blacksquare$
\end{propriete}

\begin{aretenir}[À retenir absolument]
L'équation $\ddot{q} + \omega_0^2 q = 0$ avec $\omega_0 = \dfrac{1}{\sqrt{LC}}$ gouverne la charge du condensateur. Elle est \emph{strictement analogue} à l'équation mécanique $\ddot{x} + \dfrac{k}{m}x = 0$ de l'oscillateur élastique : la charge $q$ joue le rôle de la position $x$, l'inductance $L$ celui de la masse $m$ (inertie électrique), et $1/C$ celui de la raideur $k$ (raideur électrique du condensateur).
\end{aretenir}

\courssub{Résolution de l'équation différentielle}

\textbf{Forme générale de la solution et vérification}\par

L'équation différentielle $\ddot{q} + \omega_0^2 q = 0$ admet pour solution générale :
\[ q(t) = Q_m\cos(\omega_0 t + \varphi), \]
où $Q_m > 0$ est l'\textbf{amplitude} (charge maximale, en coulombs) et $\varphi$ la \textbf{phase à l'origine} (en radians), toutes deux déterminées par les conditions initiales.

\begin{experience}[Vérification exigible : la solution proposée satisfait-elle l'équation ?]
Injectons $q(t) = Q_m\cos(\omega_0 t + \varphi)$ dans l'équation différentielle.
\begin{align*}
q(t) &= Q_m\cos(\omega_0 t + \varphi) \\
\dot{q}(t) &= -\omega_0 Q_m\sin(\omega_0 t + \varphi) \\
\ddot{q}(t) &= -\omega_0^2 Q_m\cos(\omega_0 t + \varphi) = -\omega_0^2\, q(t).
\end{align*}
En reportant dans $\ddot{q} + \dfrac{1}{LC}q$ :
\[ \ddot{q} + \dfrac{1}{LC}q = \left(\dfrac{1}{LC} - \omega_0^2\right) q(t). \]
Cette quantité est nulle \emph{pour tout} $t$ si et seulement si :
\[ \omega_0^2 = \dfrac{1}{LC} \quad \Longleftrightarrow \quad \omega_0 = \dfrac{1}{\sqrt{LC}}. \]
La solution sinusoïdale convient donc, à condition d'identifier la pulsation à la pulsation propre du circuit. $\blacksquare$
\end{experience}

\textbf{Détermination des constantes à partir des conditions initiales}\par

Dans notre expérience, à l'instant $t = 0$ où l'on bascule l'interrupteur en position 2 :
\begin{itemize}
\item le condensateur est chargé : $q(0) = Q_0 = CE$ ;
\item le courant n'a pas encore eu le temps de s'établir (la bobine s'oppose à toute variation brutale de $i$ : l'intensité dans une bobine est nécessairement continue) : $i(0) = 0$.
\end{itemize}

De $q(t) = Q_m\cos(\omega_0 t + \varphi)$ et $i(t) = \dot{q}(t) = -\omega_0 Q_m\sin(\omega_0 t + \varphi)$, on tire à $t = 0$ :
\[
\begin{cases}
q(0) = Q_m\cos\varphi = Q_0 \\
i(0) = -\omega_0 Q_m\sin\varphi = 0
\end{cases}
\quad\Longrightarrow\quad
\sin\varphi = 0 \ \text{et}\ \cos\varphi > 0 \ \Longrightarrow\ \varphi = 0,\quad Q_m = Q_0.
\]

\begin{propriete}[Charge et intensité dans le circuit LC]
Avec les conditions initiales $q(0) = Q_0$ et $i(0) = 0$, les grandeurs électriques du circuit $LC$ idéal évoluent selon :
\[ q(t) = Q_0\cos(\omega_0 t) \qquad \text{et} \qquad i(t) = -\omega_0 Q_0\sin(\omega_0 t) = -I_m\sin(\omega_0 t), \]
avec l'\textbf{amplitude de l'intensité} :
\[ I_m = \omega_0 Q_0 = \dfrac{Q_0}{\sqrt{LC}}. \]
\end{propriete}

\courssub{Quadrature de phase entre q et i}

Réécrivons l'intensité sous forme de cosinus, grâce à $-\sin\theta = \cos\left(\theta + \dfrac{\pi}{2}\right)$ :
\[ i(t) = I_m\cos\left(\omega_0 t + \dfrac{\pi}{2}\right). \]
Comparons avec $q(t) = Q_0\cos(\omega_0 t)$ : l'intensité présente une phase avancée de $\dfrac{\pi}{2}$ par rapport à la charge.

\begin{aretenir}[Déphasage entre $q$ et $i$]
Dans un circuit $LC$ idéal, $q$ et $i$ oscillent en \cle{quadrature de phase} :
\begin{itemize}
\item l'intensité $i$ est \textbf{en avance de $\dfrac{\pi}{2}$} sur la charge $q$ (ou, de façon équivalente, $q$ est en retard de $\dfrac{\pi}{2}$ sur $i$) ;
\item conséquences physiques immédiates : quand le condensateur est \emph{chargé au maximum} ($q = \pm Q_0$), le courant est \emph{nul} ; quand le condensateur est \emph{déchargé} ($q = 0$), le courant est \emph{maximal en valeur absolue} ($i = \pm I_m$).
\end{itemize}
C'est exactement le comportement du pendule : la vitesse est nulle aux extrémités (position extrême) et maximale au passage par la position d'équilibre.
\end{aretenir}

\begin{popfigure}
\begin{center}
\begin{tikzpicture}
\begin{axis}[
  width=0.95\linewidth, height=6.2cm,
  axis lines=middle,
  xlabel={$t$}, ylabel={},
  xmin=-0.05, xmax=2.15,
  ymin=-1.35, ymax=1.35,
  xtick={0,0.5,1,1.5,2},
  xticklabels={$0$,$T_0/2$,$T_0$,$3T_0/2$,$2T_0$},
  ytick={-1,1},
  yticklabels={$-Q_0$,$Q_0$},
  legend style={at={(0.5,1.02)}, anchor=south, draw=none, font=\small},
  legend columns=2,
  grid=both, grid style={popline!40, very thin},
]
\addplot[popcurve, domain=0:2, samples=200] {cos(360*x)};
\addlegendentry{$q(t) = Q_0\cos(\omega_0 t)$}
\addplot[poporange, very thick, domain=0:2, samples=200] {-sin(360*x)};
\addlegendentry{$i(t)/\omega_0 = -Q_0\sin(\omega_0 t)$}
% Mise en évidence du déphasage pi/2
\draw[{Stealth}-{Stealth}, poppurple, thick] (axis cs:0,1.18) -- (axis cs:0.25,1.18) node[midway, above, poppurple, font=\small]{$\Delta\varphi = \pi/2$};
\draw[popmuted, dashed] (axis cs:0.25,-1.3) -- (axis cs:0.25,1.25);
\draw[popmuted, dashed] (axis cs:0,-1.3) -- (axis cs:0,1.25);
\end{axis}
\end{tikzpicture}
\end{center}
\legende{Évolutions temporelles de la charge $q(t)$ (bleu) et de l'intensité $i(t)$ (orange, normalisée par $\omega_0$ pour être comparée à $q$) sur deux périodes. Le courant atteint son extremum un quart de période avant la charge : $i$ est en avance de $\pi/2$ sur $q$.}
\end{popfigure}

\begin{exemplebox}[Oscillations libres : un exemple chiffré complet]
Un condensateur de capacité $C = \SI{10}{\micro F}$ est chargé sous une tension $E = \SI{12}{V}$, puis branché à $t = 0$ aux bornes d'une bobine idéale d'inductance $L = \SI{25}{mH}$.

\textbf{Charge initiale :}
\[ Q_0 = CE = \num{10e-6} \times 12 = \num{1,2e-4}\ \mathrm{C} = \SI{120}{\micro C}. \]

\textbf{Pulsation propre :}
\[ \omega_0 = \dfrac{1}{\sqrt{LC}} = \dfrac{1}{\sqrt{\num{25e-3} \times \num{10e-6}}} = \dfrac{1}{\sqrt{\num{2,5e-7}}} = \SI{2,0e3}{rad.s^{-1}}. \]

\textbf{Expressions temporelles :}
\[ q(t) = \num{1,2e-4}\cos(2000\,t)\ \mathrm{(C)}, \qquad i(t) = -\num{0,24}\sin(2000\,t)\ \mathrm{(A)}. \]

\textbf{Amplitude du courant :}
\[ I_m = \omega_0 Q_0 = 2000 \times \num{1,2e-4} = \SI{0,24}{A}. \]

Le courant atteint donc sa valeur maximale en valeur absolue, $\SI{0,24}{A}$, chaque fois que le condensateur est totalement déchargé, soit deux fois par période (une fois dans chaque sens).
\end{exemplebox}

\begin{exoresolu}[Exploiter l'équation différentielle d'un circuit LC]
\textbf{Énoncé.} Dans un circuit $LC$ idéal, on mesure à l'oscilloscope une tension $u_C$ maximale de $\SI{6,0}{V}$ et une période des oscillations $T_0 = \SI{0,50}{ms}$. La bobine a une inductance $L = \SI{10}{mH}$.
\begin{enumerate}
\item Établir l'équation différentielle vérifiée par $u_C$.
\item En déduire la capacité $C$ du condensateur.
\item Calculer l'amplitude $I_m$ de l'intensité.
\end{enumerate}
\tcblower
\textbf{Solution détaillée.}

\textbf{1.} Loi des mailles : $u_C + u_L = 0$ avec $u_L = L\dfrac{di}{dt}$ et $i = \dfrac{dq}{dt} = C\dfrac{du_C}{dt}$. Donc :
\[ u_C + LC\,\dfrac{d^2u_C}{dt^2} = 0 \quad\Longleftrightarrow\quad \ddot{u}_C + \dfrac{1}{LC}u_C = 0. \]

\textbf{2.} La pulsation propre vaut $\omega_0 = \dfrac{2\pi}{T_0} = \dfrac{2\pi}{\num{0,50e-3}} \approx \SI{1,26e4}{rad.s^{-1}}$. De $\omega_0^2 = \dfrac{1}{LC}$ on tire :
\[ C = \dfrac{1}{L\omega_0^2} = \dfrac{1}{\num{10e-3}\times(\num{1,26e4})^2} \approx \SI{6,3e-7}{F} = \SI{0,63}{\micro F}. \]

\textbf{3.} L'amplitude de la charge est $Q_m = CU_m = \num{6,3e-7}\times 6{,}0 \approx \num{3,8e-6}\ \mathrm{C}$, d'où :
\[ I_m = \omega_0 Q_m = \num{1,26e4}\times\num{3,8e-6} \approx \SI{4,8e-2}{A} = \SI{48}{mA}. \]
\end{exoresolu}

\begin{attention}[Pièges de rédaction au Baccalauréat]
\begin{itemize}
\item Ne confonds pas $q$ (charge instantanée, variable) et $Q_0$ ou $Q_m$ (amplitudes, constantes). Écrire « $q = CE$ » à tout instant est faux : seule la charge \emph{initiale} vaut $CE$.
\item L'équation différentielle s'écrit avec un signe $+$ : $\ddot{q} + \dfrac{q}{LC} = 0$. Un signe $-$ trahit une erreur de convention d'orientation.
\item La pulsation $\omega_0$ s'exprime en $\mathrm{rad\cdot s^{-1}}$, pas en hertz : ne confonds pas $\omega_0 = \dfrac{1}{\sqrt{LC}}$ et $f_0 = \dfrac{1}{2\pi\sqrt{LC}}$ (nous y reviendrons dans la section suivante).
\item Dans $i = \dfrac{dq}{dt}$, le signe de $i(t) = -\omega_0 Q_0\sin(\omega_0 t)$ est \emph{cohérent} : juste après $t = 0$, le condensateur se décharge, $q$ diminue, donc $i < 0$ dans le sens choisi — c'est normal et attendu.
\end{itemize}
\end{attention}

\begin{savaistu}[Lord Kelvin et la naissance des oscillations électriques]
C'est le physicien britannique \textbf{William Thomson}, devenu Lord Kelvin, qui établit en 1853 la théorie mathématique des oscillations électriques d'un circuit composé d'une bobine et d'un condensateur (à l'époque, une « bouteille de Leyde »). Il montra que la décharge n'est pas un simple courant qui s'éteint, mais une succession d'oscillations dont il donna la période $T_0 = 2\pi\sqrt{LC}$ — la \emph{formule de Thomson} que tu étudies aujourd'hui. Plus de trente ans plus tard, en 1887, Heinrich Hertz produisit et détecta les premières ondes électromagnétiques grâce à un circuit oscillant de ce type : la radio était née. Aujourd'hui, chaque téléphone portable connecté aux réseaux MTN ou Orange au Cameroun contient des circuits oscillants $LC$ accordés avec une précision extrême pour sélectionner les fréquences de communication !
\end{savaistu}

En résumé, le circuit $LC$ idéal est un oscillateur harmonique électrique : la charge et l'intensité y oscillent sinusoïdalement, en quadrature de phase, à la pulsation propre $\omega_0 = \dfrac{1}{\sqrt{LC}}$. Reste une question essentielle : d'où vient l'énergie qui alimente ces oscillations éternelles, et comment se transforme-t-elle ? C'est l'objet de l'étude énergétique, après avoir examiné dans la prochaine section la période propre et ses applications concrètes.

\coursec{Période propre, fréquence et applications}

Dans la section précédente, nous avons établi que la charge du condensateur d'un circuit \cle{LC} idéal obéit à l'équation différentielle
\[
\ddot{q}+\frac{1}{LC}\,q=0,
\]
dont la solution générale est sinusoïdale : $q(t)=Q_m\cos(\omega_0 t+\varphi)$, avec $\omega_0=\dfrac{1}{\sqrt{LC}}$. Le circuit oscille donc \emph{spontanément}, sans aucune intervention extérieure, à une cadence qui ne dépend que de ses deux constituants. Cette cadence propre — la \cle{période propre} $T_0$ — est l'objet de cette section. C'est elle qui explique pourquoi un poste radio capte une station et pas une autre, et c'est elle que l'on mesure à l'oscilloscope au laboratoire.

\courssub{Période propre, fréquence propre, pulsation propre}

\textbf{Intuition.} Imagine une balançoire : sa période d'oscillation ne dépend que de sa longueur, pas de la personne assise dessus ni de l'amplitude (pour de petites oscillations). De même, un circuit LC possède un « rythme naturel » : une fois le condensateur chargé puis le circuit refermé, la charge oscille entre l'armature positive et l'armature négative avec une période parfaitement définie par $L$ et $C$. Une grosse bobine « freine » les variations de courant (inertie électrique), un gros condensateur met longtemps à se charger et se décharger (réservoir électrique) : dans les deux cas, les oscillations sont plus lentes.

\begin{propriete}[Grandeurs caractéristiques des oscillations libres (admise)]
La solution $q(t)=Q_m\cos(\omega_0 t+\varphi)$ de l'équation différentielle du circuit LC impose les relations suivantes, appelées \cle{formules de Thomson} :
\[
\omega_0=\frac{1}{\sqrt{LC}} \qquad ; \qquad T_0=\frac{2\pi}{\omega_0}=2\pi\sqrt{LC} \qquad ; \qquad f_0=\frac{1}{T_0}=\frac{1}{2\pi\sqrt{LC}}.
\]
\begin{itemize}
\item $\omega_0$ : \cle{pulsation propre}, en \si{\radian\per\second} ;
\item $T_0$ : \cle{période propre}, en secondes (\si{\second}) ;
\item $f_0$ : \cle{fréquence propre}, en hertz (\si{\hertz}).
\end{itemize}
Ces grandeurs ne dépendent \textbf{que} de $L$ (henry) et $C$ (farad), jamais de la charge initiale $Q_m$ ni de l'amplitude des oscillations.
\end{propriete}

\textbf{D'où vient le facteur $2\pi$ ?} La fonction $\cos(\omega_0 t+\varphi)$ retrouve la même valeur chaque fois que son argument augmente de $2\pi$. La période $T_0$ est donc définie par $\omega_0 T_0=2\pi$, d'où $T_0=\dfrac{2\pi}{\omega_0}=2\pi\sqrt{LC}$. Le $2\pi$ n'est pas un détail : l'oublier, c'est se tromper d'un facteur $\approx 6{,}28$ !

\begin{attention}[Ne pas confondre $\omega_0$ et $f_0$]
L'erreur la plus fréquente au Baccalauréat est de confondre :
\begin{itemize}
\item la \textbf{pulsation propre} $\omega_0=\dfrac{1}{\sqrt{LC}}$, en \si{\radian\per\second} ;
\item la \textbf{fréquence propre} $f_0=\dfrac{1}{2\pi\sqrt{LC}}$, en \si{\hertz}.
\end{itemize}
Elles sont reliées par $\omega_0=2\pi f_0$. Si l'énoncé donne une fréquence en \si{\hertz} (par exemple une station FM) et que tu dois calculer $C$, tu dois impérativement écrire $\omega_0=2\pi f_0$ \emph{avant} d'utiliser $\omega_0^2=\dfrac{1}{LC}$. Vérifie toujours l'unité demandée dans la question.
\end{attention}

\courssub{Analyse dimensionnelle : $\sqrt{LC}$ est homogène à un temps}

Vérifier l'homogénéité de la formule de Thomson est une compétence exigible. Partons des relations de définition des deux dipôles :
\[
u_L=L\,\frac{di}{dt}\quad\Longrightarrow\quad [L]=\frac{[U]\cdot[T]}{[I]},
\qquad
i=C\,\frac{du_C}{dt}\quad\Longrightarrow\quad [C]=\frac{[I]\cdot[T]}{[U]}.
\]
Le produit donne :
\[
[LC]=[L]\cdot[C]=\frac{[U]\cdot[T]}{[I]}\times\frac{[I]\cdot[T]}{[U]}=[T]^2.
\]
Donc $[\sqrt{LC}]=[T]$ : la quantité $\sqrt{LC}$ est bien homogène à un \textbf{temps}, et $T_0=2\pi\sqrt{LC}$ est cohérente ($2\pi$ est sans dimension). Retiens ce résultat : \emph{le produit $LC$ est homogène à un temps au carré}, exactement comme le rapport $m/k$ en mécanique.

\courssub{Influence de $L$ et $C$ sur la période propre}

La formule $T_0=2\pi\sqrt{LC}$ se lit directement :
\begin{itemize}
\item si l'on \textbf{augmente $L$} (bobine plus longue, noyau de fer doux inséré), $T_0$ augmente : les oscillations \textbf{ralentissent}. La bobine s'oppose plus fortement aux variations de $i$ : le courant met plus de temps à s'établir et à s'inverser ;
\item si l'on \textbf{augmente $C$}, $T_0$ augmente également : le condensateur stocke davantage de charge sous une tension donnée, sa charge et sa décharge prennent plus de temps ;
\item $T_0$ ne croît qu'en $\sqrt{\phantom{L}}$ : pour \textbf{doubler} la période, il faut \textbf{quadrupler} $L$ ou $C$.
\end{itemize}

\begin{popfigure}
\begin{center}
\begin{tikzpicture}
\begin{loglogaxis}[
  width=0.92\linewidth, height=7.2cm,
  xlabel={$C$ (F)}, ylabel={$T_0$ (s)},
  xmin=1e-12, xmax=1e-6, ymin=1e-8, ymax=1e-2,
  grid=both, grid style={popline!50},
  legend style={at={(0.03,0.97)}, anchor=north west, font=\small},
  label style={font=\small}, tick label style={font=\small}
]
\addplot[popcurve, poporange, domain=1e-12:1e-6, samples=100] {6.2832*sqrt(max(0,1e-4*x))};
\addlegendentry{$L=\SI{0,10}{mH}$}
\addplot[popcurve, popgreen, domain=1e-12:1e-6, samples=100] {6.2832*sqrt(max(0,1e-2*x))};
\addlegendentry{$L=\SI{10}{mH}$}
\addplot[popcurve, popblue, domain=1e-12:1e-6, samples=100] {6.2832*sqrt(max(0,1*x))};
\addlegendentry{$L=\SI{1}{H}$}
\end{loglogaxis}
\end{tikzpicture}
\end{center}
\legende{Période propre $T_0=2\pi\sqrt{LC}$ en fonction de la capacité $C$ (échelles logarithmiques) pour trois valeurs de l'inductance $L$. En échelle log-log, les courbes sont des droites de pente $1/2$ : $T_0$ croît comme $\sqrt{C}$, et chaque valeur de $L$ translate la droite vers le haut.}
\end{popfigure}

\courssub{Application : le circuit d'accord d'un poste radio}

Chaque station de radio émet une onde électromagnétique à une fréquence bien précise (par exemple \SI{94,5}{MHz} pour une station FM à Yaoundé). L'antenne d'un récepteur reçoit \emph{toutes} ces ondes simultanément : comment isoler une seule station ?

La réponse est le \cle{circuit d'accord} : un circuit LC dont on règle la fréquence propre $f_0$ pour qu'elle coïncide avec la fréquence $f$ de la station désirée. Lorsque $f_0=f$, le circuit entre en \emph{résonance} avec l'onde reçue : l'amplitude des oscillations devient maximale pour cette station, et négligeable pour les autres. On dit que le circuit est \textbf{accordé} (ou « syntonisé ») sur la station.

Comme il est techniquement difficile de fabriquer une bobine variable, on fixe $L$ et l'on fait varier $C$ à l'aide d'un \cle{condensateur variable}. La condition d'accord s'écrit :
\[
f_0=\frac{1}{2\pi\sqrt{LC}}=f
\quad\Longrightarrow\quad
\boxed{\;C=\frac{1}{4\pi^2 f^2 L}\;}
\]

\begin{popfigure}
\begin{center}
\begin{circuitikz}[european, scale=0.95, transform shape]
\% Antenne
\draw (0,0) -- (0,1.2);
\draw[thick, popink] (0,1.2) -- (-0.5,2.1);
\draw[thick, popink] (0,1.2) -- (0,2.3);
\draw[thick, popink] (0,1.2) -- (0.5,2.1);
\node[font=\small, popink] at (0,2.7) {antenne};
\% Circuit LC
\draw (0,0) to[L, l_={$L$}] (3.2,0);
\draw (3.2,0) to[vC, l_={$C$} variable, v<=$u_C$] (3.2,-2.6);
\draw (3.2,-2.6) -- (0,-2.6) -- (0,0);
\% Sortie vers amplificateur
\draw (3.2,-1.3) to[short, -o] (4.9,-1.3);
\draw (0,-1.3) to[short, -o] (-1.6,-1.3);
\node[font=\small, right, popdark] at (4.9,-1.3) {vers l'amplificateur};
\node[font=\small, left, popdark] at (-1.6,-1.3) {masse};
\% Onde reçue
\draw[decorate, decoration={snake, amplitude=1.2mm, segment length=4mm}, popblue, thick] (-2.6,1.9) -- (-0.9,1.9);
\node[font=\small, popblue] at (-1.75,2.35) {onde radio $f$};
\end{circuitikz}
\end{center}
\legende{Schéma de principe d'un circuit d'accord radio : l'onde reçue par l'antenne induit des oscillations dans le circuit LC. En tournant le condensateur variable, on ajuste $C$ pour que $f_0=\dfrac{1}{2\pi\sqrt{LC}}$ égale la fréquence $f$ de la station : le signal correspondant est alors sélectionné et envoyé vers l'amplificateur.}
\end{popfigure}

\begin{exemplebox}[Accord sur une station FM]
On veut capter une station émettant à $f_0=\SI{94,5}{MHz}$ avec une bobine d'inductance $L=\SI{0,10}{\micro H}=\num{1,0e-7}\ \SI{}{H}$. Quelle capacité $C$ faut-il ?
\begin{align*}
C&=\frac{1}{4\pi^2 f_0^2 L}
=\frac{1}{4\pi^2\times(9,45 \times 10^7)^2\times 1,0 \times 10^{-7}}\\
&=\frac{1}{4\pi^2\times 8,93 \times 10^{15}\times 1,0 \times 10^{-7}}
=\frac{1}{3,53 \times 10^{10}}
\approx 2,8 \times 10^{-11}\ \SI{}{F}.
\end{align*}
Il faut donc $C\approx \SI{28}{pF}$. Remarque l'ordre de grandeur : les fréquences FM ($\approx \SI{100}{MHz}$) s'obtiennent avec des inductances de l'ordre de \SI{0,1}{\micro H} et des capacités de quelques dizaines de picofarads — des composants minuscules, ce qui explique la taille réduite des récepteurs modernes.
\end{exemplebox}

\begin{savaistu}[Le condensateur variable des postes radio à cadran]
Les anciens postes radio, comme ceux que l'on trouve encore dans les familles camerounaises, possédaient un bouton de recherche de stations relié à un \textbf{condensateur variable à lames mobiles} : deux jeux de lames métalliques en demi-cercles s'interpénétraient plus ou moins quand on tournait le bouton, ce qui modifiait la surface des armatures en regard, donc la capacité $C$, donc la fréquence d'accord $f_0$. Tourner le bouton, c'était littéralement « parcourir » la formule de Thomson ! Aujourd'hui, les récepteurs modernes utilisent des diodes à capacité variable (\emph{varicap}) commandées électriquement, mais le principe physique est identique.
\end{savaistu}

\courssub{Méthode de résolution et lien avec la mécanique}

\begin{methode}[Déterminer une grandeur inconnue ($L$, $C$ ou $T_0$)]
\begin{enumerate}
\item \textbf{Convertir} toutes les données en unités SI : $L$ en henry, $C$ en farad (attention aux \SI{}{pF}, \SI{}{nF}, \SI{}{\micro H}, \SI{}{mH}), $f_0$ en hertz.
\item \textbf{Passer par la pulsation} : écrire $\omega_0=\dfrac{1}{\sqrt{LC}}$, soit $\omega_0^2=\dfrac{1}{LC}$.
\item \textbf{Relier} $\omega_0$ à la grandeur temporelle connue : $\omega_0=\dfrac{2\pi}{T_0}=2\pi f_0$.
\item \textbf{Isoler} l'inconnue : $L=\dfrac{1}{\omega_0^2 C}$, $C=\dfrac{1}{\omega_0^2 L}$ ou $T_0=\dfrac{2\pi}{\omega_0}$.
\item \textbf{Vérifier} l'ordre de grandeur et l'unité du résultat (analyse dimensionnelle : $\sqrt{LC}$ en secondes).
\end{enumerate}
\end{methode}

\textbf{Lien avec la mécanique.} Tu as déjà rencontré cette structure mathématique : la période propre du pendule élastique (masse $m$ accrochée à un ressort de raideur $k$) est
\[
T_0=2\pi\sqrt{\frac{m}{k}}.
\]
L'analogie est frappante : la masse $m$ (inertie mécanique) joue le rôle de l'inductance $L$ (inertie électrique), et l'inverse de la raideur $1/k$ (souplesse du ressort) joue le rôle de la capacité $C$ (capacité à stocker la charge). Dans les deux cas, la période croît avec l'« inertie » du système et avec sa « souplesse ». Cette analogie sera formalisée dans la dernière section de la leçon.

\courssub{Mesure expérimentale de $T_0$ à l'oscilloscope (complément utile au Bac)}

Au laboratoire, on visualise la tension $u_C(t)$ aux bornes du condensateur d'un circuit LC (en pratique RLC faiblement amorti) à l'oscilloscope. La courbe est une sinusoïde (pseudo-périodique si l'amortissement est visible). Pour mesurer $T_0$ :
\begin{enumerate}
\item on compte le nombre $n$ de divisions horizontales correspondant à \emph{une} période (par exemple entre deux maxima consécutifs) ;
\item on multiplie par le balayage $b$ (en \si{\second\per\division}) : $T_0=n\times b$ ;
\item pour gagner en précision, on mesure plusieurs périodes : $T_0=\dfrac{n\times b}{N}$ pour $N$ périodes.
\end{enumerate}
Connaissant $T_0$ et l'une des deux grandeurs $L$ ou $C$, on en déduit l'autre par la formule de Thomson : c'est une méthode classique de \textbf{mesure indirecte d'une inductance ou d'une capacité inconnue}.

\begin{exoresolu}[Détermination d'une inductance inconnue]
Un condensateur de capacité $C=\SI{10}{nF}$, initialement chargé, est connecté à une bobine d'inductance $L$ inconnue et de résistance négligeable. À l'oscilloscope, on observe que $N=5$ périodes complètes occupent $n=6,3$ divisions horizontales, avec un balayage $b=\SI{20}{\micro\second}$.
\begin{enumerate}
\item Déterminer la période propre $T_0$ des oscillations.
\item En déduire l'inductance $L$ de la bobine.
\end{enumerate}
\tcblower
\textbf{1. Période propre.} La durée totale balayée par 5 périodes est $n\times b=6,3\times\SI{20}{\micro\second}=\SI{126}{\micro\second}$. Donc :
\[
T_0=\frac{n\times b}{N}=\frac{\SI{126}{\micro\second}}{5}=\SI{25,2}{\micro\second}=2,52 \times 10^{-5}\ \SI{}{s}.
\]
\textbf{2. Inductance.} De $T_0=2\pi\sqrt{LC}$, on tire $T_0^2=4\pi^2 LC$, soit :
\[
L=\frac{T_0^2}{4\pi^2 C}
=\frac{(2,52 \times 10^{-5})^2}{4\pi^2\times 1,0 \times 10^{-8}}
=\frac{6,35 \times 10^{-10}}{3,95 \times 10^{-7}}
\approx 1,6 \times 10^{-3}\ \SI{}{H}.
\]
La bobine a donc une inductance $L\approx\SI{1,6}{mH}$. Vérification dimensionnelle : $\dfrac{[T]^2}{[C]}=\dfrac{[T]^2\,[U]}{[I]\,[T]}=\dfrac{[U]\,[T]}{[I]}=[L]$, conforme à $u=L\,\dfrac{di}{dt}$. Le résultat est bien homogène à une inductance.
\end{exoresolu}

\begin{aretenir}[L'essentiel de la section]
\begin{itemize}
\item Formules de Thomson : $\omega_0=\dfrac{1}{\sqrt{LC}}$, $T_0=2\pi\sqrt{LC}$, $f_0=\dfrac{1}{2\pi\sqrt{LC}}$.
\item Le produit $LC$ est homogène à un \textbf{temps au carré}.
\item Augmenter $L$ ou $C$ \textbf{ralentit} les oscillations ; $T_0$ ne dépend pas de l'amplitude.
\item Le circuit d'accord radio exploite la résonance : on règle $C$ pour que $f_0$ coïncide avec la fréquence de la station, soit $C=\dfrac{1}{4\pi^2 f^2 L}$.
\item Analogie avec le pendule élastique : $T_0=2\pi\sqrt{m/k}$ ; correspondances $L\leftrightarrow m$ et $C\leftrightarrow 1/k$.
\item À l'oscilloscope : $T_0=\dfrac{n\times b}{N}$, puis on remonte à $L$ ou $C$ par la formule de Thomson.
\end{itemize}
\end{aretenir}

Maintenant que nous connaissons la cadence propre du circuit LC, une question naturelle se pose : que devient l'énergie initialement stockée dans le condensateur au cours de ces oscillations ? C'est l'objet de la section suivante, consacrée à l'étude énergétique du circuit LC et à la conservation de l'énergie électromagnétique.

\coursec{Étude énergétique du circuit LC : conservation de l'énergie électromagnétique}

Dans la section précédente, nous avons établi que la charge du condensateur et l'intensité du courant dans un circuit LC oscillent sinusoïdalement avec la période propre $T_0 = 2\pi\sqrt{LC}$. Mais que devient l'\emph{énergie} pendant ces oscillations ? C'est la question à laquelle cette section répond — et la réponse est d'une beauté remarquable : l'énergie totale du circuit reste \cle{constante}, elle ne cesse de passer d'une forme à l'autre, comme l'eau qui s'écoule entre deux bassins communiquants sans jamais se perdre.

\courssub{Les deux réservoirs d'énergie du circuit}

\textbf{Énergie électrique stockée dans le condensateur}\par

Charger un condensateur, c'est arracher des électrons d'une armature pour les accumuler sur l'autre : il faut lutter contre la répulsion électrique, donc \emph{fournir un travail}. Ce travail n'est pas perdu : il est stocké dans le condensateur sous forme d'énergie \cle{électrique} (on dit aussi électrostatique), récupérable lors de la décharge.

\begin{definition}[Énergie électrique d'un condensateur]
Un condensateur de capacité $C$, portant la charge $q$ sous la tension $u_C = q/C$, stocke l'énergie électrique :
\[
E_C = \frac{1}{2}\,\frac{q^2}{C} = \frac{1}{2}\,C\,u_C^2 = \frac{1}{2}\,q\,u_C
\]
$E_C$ en joules (J), $q$ en coulombs (C), $C$ en farads (F), $u_C$ en volts (V).
\end{definition}

\textbf{Énergie magnétique stockée dans la bobine}\par

De même, établir un courant dans une bobine exige un travail contre la force contre-électromotrice d'auto-induction. Ce travail est stocké dans la bobine sous forme d'énergie \cle{magnétique}, associée au champ magnétique $\vec{B}$ créé par le courant.

\begin{definition}[Énergie magnétique d'une bobine]
Une bobine d'inductance $L$ parcourue par un courant d'intensité $i$ stocke l'énergie magnétique :
\[
E_L = \frac{1}{2}\,L\,i^2
\]
$E_L$ en joules (J), $L$ en henrys (H), $i$ en ampères (A).
\end{definition}

\begin{attention}[Analogie à retenir]
Ces deux expressions ont exactement la même structure que l'énergie potentielle élastique $\frac{1}{2}kx^2$ et l'énergie cinétique $\frac{1}{2}mv^2$ du système masse--ressort : le condensateur joue le rôle du \cle{ressort} (réservoir d'énergie « potentielle »), la bobine celui de la \cle{masse} (réservoir d'énergie « cinétique »). Nous y reviendrons en détail en fin de leçon.
\end{attention}

\begin{definition}[Énergie électromagnétique totale]
L'énergie électromagnétique totale du circuit LC est la somme des deux formes d'énergie :
\[
E = E_C + E_L = \frac{q^2}{2C} + \frac{1}{2}Li^2
\]
\end{definition}

\courssub{Démonstration de la conservation de l'énergie}

Voici le résultat central de cette section, dont la démonstration est \textbf{exigible au Baccalauréat}.

\begin{propriete}[Conservation de l'énergie électromagnétique — démonstration exigible]
Dans un circuit LC \emph{idéal} (bobine de résistance $r$ nulle), l'énergie électromagnétique totale se conserve au cours du temps :
\[
\frac{dE}{dt} = 0 \qquad \Longleftrightarrow \qquad E = \text{constante}
\]
Sa valeur est entièrement fixée par les conditions initiales : si le condensateur est initialement chargé ($q(0) = Q_m$) et le courant nul ($i(0) = 0$), alors
\[
E = \frac{Q_m^2}{2C} = \frac{1}{2}C\,U_m^2 = \frac{1}{2}L\,I_m^2
\]
où $I_m = \omega_0 Q_m$ est l'intensité maximale.
\end{propriete}

\begin{experience}[Démonstration guidée : montrer que $\dfrac{dE}{dt} = 0$]
Suivons la démarche pas à pas, telle qu'elle doit être rédigée.
\begin{enumerate}
\item \textbf{Écrire l'énergie totale} : $E = \dfrac{q^2}{2C} + \dfrac{1}{2}Li^2$.
\item \textbf{Dériver par rapport au temps} ($C$ et $L$ sont des constantes) :
\[
\frac{dE}{dt} = \frac{1}{2C}\cdot 2q\,\frac{dq}{dt} + \frac{1}{2}L\cdot 2i\,\frac{di}{dt} = \frac{q}{C}\,\dot{q} + L\,i\,\dot{i}
\]
\item \textbf{Exploiter les relations du circuit} : $i = \dot{q}$ donc $\dot{i} = \ddot{q}$. L'équation différentielle du circuit LC s'écrit $\ddot{q} + \dfrac{q}{LC} = 0$, soit $\ddot{q} = -\dfrac{q}{LC}$.
\item \textbf{Substituer} :
\[
\frac{dE}{dt} = \frac{q}{C}\,i + L\,i\left(-\frac{q}{LC}\right) = \frac{q\,i}{C} - \frac{q\,i}{C} = 0
\]
\end{enumerate}
\textbf{Conclusion} : $\dfrac{dE}{dt} = 0$, donc $E$ est une constante du mouvement. L'énergie électromagnétique se conserve dans un circuit LC idéal.
\end{experience}

\begin{methode}[Rédaction type attendue au Bac : montrer une conservation d'énergie]
\begin{enumerate}
\item Écrire l'expression de l'énergie totale $E = \dfrac{q^2}{2C} + \dfrac{1}{2}Li^2$.
\item Dériver $E$ par rapport à $t$ en utilisant la dérivation des fonctions composées : $\dfrac{d(q^2)}{dt} = 2q\dot{q}$.
\item Remplacer $\dot{q}$ par $i$ et $\dot{i} = \ddot{q}$ par $-\dfrac{q}{LC}$ (équation différentielle).
\item Simplifier et conclure : $\dfrac{dE}{dt} = 0 \Rightarrow E = \text{constante}$.
\item Préciser la valeur de la constante à l'aide des conditions initiales.
\end{enumerate}
\end{methode}

\courssub{Vérification par les expressions temporelles}

Retrouvons ce résultat à partir des solutions établies dans la section 2. Avec les conditions initiales $q(0) = Q_m$ et $i(0) = 0$ :
\[
q(t) = Q_m\cos(\omega_0 t) \qquad \text{et} \qquad i(t) = \dot{q} = -\omega_0 Q_m\sin(\omega_0 t) = -I_m\sin(\omega_0 t)
\]
avec $I_m = \omega_0 Q_m$. Calculons chaque énergie :
\begin{align*}
E_C(t) &= \frac{q^2}{2C} = \frac{Q_m^2}{2C}\cos^2(\omega_0 t)\\
E_L(t) &= \frac{1}{2}Li^2 = \frac{1}{2}L\omega_0^2 Q_m^2\sin^2(\omega_0 t)
\end{align*}
Or $\omega_0^2 = \dfrac{1}{LC}$, donc $\dfrac{1}{2}L\omega_0^2 Q_m^2 = \dfrac{1}{2}L\cdot\dfrac{1}{LC}\cdot Q_m^2 = \dfrac{Q_m^2}{2C}$. Les deux énergies ont donc la \emph{même valeur maximale} $E_{\max} = \dfrac{Q_m^2}{2C}$, et :
\[
E(t) = \frac{Q_m^2}{2C}\bigl[\cos^2(\omega_0 t) + \sin^2(\omega_0 t)\bigr] = \frac{Q_m^2}{2C} = \text{constante}
\]
grâce à l'identité fondamentale $\cos^2\theta + \sin^2\theta = 1$. La conservation est vérifiée explicitement.

\begin{attention}[Piège classique : la période des échanges énergétiques]
Erreur fréquente au Bac : affirmer que $E_C$ et $E_L$ oscillent avec la période $T_0$. \textbf{C'est faux !} Les fonctions $\cos^2(\omega_0 t)$ et $\sin^2(\omega_0 t)$ ont pour période $\dfrac{T_0}{2}$, car $\cos^2(\omega_0 t) = \dfrac{1 + \cos(2\omega_0 t)}{2}$ : la pulsation des énergies est $2\omega_0$. Les échanges d'énergie entre condensateur et bobine se font donc avec la période $\boxed{T_E = \dfrac{T_0}{2}}$, alors que $q$ et $i$ oscillent avec la période $T_0$. Retiens : en une période $T_0$ des oscillations électriques, l'énergie fait \emph{deux} allers-retours complets.
\end{attention}

\courssub{Description des échanges d'énergie au cours d'une période}

Suivons le circuit instant par instant, en partant d'un condensateur chargé et d'un courant nul. On définit le sens positif du courant $i$ dans le sens anti-horaire (du $+$ du condensateur vers la bobine, puis vers le $-$ du condensateur).

\begin{itemize}
\item \textbf{À $t = 0$} : $q = Q_m$, $i = 0$. Toute l'énergie est \emph{électrique} : $E_C = \dfrac{Q_m^2}{2C}$, $E_L = 0$. Le condensateur est le « réservoir plein ».
\item \textbf{De $0$ à $T_0/4$} : le condensateur se décharge dans la bobine ; le courant croît dans le sens horaire (opposé au sens positif). $E_C$ diminue, $E_L$ augmente : l'énergie électrique se \emph{transforme} en énergie magnétique.
\item \textbf{À $t = T_0/4$} : $q = 0$, $i = -I_m$ (courant horaire maximal). Toute l'énergie est \emph{magnétique} : $E_L = \dfrac{1}{2}LI_m^2 = \dfrac{Q_m^2}{2C}$, $E_C = 0$.
\item \textbf{De $T_0/4$ à $T_0/2$} : le courant persiste (par auto-induction, la bobine « s'oppose » à sa disparition) et recharge le condensateur en sens inverse. $E_L$ rediminue, $E_C$ croît.
\item \textbf{À $t = T_0/2$} : $q = -Q_m$, $i = 0$. L'énergie est de nouveau entièrement électrique, le condensateur étant chargé en sens opposé.
\item \textbf{À $t = 3T_0/4$} : $q = 0$, $i = +I_m$ (courant anti-horaire maximal) : énergie entièrement magnétique. Puis le cycle recommence.
\end{itemize}

L'interprétation mécanique est éclairante : le condensateur joue le rôle du \cle{ressort} (il stocke une énergie de type « potentielle », liée à un état de charge, comme le ressort stocke $\frac{1}{2}kx^2$ liée à l'allongement), et la bobine celui de la \cle{masse} (elle stocke une énergie de type « cinétique », liée au courant, comme la masse stocke $\frac{1}{2}mv^2$ liée à la vitesse). L'oscillation électrique est le va-et-vient exact de l'oscillation mécanique : décharge $\leftrightarrow$ détente du ressort, courant maximal $\leftrightarrow$ vitesse maximale.

\begin{popfigure}
\begin{tikzpicture}[scale=0.92]
% Vignette 1 : t = 0
\begin{scope}[shift={(0,0)}]
\draw[popblue, thick] (-1.1,-0.5) rectangle (1.1,0.9);
% Condensateur
\draw[thick] (-0.5,0.2) -- (-0.12,0.2); \draw[thick] (0.12,0.2) -- (0.5,0.2);
\draw[thick] (-0.12,0.05) -- (-0.12,0.35); \draw[thick] (0.12,0.05) -- (0.12,0.35);
% Fils et bobine (branche basse)
\draw[thick] (-0.5,0.2) -- (-0.5,-0.2) -- (0.1,-0.2);
\draw[thick, decoration={coil, segment length=4pt, amplitude=4pt}, decorate] (0.1,-0.2) -- (0.5,-0.2);
\draw[thick] (0.5,-0.2) -- (0.5,0.2);
% Charges
\node[popink] at (-0.35,0.55) {\scriptsize $+Q_m$};
\node[popink] at (0.38,0.55) {\scriptsize $-Q_m$};
% Sens positif du courant (référence)
\draw[popmuted, dashed, -{Stealth[length=2mm]}] (-0.8,0.6) arc (180:-180:0.5) node[pos=0.5, above] {\scriptsize $i>0$};
\node[popmuted] at (0,-0.45) {\scriptsize $i=0$};
\node[font=\scriptsize\bfseries, popdark] at (0,-0.75) {$t=0$ : $E=E_C$};
\end{scope}
% Vignette 2 : t = T0/4
\begin{scope}[shift={(3.4,0)}]
\draw[poporange, thick] (-1.1,-0.5) rectangle (1.1,0.9);
\draw[thick] (-0.5,0.2) -- (-0.12,0.2); \draw[thick] (0.12,0.2) -- (0.5,0.2);
\draw[thick] (-0.12,0.05) -- (-0.12,0.35); \draw[thick] (0.12,0.05) -- (0.12,0.35);
\draw[thick] (-0.5,0.2) -- (-0.5,-0.2) -- (0.1,-0.2);
\draw[thick, decoration={coil, segment length=4pt, amplitude=4pt}, decorate] (0.1,-0.2) -- (0.5,-0.2);
\draw[thick] (0.5,-0.2) -- (0.5,0.2);
\node[popmuted] at (0,0.55) {\scriptsize $q=0$};
% Courant physique horaire (i < 0)
\draw[-{Stealth[length=3mm]}, poporange, very thick] (0.3,-0.35) arc (0:-180:0.5);
\node[poporange] at (0,-0.5) {\scriptsize $|i|=I_m\ (i=-I_m)$};
\node[font=\scriptsize\bfseries, popdark] at (0,-0.75) {$t=T_0/4$ : $E=E_L$};
\end{scope}
% Vignette 3 : t = T0/2
\begin{scope}[shift={(6.8,0)}]
\draw[popblue, thick] (-1.1,-0.5) rectangle (1.1,0.9);
\draw[thick] (-0.5,0.2) -- (-0.12,0.2); \draw[thick] (0.12,0.2) -- (0.5,0.2);
\draw[thick] (-0.12,0.05) -- (-0.12,0.35); \draw[thick] (0.12,0.05) -- (0.12,0.35);
\draw[thick] (-0.5,0.2) -- (-0.5,-0.2) -- (0.1,-0.2);
\draw[thick, decoration={coil, segment length=4pt, amplitude=4pt}, decorate] (0.1,-0.2) -- (0.5,-0.2);
\draw[thick] (0.5,-0.2) -- (0.5,0.2);
\node[popink] at (-0.35,0.55) {\scriptsize $-Q_m$};
\node[popink] at (0.38,0.55) {\scriptsize $+Q_m$};
\node[popmuted] at (0,-0.45) {\scriptsize $i=0$};
\node[font=\scriptsize\bfseries, popdark] at (0,-0.75) {$t=T_0/2$ : $E=E_C$};
\end{scope}
% Vignette 4 : t = 3T0/4
\begin{scope}[shift={(10.2,0)}]
\draw[poporange, thick] (-1.1,-0.5) rectangle (1.1,0.9);
\draw[thick] (-0.5,0.2) -- (-0.12,0.2); \draw[thick] (0.12,0.2) -- (0.5,0.2);
\draw[thick] (-0.12,0.05) -- (-0.12,0.35); \draw[thick] (0.12,0.05) -- (0.12,0.35);
\draw[thick] (-0.5,0.2) -- (-0.5,-0.2) -- (0.1,-0.2);
\draw[thick, decoration={coil, segment length=4pt, amplitude=4pt}, decorate] (0.1,-0.2) -- (0.5,-0.2);
\draw[thick] (0.5,-0.2) -- (0.5,0.2);
\node[popmuted] at (0,0.55) {\scriptsize $q=0$};
% Courant physique anti-horaire (i > 0)
\draw[-{Stealth[length=3mm]}, poporange, very thick] (-0.4,-0.35) arc (180:0:0.5);
\node[poporange] at (0,-0.5) {\scriptsize $|i|=I_m\ (i=+I_m)$};
\node[font=\scriptsize\bfseries, popdark] at (0,-0.75) {$t=3T_0/4$ : $E=E_L$};
\end{scope}
\end{tikzpicture}
\legende{Quatre états du circuit LC au cours d'une période. Le sens positif de $i$ (flèche pointillée, $t=0$) est anti-horaire. L'énergie passe alternativement du condensateur (charges) à la bobine (courant), avec un aller-retour complet toutes les $T_0/2$.}
\end{popfigure}

\courssub{Exploitation graphique : courbes $E_C(t)$, $E_L(t)$ et $E(t)$}

\begin{popfigure}
\begin{tikzpicture}
\begin{axis}[
  width=\linewidth, height=6.5cm,
  xlabel={$t$}, ylabel={Énergie (J)},
  xmin=0, xmax=1.02, ymin=0, ymax=1.25,
  xtick={0,0.25,0.5,0.75,1},
  xticklabels={$0$,$T_0/4$,$T_0/2$,$3T_0/4$,$T_0$},
  ytick={0.5,1},
  yticklabels={$E/2$,$E$},
  grid=both, grid style={popline!40},
  legend style={at={(0.5,1.02)}, anchor=south, legend columns=3, font=\scriptsize, draw=none}
]
\addplot[popblue, very thick, domain=0:1, samples=200] {cos(360*x)^2};
\addlegendentry{$E_C(t)=\frac{Q_m^2}{2C}\cos^2(\omega_0 t)$}
\addplot[poporange, very thick, domain=0:1, samples=200] {sin(360*x)^2};
\addlegendentry{$E_L(t)=\frac{Q_m^2}{2C}\sin^2(\omega_0 t)$}
\addplot[popgreen, very thick, dashed, domain=0:1, samples=10] {1};
\addlegendentry{$E=E_C+E_L=$ cte}
\draw[popmuted, dashed] (axis cs:0.25,0) -- (axis cs:0.25,1.15);
\draw[popmuted, dashed] (axis cs:0.5,0) -- (axis cs:0.5,1.15);
\draw[popmuted, dashed] (axis cs:0.75,0) -- (axis cs:0.75,1.15);
\node[popblue, font=\scriptsize] at (axis cs:0.02,1.08) [anchor=west] {tout électrique};
\node[poporange, font=\scriptsize] at (axis cs:0.27,1.08) [anchor=west] {tout magnétique};
\end{axis}
\end{tikzpicture}
\legende{Échanges d'énergie dans le circuit LC idéal. $E_C$ et $E_L$ oscillent en opposition de phase avec la période $T_0/2$ ; leur somme $E$ reste constante. Les deux courbes se croisent à $E/2$ aux instants $T_0/8$, $3T_0/8$, \dots}
\end{popfigure}

Lecture du graphique : les deux courbes sont des « bosses » de période $T_0/2$, décalées d'un quart de période électrique ; quand l'une est maximale, l'autre est nulle, et elles se croisent toujours à la moitié de l'énergie totale. La droite horizontale verte matérialise la conservation : à tout instant, $E_C + E_L = E$.

\begin{exoresolu}[Calcul d'énergie dans un circuit LC]
\textbf{Énoncé.} Un condensateur de capacité $C = \SI{10}{\micro F}$ est chargé sous la tension $U_0 = \SI{12}{V}$, puis connecté à l'instant $t = 0$ à une bobine idéale d'inductance $L = \SI{50}{mH}$.
\begin{enumerate}
\item Calculer la charge initiale $Q_0$ et l'énergie électromagnétique $E$ du circuit.
\item En déduire l'intensité maximale $I_m$ et retrouver $E$ à partir de $\frac{1}{2}LI_m^2$.
\item Quelle est l'énergie magnétique à l'instant où $q = Q_0/2$ ?
\end{enumerate}
\tcblower
\textbf{Solution détaillée.}
\begin{enumerate}
\item Charge initiale : $Q_0 = CU_0 = \num{10e-6} \times 12 = \num{1,2e-4}\ \text{C} = \SI{0,12}{mC}$.

À $t = 0$, le courant est nul : toute l'énergie est électrique.
\[
E = \frac{Q_0^2}{2C} = \frac{(\num{1,2e-4})^2}{2\times \num{10e-6}} = \frac{\num{1,44e-8}}{\num{2e-5}} = \SI{7,2e-4}{J}
\]
\item Par conservation de l'énergie, quand $E_L$ est maximale, $E_L = E$ :
\[
\frac{1}{2}LI_m^2 = E \quad\Longrightarrow\quad I_m = \sqrt{\frac{2E}{L}} = \sqrt{\frac{2\times\num{7,2e-4}}{\num{50e-3}}} = \sqrt{\num{2,88e-2}} \approx \SI{0,17}{A}
\]
Vérification : $\frac{1}{2}LI_m^2 = \frac{1}{2}\times\num{50e-3}\times(0{,}17)^2 \approx \num{7,2e-4}\ \text{J} = E$. \checkmark

(On pouvait aussi calculer $\omega_0 = 1/\sqrt{LC} \approx \SI{1,41e3}{rad.s^{-1}}$ puis $I_m = \omega_0 Q_0 \approx \SI{0,17}{A}$.)
\item Conservation de l'énergie : $E = \dfrac{q^2}{2C} + E_L$. Avec $q = Q_0/2$ :
\[
E_L = E - \frac{(Q_0/2)^2}{2C} = E - \frac{E}{4} = \frac{3E}{4} = \SI{5,4e-4}{J}
\]
L'énergie est alors répartie : $25\ \%$ électrique, $75\ \%$ magnétique.
\end{enumerate}
\end{exoresolu}

\begin{exemplebox}[Analyse dimensionnelle : une sécurité avant tout calcul]
Vérifions l'homogénéité de $E = \dfrac{q^2}{2C}$. D'après $u_C = q/C$, on a $[C] = \dfrac{[q]}{[u]}$, donc $\left[\dfrac{q^2}{C}\right] = [q]\,[u]$. Or le travail électrique s'écrit $W = q\,U$ : $[q]\,[u] = \text{J}$. \checkmark\ De même, avec $u_L = L\,\dfrac{di}{dt}$, $[L] = \dfrac{[u]\,[t]}{[i]}$, donc $[Li^2] = [u]\,[i]\,[t] = \text{W}\times\text{s} = \text{J}$. \checkmark\ Les deux termes de $E$ sont bien des énergies.
\end{exemplebox}

\begin{savaistu}[Des oscillations LC aux ondes radio]
C'est exactement ce va-et-vient d'énergie entre condensateur et bobine qui, couplé à une antenne, permet d'\emph{émettre des ondes électromagnétiques}. Un circuit LC ouvert sur une antenne rayonne une petite partie de son énergie à chaque oscillation : c'est le principe des premiers émetteurs radio (Hertz, 1887) et, aujourd'hui encore, des étages d'accord des émetteurs qui alimentent les antennes de radio et de télévision, comme celles de la CRTV ou des relais MTN et Orange disséminés sur le territoire camerounais. Choisir $L$ et $C$, c'est choisir la fréquence d'émission $f_0 = \dfrac{1}{2\pi\sqrt{LC}}$ : quand tu tournes le bouton d'un poste radio, tu modifies en réalité une capacité $C$ pour accorder le circuit LC du récepteur sur la fréquence de ta station préférée !
\end{savaistu}

\begin{aretenir}[L'essentiel de la section]
\begin{itemize}
\item Condensateur : $E_C = \dfrac{q^2}{2C} = \dfrac{1}{2}Cu_C^2$ ; bobine : $E_L = \dfrac{1}{2}Li^2$.
\item Énergie électromagnétique totale : $E = E_C + E_L$.
\item Dans un circuit LC \textbf{idéal}, $E$ se conserve : $\dfrac{dE}{dt} = 0$ (démonstration exigible : dériver puis utiliser $\ddot{q} = -\dfrac{q}{LC}$).
\item Valeur de la constante : $E = \dfrac{Q_m^2}{2C} = \dfrac{1}{2}CU_m^2 = \dfrac{1}{2}LI_m^2$.
\item Les échanges d'énergie ont lieu avec la période $\dfrac{T_0}{2}$ (et non $T_0$ !) : à $t=0$ tout est électrique, à $t = T_0/4$ tout est magnétique, etc.
\item Analogie mécanique : condensateur $\leftrightarrow$ ressort, bobine $\leftrightarrow$ masse.
\end{itemize}
\end{aretenir}

\coursec{Le circuit RLC série : oscillations amorties}

Dans la section précédente, nous avons établi un résultat remarquable : dans un circuit \cle{LC} idéal, l'énergie électromagnétique totale $E = \dfrac{q^2}{2C} + \dfrac{1}{2}Li^2$ se conserve rigoureusement. Les oscillations électriques se poursuivent alors indéfiniment, avec une amplitude constante et une période propre $T_0 = 2\pi\sqrt{LC}$.

Mais cette situation idéale n'existe pas dans la réalité. Toute bobine réelle possède une \cle{résistance interne} $r$ non nulle, et l'on ajoute souvent volontairement un résistor de résistance $R'$. Si l'on observe à l'oscilloscope la tension aux bornes du condensateur d'un tel circuit, on ne voit jamais une sinusoïde éternelle : on voit des oscillations dont l'amplitude \emph{décroît} progressivement jusqu'à s'annuler. Où passe l'énergie ? Comment décrire mathématiquement cet amortissement ? C'est tout l'objet de cette section.

\courssub{Le circuit RLC série : présentation et équation différentielle}

\textbf{Description du circuit}\par

On considère un circuit série comportant un condensateur de capacité $C$, une bobine d'inductance $L$ et de résistance interne $r$, et un résistor de résistance $R'$. On note $R = R' + r$ la \cle{résistance totale} du circuit. Le condensateur a été préalablement chargé sous une tension $E$ (par une source de tension puis déconnexion, ou par un interrupteur à deux positions). À l'instant $t=0$, on connecte le condensateur chargé à l'ensemble bobine--résistor : c'est la \emph{décharge du condensateur} dans le dipôle $RL$ qui produit les oscillations.

\begin{popfigure}
\begin{center}
\begin{circuitikz}[european, cute inductors]
\draw (0,0) to[C, l_={$C$}, v^=$u_C$, *-*] (4,0)
      to[L, l_={$L$}] (4,-2)
      to[R, l_={$R$}] (0,-2)
      to[nos, l_={$K$}] (0,0);
\node at (2,0.5) {\footnotesize $t<0$ : $K$ ouvert, $u_C=E$};
\node at (2,-2.5) {\footnotesize $t\ge 0$ : $K$ fermé, décharge};
\end{circuitikz}
\end{center}
\legende{Circuit RLC série en décharge (position 2). Le condensateur, préalablement chargé à la tension $E$ ($u_C(0)=E$), se décharge dans la bobine et le résistor via l'interrupteur $K$. La résistance $R$ regroupe le résistor $R'$ et la résistance interne $r$ de la bobine.}
\end{popfigure}

\textbf{Établissement de l'équation différentielle}\par

Plaçons-nous en convention récepteur pour les trois dipôles traversés par le courant $i = \dfrac{dq}{dt}$ (le courant entre par la borne $+$ de $u_C$). La loi des mailles (additivité des tensions) s'écrit :
\[ u_C + u_R + u_L = 0 \]
avec :
\begin{itemize}
\item $u_C = \dfrac{q}{C}$ (tension aux bornes du condensateur) ;
\item $u_R = R\,i = R\,\dfrac{dq}{dt}$ (loi d'Ohm pour le résistor et la résistance interne) ;
\item $u_L = L\,\dfrac{di}{dt} = L\,\dfrac{d^2q}{dt^2}$ (tension aux bornes de la bobine idéale).
\end{itemize}
En substituant :
\[ L\,\ddot{q} + R\,\dot{q} + \frac{q}{C} = 0 \]

\begin{propriete}[Équation différentielle des oscillations amorties]
La charge $q(t)$ du condensateur dans un circuit RLC série vérifie l'équation différentielle :
\[ \ddot{q} + \frac{R}{L}\,\dot{q} + \frac{1}{LC}\,q = 0 \]
C'est une équation différentielle linéaire du second ordre à coefficients constants, \emph{avec} terme d'amortissement $\dfrac{R}{L}\dot{q}$. Le cas $R = 0$ redonne l'équation du circuit LC idéal : $\ddot{q} + \dfrac{q}{LC} = 0$.
\end{propriete}

Vérifions l'homogénéité : $[R/L] = \dfrac{\Omega}{\text{H}} = \dfrac{\text{V}\cdot\text{A}^{-1}}{\text{V}\cdot\text{s}\cdot\text{A}^{-1}} = \text{s}^{-1}$, donc $\dfrac{R}{L}\dot{q}$ s'exprime en $\text{C}\cdot\text{s}^{-2}$, comme $\ddot{q}$ et $\dfrac{q}{LC}$. L'équation est homogène.

\courssub{Étude énergétique : pourquoi l'amplitude décroît}

L'intuition est simple : le résistor convertit l'énergie électrique en chaleur par \cle{effet Joule}. À chaque échange d'énergie entre le condensateur et la bobine, une partie est irréversiblement perdue. L'énergie électromagnétique ne peut donc que diminuer. Démontrons-le proprement.

\begin{propriete}[Décroissance de l'énergie électromagnétique — démonstration exigible]
Dans un circuit RLC série, l'énergie électromagnétique totale $E = \dfrac{q^2}{2C} + \dfrac{1}{2}Li^2$ vérifie :
\[ \frac{dE}{dt} = -R\,i^2 \leq 0 \]
L'énergie électromagnétique \textbf{ne peut que décroître} : elle est dissipée par effet Joule dans la résistance totale $R$.
\end{propriete}

\begin{experience}[Démonstration guidée]
On part de l'expression de l'énergie totale et on dérive par rapport au temps :
\begin{align*}
E &= \frac{q^2}{2C} + \frac{1}{2}Li^2 \\
\frac{dE}{dt} &= \frac{q}{C}\,\frac{dq}{dt} + Li\,\frac{di}{dt} = \frac{q}{C}\,i + Li\,\frac{di}{dt} \quad (\text{car } i = \dot{q})\\
\frac{dE}{dt} &= i\left(\frac{q}{C} + L\,\frac{di}{dt}\right)
\end{align*}
Or la loi des mailles donne $\dfrac{q}{C} + Ri + L\dfrac{di}{dt} = 0$, c'est-à-dire $\dfrac{q}{C} + L\dfrac{di}{dt} = -Ri$. En reportant :
\[ \frac{dE}{dt} = i \times (-Ri) = -R\,i^2 \]
Comme $R > 0$ et $i^2 \geq 0$, on a bien $\dfrac{dE}{dt} \leq 0$ : l'énergie décroît. Le terme $-Ri^2$ est exactement l'opposé de la puissance dissipée par effet Joule $P_J = Ri^2$ : l'énergie perdue par le circuit est intégralement convertie en chaleur.
\end{experience}

\begin{attention}[Ce que dit — et ne dit pas — ce bilan]
La décroissance de $E$ n'est pas uniforme : $\dfrac{dE}{dt} = -Ri^2$ est nul aux instants où $i = 0$ (charge extrémale du condensateur) et maximal en valeur absolue quand $|i|$ est maximal. L'énergie décroît donc \og par paliers adoucis \fg{}. En revanche, elle ne peut \emph{jamais} augmenter : aucun apport extérieur ne compense les pertes. C'est la signature d'un système \cle{non conservatif}.
\end{attention}

\courssub{Les trois régimes d'évolution}

L'équation différentielle $\ddot{q} + \dfrac{R}{L}\dot{q} + \dfrac{q}{LC} = 0$ admet des solutions dont la nature dépend de la valeur de $R$ par rapport à une valeur particulière appelée \cle{résistance critique}. On cherche les solutions de la forme $q(t) = e^{pt}$, ce qui mène à l'équation caractéristique $Lp^2 + Rp + 1/C = 0$. Le discriminant $\Delta = R^2 - \dfrac{4L}{C}$ dicte le régime.

\begin{propriete}[Résistance critique et nature des régimes — admis]
On définit la résistance critique du circuit :
\[ R_C = 2\sqrt{\frac{L}{C}} \]
La nature du régime dépend de la comparaison de $R$ à $R_C$ :
\begin{itemize}
\item $R < R_C$ : régime \cle{pseudo-périodique} — oscillations dont l'amplitude décroît exponentiellement ;
\item $R = R_C$ : régime \cle{critique} — retour vers l'équilibre le plus rapide possible, sans oscillation ;
\item $R > R_C$ : régime \cle{apériodique} — retour lent vers l'équilibre, sans oscillation.
\end{itemize}
Vérification dimensionnelle : $[L/C] = \dfrac{\text{H}}{\text{F}} = \dfrac{\text{V}\cdot\text{s}\cdot\text{A}^{-1}}{\text{A}\cdot\text{s}\cdot\text{V}^{-1}} = \Omega^2$, donc $\sqrt{L/C}$ est bien homogène à une résistance.
\end{propriete}

\textbf{Régime pseudo-périodique ($R < R_C$)}\par

C'est le régime observé dans la plupart des circuits de laboratoire. Les racines de l'équation caractéristique sont complexes conjuguées. La charge s'écrit sous la forme :
\[ q(t) = Q_m\,e^{-\frac{R}{2L}t}\cos(\omega t + \varphi) \quad \text{avec} \quad \omega = \sqrt{\frac{1}{LC} - \frac{R^2}{4L^2}} \]
Le facteur exponentiel $e^{-\frac{R}{2L}t}$ joue le rôle d'\cle{enveloppe} : l'amplitude des oscillations décroît exponentiellement. Les oscillations ne sont pas rigoureusement périodiques, mais les passages successifs par zéro sont régulièrement espacés : on définit la \cle{pseudo-période} $T = \dfrac{2\pi}{\omega}$, durée séparant deux passages consécutifs par zéro dans le même sens. On admet que si l'amortissement est faible ($R \ll R_C$), alors $\omega \approx \omega_0 = \dfrac{1}{\sqrt{LC}}$ et $T \approx T_0 = 2\pi\sqrt{LC}$.

\textbf{Régime critique ($R = R_C$)}\par

La racine de l'équation caractéristique est double : $p = -\dfrac{R}{2L} = -\dfrac{1}{\sqrt{LC}} = -\omega_0$. La charge décroît vers zéro sans osciller, le plus rapidement possible :
\[ q(t) = (A + Bt)\,e^{-\omega_0 t} \]
C'est le régime idéal lorsqu'on veut ramener un système à l'équilibre vite et sans rebond.

\textbf{Régime apériodique ($R > R_C$)}\par

Les deux racines $p_1, p_2$ sont réelles, négatives et distinctes. La dissipation est si forte que le condensateur se décharge lentement, sans jamais changer de signe (somme de deux exponentielles décroissantes) :
\[ q(t) = A_1 e^{p_1 t} + A_2 e^{p_2 t} \quad (p_1, p_2 < 0) \]
Retour monotone et lent vers l'équilibre.

\begin{popfigure}
\begin{center}
\begin{tikzpicture}
\begin{axis}[width=\linewidth, height=7cm,
xlabel={$t$ (ms)}, ylabel={$q(t)$ (unité arbitraire)},
xmin=0, xmax=10, ymin=-1.1, ymax=1.2,
grid=major, grid style={popline!50},
legend style={at={(0.98,0.98)}, anchor=north east, font=\footnotesize},
axis lines=left]
\addplot[popcurve, domain=0:10, samples=300, color=popblue] {exp(-0.35*x)*cos(deg(2*pi*0.8*x))};
\addlegendentry{pseudo-périodique ($R < R_C$)}
\addplot[popcurve, domain=0:10, samples=200, dashed, color=popmuted] {exp(-0.35*x)};
\addlegendentry{enveloppe $Q_m e^{-\frac{R}{2L}t}$}
\addplot[popcurve, domain=0:10, samples=200, color=poporange] {(1+2.5*x)*exp(-2.5*x)};
\addlegendentry{critique ($R = R_C$)}
\addplot[popcurve, domain=0:10, samples=200, color=popgreen] {1.2*exp(-0.6*x)-0.2*exp(-4*x)};
\addlegendentry{apériodique ($R > R_C$)}
\end{axis}
\end{tikzpicture}
\end{center}
\legende{Évolution de la charge $q(t)$ pour les trois régimes. En bleu, le régime pseudo-périodique : oscillations amorties dans l'enveloppe exponentielle (pointillés). En orange, le régime critique : retour le plus rapide sans oscillation. En vert, le régime apériodique : retour lent.}
\end{popfigure}

\begin{methode}[Identifier un régime sur un enregistrement]
Face à une courbe $u_C(t)$ ou $q(t)$ obtenue à l'oscilloscope :
\begin{enumerate}
\item Observe-t-on des \textbf{oscillations} (passages répétés par zéro, dépassements) ? $\Rightarrow$ régime \textbf{pseudo-périodique}. Mesure alors la pseudo-période $T$ entre deux maxima consécutifs (ou deux passages par zéro dans le même sens).
\item La courbe revient-elle vers zéro \textbf{sans dépassement} et \textbf{rapidement} ? $\Rightarrow$ régime \textbf{critique}.
\item La courbe revient-elle vers zéro \textbf{sans dépassement} mais \textbf{lentement} ? $\Rightarrow$ régime \textbf{apériodique}.
\end{enumerate}
Pour explorer les régimes au laboratoire, on fait varier $R'$ : en augmentant progressivement la résistance, on passe du pseudo-périodique (oscillations nombreuses), à un amortissement de plus en plus marqué, puis au critique, puis à l'apériodique.
\end{methode}

\begin{attention}[Pseudo-période et période propre : ne pas confondre]
Erreur fréquente au Bac : écrire que les oscillations amorties ont pour période $T_0 = 2\pi\sqrt{LC}$. En toute rigueur :
\begin{itemize}
\item la \textbf{période propre} $T_0 = 2\pi\sqrt{LC}$ ne dépend que de $L$ et $C$ ; c'est la période des oscillations \emph{non amorties} ($R = 0$) ;
\item la \textbf{pseudo-période} $T$ des oscillations amorties vérifie $T > T_0$ et dépend de $R$ via $\omega = \sqrt{\omega_0^2 - (R/2L)^2}$.
\end{itemize}
Cependant, si l'amortissement est faible ($R \ll R_C$), on admet l'approximation $T \approx T_0$, que l'on utilise pour les mesures à l'oscilloscope. Rédige ainsi : \og la pseudo-période mesurée est voisine de la période propre car l'amortissement est faible \fg{}.
\end{attention}

\begin{exemplebox}[Calcul de la résistance critique et discussion des régimes]
Un circuit RLC série comporte une bobine d'inductance $L = \SI{25}{mH}$ et un condensateur de capacité $C = \SI{10}{\micro F}$. La résistance critique vaut :
\[ R_C = 2\sqrt{\frac{L}{C}} = 2\sqrt{\frac{25\times 10^{-3}}{10\times 10^{-6}}} = 2\sqrt{2500} = 2\times 50 = \SI{100}{\ohm} \]
Discutons le régime pour trois valeurs de la résistance totale :
\begin{itemize}
\item $R = \SI{10}{\ohm} < R_C$ : régime \textbf{pseudo-périodique}, amortissement faible ; la pseudo-période est voisine de $T_0 = 2\pi\sqrt{LC} = 2\pi\sqrt{25\times 10^{-3}\times 10^{-5}} \approx \SI{3,1}{ms}$ ;
\item $R = \SI{100}{\ohm} = R_C$ : régime \textbf{critique} ;
\item $R = \SI{500}{\ohm} > R_C$ : régime \textbf{apériodique}.
\end{itemize}
\end{exemplebox}

\begin{exoresolu}[Observation à l'oscilloscope et exploitation]
On réalise le circuit RLC série avec $L = \SI{0,10}{H}$, $C = \SI{1,0}{\micro F}$, et une résistance totale $R$ réglable. Le condensateur est chargé sous $E = \SI{6,0}{V}$ puis déchargé dans le dipôle $RL$. On visualise $u_C(t)$ à l'oscilloscope (base de temps : $\SI{1}{ms/div}$).
\begin{enumerate}
\item Avec $R = \SI{40}{\ohm}$, on observe des oscillations amorties ; deux maxima consécutifs sont séparés de $2\ \text{div}$. Déterminer la pseudo-période $T$ et la comparer à $T_0$.
\item Calculer la résistance critique $R_C$ et préciser le régime observé.
\item Que devient la courbe si l'on règle $R = \SI{1,5}{k\ohm}$ ?
\end{enumerate}
\tcblower
\textbf{1.} La pseudo-période est la durée entre deux maxima consécutifs :
\[ T = 2\ \text{div} \times \SI{1}{ms/div} = \SI{2,0}{ms} \]
La période propre vaut :
\[ T_0 = 2\pi\sqrt{LC} = 2\pi\sqrt{0{,}10\times 1{,}0\times 10^{-6}} = 2\pi\sqrt{10^{-7}} \approx \SI{2,0}{ms} \]
On constate que $T \approx T_0$ : l'amortissement est faible, l'approximation est légitime.

\textbf{2.} Résistance critique :
\[ R_C = 2\sqrt{\frac{L}{C}} = 2\sqrt{\frac{0{,}10}{1{,}0\times 10^{-6}}} = 2\sqrt{10^{5}} \approx \SI{632}{\ohm} \]
Comme $R = \SI{40}{\ohm} \ll R_C \approx \SI{632}{\ohm}$, le régime est \textbf{pseudo-périodique faiblement amorti}, ce qui confirme l'observation et justifie $T \approx T_0$.

\textbf{3.} Avec $R = \SI{1,5}{k\ohm} = \SI{1500}{\ohm} > R_C$, le régime devient \textbf{apériodique} : la tension $u_C$ décroît lentement vers zéro sans osciller ni changer de signe.
\end{exoresolu}

\begin{savaistu}[Le régime critique dans la vie quotidienne]
Lorsqu'on veut qu'un système revienne à l'équilibre vite et sans rebond, on le règle au voisinage du régime critique. C'est le principe des \textbf{fermetures douces} des portes de placard et des tiroirs de cuisine : un amortisseur calibré empêche le claquement. Même chose pour les \textbf{suspensions} des motos et des voitures — très sollicitées sur les pistes déformées entre Douala et Yaoundé : un amortisseur bien réglé ramène la roue au contact du sol en une seule fois, sans oscillation, ce qui garantit l'adhérence et le confort. En électronique, les circuits d'accord des postes radio (réception des stations comme CRTV ou les radios FM privées) exploitent au contraire le régime pseudo-périodique très faiblement amorti pour sélectionner une fréquence précise.
\end{savaistu}

\begin{aretenir}[L'essentiel de la section]
\begin{itemize}
\item La résistance totale $R$ (résistor + résistance interne de la bobine) dissipe l'énergie par effet Joule : les oscillations sont \textbf{amorties}.
\item Équation différentielle : $\ddot{q} + \dfrac{R}{L}\dot{q} + \dfrac{q}{LC} = 0$ (loi des mailles).
\item Bilan énergétique (démonstration exigible) : $\dfrac{dE}{dt} = -Ri^2 \leq 0$ : l'énergie électromagnétique ne peut que décroître.
\item Trois régimes selon la position de $R$ par rapport à $R_C = 2\sqrt{L/C}$ : pseudo-périodique ($R < R_C$), critique ($R = R_C$), apériodique ($R > R_C$).
\item En régime pseudo-périodique faiblement amorti : $T \approx T_0 = 2\pi\sqrt{LC}$, avec en toute rigueur $T > T_0$.
\end{itemize}
\end{aretenir}

\coursec{La bobine : inductance, tension et énergie magnétique}

\courssub{Observation : l'inertie électrique}

Lorsque l'on ouvre brutalement un circuit alimentant une bobine (relais de démarrage d'un moteur de pompe à eau, transformateur de poste de soudage), on observe souvent un arc électrique intense aux bornes de l'interrupteur. Ce phénomène trahit une propriété fondamentale : \emph{la bobine s'oppose à toute variation brutale du courant qui la traverse}. Tout comme une masse en mouvement possède une inertie mécanique qui s'oppose à la variation de sa vitesse, la bobine possède une \cle{inertie électrique} liée à son champ magnétique.

\begin{experience}[Caractéristique tension-courant d'une bobine réelle]
\textbf{Matériel :} Bobine d'inductance $L$ (ex. $L = 50~\mathrm{mH}$) et résistance interne $r$ (ex. $r = 10~\Omega$), générateur de tension variable $u(t)$, multimètres (voltmètre, ampèremètre) ou oscilloscope.
\textbf{Protocole :} On impose une tension $u(t)$ variable (rampe, créneau) et on mesure $i(t)$.
\textbf{Observations :}
\begin{itemize}
\item Si $u$ est constante, $i$ est constante : la bobine se comporte comme une simple résistance $r$ ($u = ri$).
\item Si $u$ varie, la relation $u = ri$ ne suffit plus : il faut ajouter un terme proportionnel à $\dfrac{di}{dt}$.
\end{itemize}
\textbf{Interprétation :} La bobine stocke de l'énergie magnétique. La variation de ce stock d'énergie nécessite une puissance supplémentaire, d'où le terme $L\,di/dt$.
\end{experience}

\begin{definition}[Inductance propre $L$]
L'\cle{inductance propre} $L$ (unité : le \cle{henry} \SI{}{H} = \SI{}{V\,s\,A^{-1}}) d'un dipôle bobine caractérise sa capacité à stocker de l'énergie magnétique et à s'opposer aux variations de courant. Pour une bobine réelle de résistance interne $r$, la tension $u$ à ses bornes (recepteur) et le courant $i$ qui la traverse (entrant par la borne $+$) sont liés par :
\[
u = r\,i + L\,\frac{di}{dt}.
\]
Si la bobine est idéale ($r = 0$), $u = L\,\dfrac{di}{dt}$.
\end{definition}

\begin{propriete}[Énergie magnétique stockée]
La puissance reçue par la bobine est $P = u\,i = r i^2 + L\,i\,\dfrac{di}{dt}$. Le terme $r i^2$ est l'énergie dissipée par effet Joule (chaleur). Le terme $L\,i\,\dfrac{di}{dt} = \dfrac{d}{dt}\left(\dfrac{1}{2}L i^2\right)$ correspond à la puissance stockée sous forme d'énergie magnétique.
\[
E_m = \frac{1}{2}\,L\,i^2 \qquad (\text{en joules, } \SI{}{J}).
\]
\begin{attention}[Sens du courant et signe de l'énergie]
L'énergie $E_m = \frac{1}{2}Li^2$ est toujours positive (proportionnelle à $i^2$). Elle ne dépend que de l'intensité du courant, pas de son sens. En revanche, la tension d'induction $L\,di/dt$ change de signe selon que le courant croît ($di/dt > 0$, la bobine \emph{absorbe} de l'énergie, $u$ et $i$ même signe) ou décroît ($di/dt < 0$, la bobine \emph{restitue} de l'énergie, $u$ et $i$ signes opposés).
\end{attention}
\end{propriete}

\begin{exemplebox}[Bobine de relais de démarrage]
Un relais de pompe agricole possède une bobine d'inductance $L = 0,50~\mathrm{H}$ et de résistance $r = 20~\Omega$. Au démarrage, le courant passe de $0$ à $I_{\max} = 0,50~\mathrm{A}$ en $\Delta t = 10~\mathrm{ms}$.
\begin{itemize}
\item Énergie magnétique stockée in fine : $E_m = \frac{1}{2}\times 0,50 \times (0,50)^2 = 0,0625~\mathrm{J}$.
\item Tension d'induction moyenne pendant la montée : $\langle u_L \rangle = L\,\dfrac{\Delta i}{\Delta t} = 0,50 \times \dfrac{0,50}{0,010} = 25~\mathrm{V}$.
\item Tension totale aux bornes : $u = r I_{\max} + \langle u_L \rangle = 10 + 25 = 35~\mathrm{V}$.
\end{itemize}
C'est cette surtension transitoire qui use les contacts des interrupteurs ; on la limite souvent par une diode de roue libre (parallèle à la bobine) ou un varistance.
\end{exemplebox}

\coursec{Le circuit LC idéal : équation différentielle des oscillations libres}

\courssub{Mise en situation : décharge d'un condensateur dans une bobine}

On charge un condensateur de capacité $C$ à la tension $U_0$ (charge $Q_0 = C U_0$). On le déconnecte de la source et on le branche sur une bobine \emph{idéale} d'inductance $L$ (résistance nulle, $r=0$, pas de résistance externe). Le circuit est fermé à l'instant $t=0$ : $q(0)=Q_0$, $i(0)=0$.

\begin{center}
\begin{popfigure}
\begin{tikzpicture}[european, font=\small, scale=1.2]
\draw (0,0) to[C=$C$, v^={$u_C$}, l_={$q$}] (0,-2)
      to[L=$L$, v^={$u_L$}, i={$i$}] (0,0);
\node[align=center] at (1.5,-1) {Circuit LC série\\$u_L + u_C = 0$};
\end{tikzpicture}
\legende{Circuit LC série idéal : le condensateur chargé se décharge dans la bobine.}
\end{popfigure}
\end{center}

\courssub{Établissement de l'équation différentielle}

La loi des mailles (tension algébrique nulle dans la maille unique) s'écrit :
\[
u_L + u_C = 0.
\]
Avec les relations constitutives : $u_L = L\,\dfrac{di}{dt}$ (bobine idéale) et $u_C = \dfrac{q}{C}$ (condensateur, convention récepteur : $i = \dfrac{dq}{dt}$ entrant par la borne $+$ du condensateur).
\[
L\,\frac{di}{dt} + \frac{q}{C} = 0.
\]
Puisque $i = \dfrac{dq}{dt}$, on obtient une équation différentielle du second ordre en la charge $q(t)$ :
\[
L\,\frac{d^2q}{dt^2} + \frac{q}{C} = 0
\quad\Longleftrightarrow\quad
\ddot{q} + \frac{1}{LC}\,q = 0.
\]

\begin{propriete}[Équation différentielle du circuit LC idéal]
La charge $q(t)$ du condensateur dans un circuit LC série idéal (sans résistance) obéit à l'équation différentielle linéaire d'ordre 2 à coefficients constants :
\[
\ddot{q} + \omega_0^2\,q = 0 \qquad\text{avec}\qquad \omega_0 = \frac{1}{\sqrt{LC}}.
\]
$\omega_0$ est la \cle{pulsation propre} du circuit (en $\mathrm{rad\,s^{-1}}$).
\end{propriete}

\begin{aretenir}[Analyse dimensionnelle de $\omega_0$]
$[L] = \mathrm{H} = \mathrm{V\,s\,A^{-1}} = \mathrm{kg\,m^2\,s^{-2}\,A^{-2}}$.
$[C] = \mathrm{F} = \mathrm{C\,V^{-1}} = \mathrm{A^2\,s^4\,kg^{-1}\,m^{-2}}$.
$[LC] = \mathrm{s^2}$, donc $[\omega_0] = [1/\sqrt{LC}] = \mathrm{s^{-1}}$ (rad/s). \checkmark
\end{aretenir}

\courssub{Solution générale : oscillations sinusoïdales}

L'équation $\ddot{q} + \omega_0^2 q = 0$ est celle de l'\cle{oscillateur harmonique}. Sa solution générale est :
\[
q(t) = Q_m \cos(\omega_0 t + \varphi),
\]
où $Q_m$ (amplitude de la charge, en $\mathrm{C}$) et $\varphi$ (phase à l'origine, en $\mathrm{rad}$) sont déterminés par les conditions initiales.
L'intensité instantanée est :
\[
i(t) = \frac{dq}{dt} = -Q_m \omega_0 \sin(\omega_0 t + \varphi) = I_m \cos(\omega_0 t + \varphi + \tfrac{\pi}{2}),
\]
avec $I_m = Q_m \omega_0 = \dfrac{Q_m}{\sqrt{LC}}$ (amplitude du courant, en $\mathrm{A}$).

\begin{exemplebox}[Conditions initiales classiques : décharge du condensateur]
À $t=0$ : $q(0) = Q_0$ (condensateur chargé), $i(0) = 0$ (courant nul au départ).
\begin{align*}
q(0) &= Q_m \cos\varphi = Q_0, \\
i(0) &= -Q_m \omega_0 \sin\varphi = 0 \;\Rightarrow\; \sin\varphi = 0 \;\Rightarrow\; \varphi = 0 \text{ (choisi)}.
\end{align*}
Donc $Q_m = Q_0$ et finalement :
\[
\boxed{q(t) = Q_0 \cos(\omega_0 t)}, \qquad
\boxed{i(t) = -Q_0 \omega_0 \sin(\omega_0 t) = -\frac{Q_0}{\sqrt{LC}} \sin(\omega_0 t)}.
\]
Le courant est en \emph{avance de phase de $\pi/2$} (un quart de période) sur la charge : quand $q$ est maximum, $i=0$ ; quand $q=0$, $|i|$ est maximum.
\end{exemplebox}

\begin{popfigure}
\begin{center}
\begin{tikzpicture}[font=\small]
\begin{axis}[
    width=\linewidth, height=6cm,
    axis lines=middle, xlabel=$t$, ylabel={$q, i$},
    xtick={0,1.57,3.14,4.71,6.28}, xticklabels={$0$, $\frac{T_0}{4}$, $\frac{T_0}{2}$, $\frac{3T_0}{4}$, $T_0$},
    ytick={-1,1}, yticklabels={$-1$, $1$}, ymin=-1.2, ymax=1.2,
    domain=0:6.28, samples=200,
    legend style={at={(0.5,1.1)}, anchor=north, legend columns=2}
]
\addplot[popblue, thick] {cos(deg(x))}; \addlegendentry{$q(t)/Q_0$}
\addplot[poporange, thick] {-sin(deg(x))}; \addlegendentry{$i(t)/I_m$}
\end{axis}
\end{tikzpicture}
\end{center}
\legende{Oscillations de la charge $q(t)$ (bleu) et de l'intensité $i(t)$ (orange) dans un circuit LC idéal. $q$ et $i$ sont déphasés de $\pi/2$.}
\end{popfigure}

\coursec{Période propre, fréquence et applications}

\courssub{Formule de Thomson}

La période propre $T_0$ et la fréquence propre $f_0$ des oscillations libres sont déduites de $\omega_0$ :
\[
T_0 = \frac{2\pi}{\omega_0} = 2\pi\sqrt{LC}, \qquad
f_0 = \frac{1}{T_0} = \frac{1}{2\pi\sqrt{LC}}.
\]
C'est la célèbre \cle{formule de Thomson} (1853), base de la radiocommunication.

\begin{methode}[Calcul de $L$, $C$, $f_0$ ou $T_0$]
\begin{enumerate}
\item Identifier les grandeurs connues et l'inconnue.
\item Écrire la formule adaptée : $f_0 = \dfrac{1}{2\pi\sqrt{LC}}$ ou $T_0 = 2\pi\sqrt{LC}$.
\item Vérifier les unités : $L$ en $\mathrm{H}$, $C$ en $\mathrm{F}$, $f_0$ en $\mathrm{Hz}$, $T_0$ en $\mathrm{s}$.
\item Calculer et exprimer le résultat avec l'unité SI et 2 à 3 chiffres significatifs.
\end{enumerate}
\end{methode}

\begin{exoresolu}[Accord d'un circuit LC sur une station FM]
On souhaite recevoir la station \textbf{Radio Cameroon} à $f_0 = 94,5~\mathrm{MHz}$ avec un condensateur variable $C_{\max} = 100~\mathrm{pF}$. Quelle inductance $L$ faut-il choisir pour la bobine d'accord ?
\tcblower
On utilise $f_0 = \dfrac{1}{2\pi\sqrt{LC}} \;\Rightarrow\; L = \dfrac{1}{(2\pi f_0)^2 C}$.
Avec $f_0 = 94,5\times 10^6~\mathrm{Hz}$ et $C = 100\times 10^{-12}~\mathrm{F}$ :
\[
L = \frac{1}{(2\pi \times 94,5\times 10^6)^2 \times 100\times 10^{-12}}
  \approx \frac{1}{(5,94\times 10^8)^2 \times 10^{-10}}
  \approx \frac{1}{3,53\times 10^7}
  \approx 2,83\times 10^{-8}~\mathrm{H} = 28,3~\mathrm{nH}.
\]
Une bobine de quelques spires (diamètre $\approx 5~\mathrm{mm}$, longueur $\approx 10~\mathrm{mm}$) convient. En pratique, on utilise un condensateur variable pour balayer la bande FM ($87,5$--$108~\mathrm{MHz}$).
\end{exoresolu}

\begin{savaistu}[Le quartz : un LC mécanique de très haute précision]
Dans les montres, les téléphones (MTN, Orange, CAMTEL) et les ordinateurs, l'oscillateur de référence n'est pas un circuit LC classique (trop instable en température), mais un \emph{quartz piézoélectrique}. Mécaniquement, une lame de quartz taillée vibre sur son mode propre comme un système masse-ressort : sa fréquence propre $f_0 = \frac{1}{2\pi}\sqrt{k/m}$ est extrêmement stable ($Q > 10^5$). Électriquement, il se modélise par un circuit RLC série de très haut facteur de qualité, en parallèle avec une capacité parasite. C'est l'application industrielle majeure de l'analogie électromécanique !
\end{savaistu}

\coursec{Étude énergétique du circuit LC : conservation de l'énergie électromagnétique}

\courssub{Énergie électromagnétique totale}

Dans le circuit LC idéal, il n'y a aucun élément dissipatif ($R=0$). L'énergie totale échangée entre le condensateur et la bobine doit se conserver.

\begin{itemize}
\item \textbf{Énergie électrique du condensateur :} $E_e(t) = \dfrac{q^2(t)}{2C} = \dfrac{Q_m^2}{2C}\cos^2(\omega_0 t + \varphi)$.
\item \textbf{Énergie magnétique de la bobine :} $E_m(t) = \dfrac{1}{2}L\,i^2(t) = \dfrac{1}{2}L\,Q_m^2\omega_0^2\sin^2(\omega_0 t + \varphi)$.
\end{itemize}
Rappelons que $\omega_0^2 = 1/(LC)$, donc $L\omega_0^2 = 1/C$. L'énergie magnétique s'écrit aussi :
\[
E_m(t) = \frac{Q_m^2}{2C}\sin^2(\omega_0 t + \varphi).
\]

\begin{propriete}[Conservation de l'énergie électromagnétique dans un circuit LC idéal]
L'énergie électromagnétique totale $E(t) = E_e(t) + E_m(t)$ est constante :
\[
E(t) = \frac{Q_m^2}{2C}\bigl[\cos^2(\omega_0 t + \varphi) + \sin^2(\omega_0 t + \varphi)\bigr] = \frac{Q_m^2}{2C} = \frac{1}{2}L I_m^2 = \text{constante}.
\]
\begin{attention}[Énergie totale vs énergies partielles]
$E_e$ et $E_m$ oscillent en \emph{opposition de phase} : quand l'une est maximale, l'autre est nulle. Leur \emph{somme} est constante. C'est le transfert perpétuel d'énergie entre le champ électrique (condensateur) et le champ magnétique (bobine) qui entretient l'oscillation.
\end{attention}
\end{propriete}

\begin{popfigure}
\begin{center}
\begin{tikzpicture}[font=\small]
\begin{axis}[
    width=\linewidth, height=6cm,
    axis lines=middle, xlabel=$t$, ylabel={Énergie},
    xtick={0,1.57,3.14,4.71,6.28}, xticklabels={$0$, $\frac{T_0}{4}$, $\frac{T_0}{2}$, $\frac{3T_0}{4}$, $T_0$},
    ytick={0,0.5,1}, yticklabels={$0$, $\frac{E_{\max}}{2}$, $E_{\max}$}, ymin=-0.1, ymax=1.1,
    domain=0:6.28, samples=200,
    legend style={at={(0.5,1.1)}, anchor=north, legend columns=3}
]
\addplot[popblue, thick] {cos(deg(x))^2}; \addlegendentry{$E_e(t)$}
\addplot[poporange, thick] {sin(deg(x))^2}; \addlegendentry{$E_m(t)$}
\addplot[popgreen, thick, dashed] {1}; \addlegendentry{$E_{\text{tot}}$}
\end{axis}
\end{tikzpicture}
\end{center}
\legende{Échanges d'énergie dans le circuit LC : $E_e$ (bleu) et $E_m$ (orange) oscillent en opposition de phase, leur somme (vert pointillé) est constante.}
\end{popfigure}

\courssub{Démonstration par la dérivée (méthode générale)}

On peut prouver la conservation sans connaître la solution explicite, en dérivant l'énergie totale :
\[
E = \frac{1}{2}L i^2 + \frac{q^2}{2C}, \quad i = \dot{q}.
\]
\[
\frac{dE}{dt} = L i \dot{i} + \frac{q}{C}\dot{q} = i\left(L \ddot{q} + \frac{q}{C}\right).
\]
Or l'équation différentielle donne $L \ddot{q} + q/C = 0$. Donc $\dfrac{dE}{dt} = 0 \;\Rightarrow\; E = \text{cte}$.
Cette méthode par bilan de puissance est généralisable au circuit RLC (où $dE/dt = -Ri^2$).

\begin{exoresolu}[Amplitudes de charge et de courant]
Un circuit LC est constitué de $L = 40~\mathrm{mH}$ et $C = 10~\mathrm{nF}$. L'énergie totale est $E = 2,0~\mathrm{mJ}$. Déterminer $Q_m$ et $I_m$.
\tcblower
$E = \dfrac{Q_m^2}{2C} \;\Rightarrow\; Q_m = \sqrt{2CE} = \sqrt{2 \times 10\times 10^{-9} \times 2,0\times 10^{-3}} = \sqrt{4,0\times 10^{-11}} = 6,32\times 10^{-6}~\mathrm{C} = 6,32~\mu\mathrm{C}$.
$E = \dfrac{1}{2}L I_m^2 \;\Rightarrow\; I_m = \sqrt{\dfrac{2E}{L}} = \sqrt{\dfrac{4,0\times 10^{-3}}{40\times 10^{-3}}} = \sqrt{0,10} = 0,316~\mathrm{A}$.
Vérification : $\omega_0 = 1/\sqrt{LC} = 1/\sqrt{40\times 10^{-3} \times 10\times 10^{-9}} = 1/\sqrt{4\times 10^{-10}} = 5,0\times 10^4~\mathrm{rad/s}$.
$I_m = Q_m \omega_0 = 6,32\times 10^{-6} \times 5,0\times 10^4 = 0,316~\mathrm{A}$. \checkmark
\end{exoresolu}

\coursec{Le circuit RLC série : oscillations amorties}

\courssub{Équation différentielle et régimes}

Tout circuit réel possède une résistance équivalente série $R$ (résistance interne de la bobine $r$ + résistance externe éventuelle). La loi des mailles devient :
\[
u_L + u_R + u_C = 0 \;\Rightarrow\; L\,\frac{di}{dt} + R\,i + \frac{q}{C} = 0.
\]
Avec $i = \dot{q}$, on obtient l'équation différentielle du \cle{circuit RLC série} :
\[
L\,\ddot{q} + R\,\dot{q} + \frac{q}{C} = 0
\quad\Longleftrightarrow\quad
\ddot{q} + \frac{R}{L}\,\dot{q} + \frac{1}{LC}\,q = 0.
\]

\begin{propriete}[Équation caractéristique et régimes]
L'équation caractéristique associée est $L\lambda^2 + R\lambda + 1/C = 0$, soit $\lambda^2 + \dfrac{R}{L}\lambda + \omega_0^2 = 0$ avec $\omega_0 = 1/\sqrt{LC}$.
Son discriminant est $\Delta = \left(\dfrac{R}{L}\right)^2 - 4\omega_0^2 = \dfrac{R^2}{L^2} - \dfrac{4}{LC} = \dfrac{1}{L^2}\left(R^2 - \dfrac{4L}{C}\right)$.
On définit la \cle{résistance critique} $R_c = 2\sqrt{\dfrac{L}{C}}$.
Trois régimes apparaissent selon le signe de $\Delta$ (ou la comparaison de $R$ à $R_c$) :
\end{propriete}

\begin{center}
\rowcolors{1}{popcream}{white}
\begin{tabularx}{\linewidth}{|l|X|X|X|}
\hline
\rowcolor{popblueL} \textbf{Régime} & \textbf{Condition} & \textbf{Racines $\lambda$} & \textbf{Allure de $q(t)$} \\
\hline
\textbf{Pseudo-périodique} & $R < R_c$ \quad ($\Delta < 0$) & $\lambda = -\alpha \pm j\omega$ & $q(t) = Q_m e^{-\alpha t}\cos(\omega t + \varphi)$ \\
\textbf{Critique} & $R = R_c$ \quad ($\Delta = 0$) & $\lambda = -\alpha$ (double) & $q(t) = (A + Bt)e^{-\alpha t}$ \\
\textbf{Apériodique} & $R > R_c$ \quad ($\Delta > 0$) & $\lambda_1, \lambda_2 < 0$ réels & $q(t) = A e^{\lambda_1 t} + B e^{\lambda_2 t}$ \\
\hline
\end{tabularx}
\end{center}
où $\alpha = \dfrac{R}{2L}$ (coefficient d'amortissement, en $\mathrm{s^{-1}}$) et $\omega = \sqrt{\omega_0^2 - \alpha^2}$ (pulsation amortie, en $\mathrm{rad/s}$).

\begin{attention}[Vocabulaire précis pour le Bac]
\begin{itemize}
\item On dit \emph{pseudo-périodique} (et non « périodique ») car l'amplitude décroît : ce n'est pas une fonction périodique au sens mathématique strict.
\item La \emph{période pseudo-périodique} est $T = \dfrac{2\pi}{\omega} > T_0$. L'amortissement \emph{allonge} la période.
\item Le \emph{régime critique} correspond au retour à l'équilibre \emph{le plus rapide sans oscillation} (porte de bâtiment, amortisseur de voiture bien réglé).
\item En régime apériodique, les deux exponentielles sont décroissantes ; la plus lente (celle avec $\lambda$ le plus proche de 0) domine à long terme.
\end{itemize}
\end{attention}

\courssub{Bilan énergétique : dissipation par effet Joule}

L'énergie électromagnétique totale $E = \frac{1}{2}Li^2 + \frac{q^2}{2C}$ n'est plus conservée.
\[
\frac{dE}{dt} = L i \dot{i} + \frac{q}{C}\dot{q} = i\left(L \ddot{q} + \frac{q}{C}\right).
\]
En utilisant l'équation différentielle $L\ddot{q} + R\dot{q} + q/C = 0 \;\Rightarrow\; L\ddot{q} + q/C = -R\dot{q} = -Ri$ :
\[
\frac{dE}{dt} = i(-Ri) = -R i^2 \leq 0.
\]
La puissance dissipée par effet Joule dans la résistance $R$ est exactement $P_J = R i^2$. L'énergie électromagnétique décroît au rythme de cette dissipation.

\begin{exoresolu}[Détermination du régime et paramètres]
Un circuit RLC série a $L = 100~\mathrm{mH}$, $C = 10~\mu\mathrm{F}$, $R = 10~\Omega$. Déterminer le régime, la pulsation propre $\omega_0$, le coefficient d'amortissement $\alpha$, et si pseudo-périodique, la pulsation $\omega$ et la période $T$.
\tcblower
$\omega_0 = \dfrac{1}{\sqrt{LC}} = \dfrac{1}{\sqrt{100\times 10^{-3} \times 10\times 10^{-6}}} = \dfrac{1}{\sqrt{10^{-6}}} = 1000~\mathrm{rad/s}$.
$R_c = 2\sqrt{\dfrac{L}{C}} = 2\sqrt{\dfrac{100\times 10^{-3}}{10\times 10^{-6}}} = 2\sqrt{10^4} = 200~\Omega$.
Comme $R = 10~\Omega < R_c = 200~\Omega$ : \textbf{régime pseudo-périodique}.
$\alpha = \dfrac{R}{2L} = \dfrac{10}{2 \times 0,10} = 50~\mathrm{s^{-1}}$.
$\omega = \sqrt{\omega_0^2 - \alpha^2} = \sqrt{10^6 - 2500} \approx \sqrt{997500} \approx 998,7~\mathrm{rad/s}$.
$T = \dfrac{2\pi}{\omega} \approx \dfrac{6,283}{998,7} \approx 6,29~\mathrm{ms}$.
Constante de temps de l'enveloppe : $\tau = 1/\alpha = 20~\mathrm{ms}$. L'amplitude est divisée par $e$ en $20~\mathrm{ms}$ (environ 3 périodes).
\end{exoresolu}

\coursec{Analogies électromécaniques et synthèse}

Dans la section précédente, nous avons établi que la charge du condensateur dans un circuit $RLC$ série obéit à l'équation différentielle
\[
L\,\ddot{q} + R\,\dot{q} + \frac{q}{C} = 0,
\]
et que, selon la valeur de la résistance $R$, le circuit évolue en régime pseudo-périodique, critique ou apériodique. Or tu as rencontré en mécanique une équation de structure \emph{strictement identique} : celle du système masse-ressort amorti. Cette ressemblance n'est pas un hasard : elle révèle une \cle{analogie} profonde entre l'électricité et la mécanique, que nous allons exploiter dans cette dernière section pour consolider, mémoriser et vérifier tous les résultats de la leçon.

\courssub{Deux systèmes, une même équation}

\textbf{Situation d'observation.} Accroche une masse $m$ à un ressort de raideur $k$, écarte-la de sa position d'équilibre et lâche-la : elle oscille, et les oscillations s'amortissent à cause des frottements. Charge un condensateur, branche-le sur une bobine réelle : la charge oscille, et les oscillations s'amortissent à cause de la résistance. Deux mondes en apparence sans rapport — un solide qui coulisse, des électrons qui circulent — produisent des courbes $x(t)$ et $q(t)$ superposables. Pourquoi ?

\textbf{Côté mécanique.} Considérons une masse $m$ mobile sur un axe horizontal, reliée à un ressort de raideur $k$ et soumise à une force de frottement fluide $\vec{F}_f = -f\,\vec{v}$ (où $f$ est le coefficient de frottement, en $\mathrm{kg\,s^{-1}}$). Le bilan des forces donne : le poids $\vec{P}$ et la réaction du support $\vec{R}_N$ se compensent ; restent la tension du ressort $\vec{T} = -kx\,\vec{\imath}$ et la force de frottement. La deuxième loi de Newton projetée sur l'axe s'écrit :
\[
m\,\ddot{x} = -kx - f\,\dot{x}
\quad\Longleftrightarrow\quad
m\,\ddot{x} + f\,\dot{x} + kx = 0.
\]

\textbf{Côté électrique.} Dans le circuit $RLC$ série, la loi des mailles $u_L + u_R + u_C = 0$ avec $u_L = L\,\dfrac{di}{dt}$, $u_R = Ri$ et $i = \dfrac{dq}{dt}$ conduit à :
\[
L\,\ddot{q} + R\,\dot{q} + \frac{q}{C} = 0.
\]

Divisons chaque équation par son premier coefficient pour faire apparaître la pulsation propre :
\[
\ddot{x} + \frac{f}{m}\,\dot{x} + \frac{k}{m}\,x = 0
\qquad\text{et}\qquad
\ddot{q} + \frac{R}{L}\,\dot{q} + \frac{1}{LC}\,q = 0.
\]
Les deux équations ont \emph{exactement} la même forme mathématique. Il suffit donc d'établir un « dictionnaire » de traduction pour passer de l'une à l'autre.

\begin{propriete}[Analogies électromécaniques (complément du programme, admise)]
Le circuit $RLC$ série et le système masse-ressort amorti sont régis par des équations différentielles de même forme. On dit qu'ils sont \cle{analogues}. La correspondance entre grandeurs est :
\[
q \leftrightarrow x, \qquad i = \dot{q} \leftrightarrow v = \dot{x}, \qquad L \leftrightarrow m, \qquad R \leftrightarrow f, \qquad \frac{1}{C} \leftrightarrow k.
\]
Toute propriété démontrée sur l'un des deux systèmes se traduit immédiatement en une propriété de l'autre. Cette analogie est présentée ici à titre de complément culturel et méthodologique ; elle n'est pas exigible en tant que telle au Baccalauréat, mais elle est un outil de vérification redoutable.
\end{propriete}

\begin{attention}[Le piège classique : $C$ ne correspond pas à $k$]
L'erreur la plus fréquente consiste à associer la capacité $C$ à la raideur $k$. C'est faux : c'est $\dfrac{1}{C}$ qui joue le rôle de $k$. Vérification sur les périodes propres :
\[
T_0 = 2\pi\sqrt{\frac{m}{k}} \;\xrightarrow{\;m\to L,\; k\to 1/C\;}\; 2\pi\sqrt{\frac{L}{1/C}} = 2\pi\sqrt{LC}. \quad\checkmark
\]
Si l'on avait posé $C \leftrightarrow k$, on aurait obtenu $2\pi\sqrt{L/C}$, dimensionnellement et numériquement faux. Retiens : un condensateur de \emph{grande} capacité se charge « mollement », comme un ressort \emph{souple} (petit $k$) ; l'inverse de la capacité mesure la « raideur » électrique.
\end{attention}

\begin{popfigure}
\begin{center}
\begin{tikzpicture}[font=\small]
% Colonne mécanique
\node[draw=popblue, fill=popblueL, rounded corners, minimum width=4.6cm, minimum height=0.8cm, font=\bfseries] (mtitle) at (-3.2,3.4) {Mécanique};
\node[draw=poporange, fill=poporangeL, rounded corners, minimum width=4.6cm, minimum height=0.8cm, font=\bfseries] (etitle) at (3.2,3.4) {Électricité};
\node[draw=popline, rounded corners, fill=white, minimum width=4.4cm, minimum height=0.65cm] (m1) at (-3.2,2.3) {élongation $x$ (m)};
\node[draw=popline, rounded corners, fill=white, minimum width=4.4cm, minimum height=0.65cm] (e1) at (3.2,2.3) {charge $q$ (C)};
\node[draw=popline, rounded corners, fill=white, minimum width=4.4cm, minimum height=0.65cm] (m2) at (-3.2,1.3) {vitesse $v=\dot{x}$ (m/s)};
\node[draw=popline, rounded corners, fill=white, minimum width=4.4cm, minimum height=0.65cm] (e2) at (3.2,1.3) {intensité $i=\dot{q}$ (A)};
\node[draw=popline, rounded corners, fill=white, minimum width=4.4cm, minimum height=0.65cm] (m3) at (-3.2,0.3) {masse $m$ (kg)};
\node[draw=popline, rounded corners, fill=white, minimum width=4.4cm, minimum height=0.65cm] (e3) at (3.2,0.3) {inductance $L$ (H)};
\node[draw=popline, rounded corners, fill=white, minimum width=4.4cm, minimum height=0.65cm] (m4) at (-3.2,-0.7) {coefficient $f$ (kg/s)};
\node[draw=popline, rounded corners, fill=white, minimum width=4.4cm, minimum height=0.65cm] (e4) at (3.2,-0.7) {résistance $R$ ($\Omega$)};
\node[draw=popline, rounded corners, fill=white, minimum width=4.4cm, minimum height=0.65cm] (m5) at (-3.2,-1.7) {raideur $k$ (N/m)};
\node[draw=popline, rounded corners, fill=white, minimum width=4.4cm, minimum height=0.65cm] (e5) at (3.2,-1.7) {$1/C$ (F$^{-1}$)};
\draw[-{Stealth}, very thick, poppurple] (m1) -- (e1);
\draw[-{Stealth}, very thick, poppurple] (m2) -- (e2);
\draw[-{Stealth}, very thick, poppurple] (m3) -- (e3);
\draw[-{Stealth}, very thick, poppurple] (m4) -- (e4);
\draw[-{Stealth}, very thick, poppurple] (m5) -- (e5);
\node[font=\itshape\small, popdark] at (0,-2.6) {$m\ddot{x}+f\dot{x}+kx=0 \;\;\longleftrightarrow\;\; L\ddot{q}+R\dot{q}+\dfrac{q}{C}=0$};
\end{tikzpicture}
\end{center}
\legende{Le dictionnaire des analogies électromécaniques : chaque grandeur mécanique possède son équivalent électrique, et les équations différentielles se correspondent terme à terme.}
\end{popfigure}

\courssub{Correspondance énergétique}

L'analogie ne s'arrête pas aux équations du mouvement : elle s'étend aux \cle{bilans énergétiques}, ce qui lui donne toute sa profondeur physique.

\begin{itemize}
\item \textbf{Énergie cinétique et énergie magnétique.} La masse en mouvement stocke l'énergie cinétique $E_c = \dfrac{1}{2}mv^2$ ; la bobine parcourue par un courant stocke l'énergie magnétique $E_m = \dfrac{1}{2}Li^2$. La correspondance $m \leftrightarrow L$, $v \leftrightarrow i$ transforme l'une en l'autre. La masse « s'oppose » aux variations de vitesse (inertie) exactement comme l'inductance s'oppose aux variations de courant (loi de Lenz).
\item \textbf{Énergie potentielle élastique et énergie électrique.} Le ressort déformé stocke $E_p = \dfrac{1}{2}kx^2$ ; le condensateur chargé stocke $E_e = \dfrac{q^2}{2C} = \dfrac{1}{2}\cdot\dfrac{1}{C}\cdot q^2$. La correspondance $k \leftrightarrow 1/C$, $x \leftrightarrow q$ est parfaite.
\item \textbf{Puissance dissipée.} Les frottements dissipent la puissance $P_f = fv^2$ ; la résistance dissipe par effet Joule $P_J = Ri^2$. Là encore, $f \leftrightarrow R$, $v \leftrightarrow i$.
\end{itemize}

Le bilan complet s'écrit de façon jumelle :
\[
\frac{d}{dt}\underbrace{\left(\tfrac{1}{2}mv^2 + \tfrac{1}{2}kx^2\right)}_{\text{énergie mécanique}} = -fv^2 \leq 0
\qquad\text{et}\qquad
\frac{d}{dt}\underbrace{\left(\tfrac{1}{2}Li^2 + \frac{q^2}{2C}\right)}_{\text{énergie électromagnétique}} = -Ri^2 \leq 0.
\]
Dans les deux cas, l'énergie totale ne peut que décroître, exactement au rythme de la dissipation. Et dans les deux cas, si l'on supprime le terme dissipatif ($f = 0$ ou $R = 0$), l'énergie totale se \emph{conserve} et les oscillations deviennent rigoureusement sinusoïdales : c'est le couple {oscillateur harmonique} $\leftrightarrow$ {circuit $LC$ idéal}.

\begin{center}
\rowcolors{1}{popcream}{white}
\begin{tabularx}{\linewidth}{|l|X|X|}
\hline
\rowcolor{popblueL} \textbf{Grandeur} & \textbf{Mécanique (masse-ressort)} & \textbf{Électricité (RLC série)} \\
\hline
Équation & $m\ddot{x}+f\dot{x}+kx=0$ & $L\ddot{q}+R\dot{q}+\dfrac{q}{C}=0$ \\
Pulsation propre & $\omega_0=\sqrt{k/m}$ & $\omega_0=\sqrt{1/(LC)}$ \\
Période propre & $T_0=2\pi\sqrt{m/k}$ & $T_0=2\pi\sqrt{LC}$ \\
Énergie « motrice » & $E_c=\tfrac12 mv^2$ & $E_m=\tfrac12 Li^2$ \\
Énergie « stockée » & $E_p=\tfrac12 kx^2$ & $E_e=\dfrac{q^2}{2C}$ \\
Puissance dissipée & $fv^2$ & $Ri^2$ \\
Régime sans amortissement & $f=0$ : oscillations entretenues & $R=0$ : circuit $LC$ idéal \\
\hline
\end{tabularx}
\end{center}

\courssub{Lecture des régimes : du ressort au circuit}

L'analogie permet de \emph{visualiser} les trois régimes du circuit $RLC$ avec les mains, sur un ressort :

\begin{itemize}
\item \textbf{Régime pseudo-périodique} ($R$ faible, $f$ faible) : le ressort oscille de nombreuses fois avant de s'immobiliser ; la charge $q(t)$ oscille avec une amplitude qui décroît exponentiellement. C'est le régime observé en TP avec une bobine de faible résistance.
\item \textbf{Régime critique} ($R = R_c = 2\sqrt{L/C}$, $f = f_c = 2\sqrt{mk}$) : le ressort revient à l'équilibre \emph{le plus vite possible sans osciller}, comme la porte d'un bâtiment public munie d'un ferme-porte bien réglé ; de même, $q(t)$ retourne à zéro en un temps minimal sans dépassement.
\item \textbf{Régime apériodique} ($R > R_c$, frottements très forts) : plonge le ressort dans du miel — il revient lentement vers l'équilibre sans jamais osciller ; le condensateur se décharge lentement, sans oscillation.
\end{itemize}

\begin{methode}[Retrouver une formule électrique grâce à l'analogie]
Pour retrouver ou vérifier une formule du circuit $RLC$ :
\begin{enumerate}
\item Écris la formule mécanique correspondante, que tu connais bien (par exemple $T_0 = 2\pi\sqrt{m/k}$).
\item Applique le dictionnaire : $m \to L$, \; $k \to 1/C$, \; $f \to R$, \; $x \to q$, \; $v \to i$.
\item Simplifie : $2\pi\sqrt{\dfrac{L}{1/C}} = 2\pi\sqrt{LC}$.
\item Vérifie le résultat par analyse dimensionnelle.
\end{enumerate}
Cette démarche est un excellent réflexe de secours le jour de l'examen : si un doute survient sur une formule électrique, la mécanique te souffle la réponse.
\end{methode}

\begin{exoresolu}[Deux oscillateurs jumeaux]
On considère un solide de masse $m = 0,25~\mathrm{kg}$ accroché à un ressort de raideur $k = 100~\mathrm{N\,m^{-1}}$, et un circuit $LC$ constitué d'une bobine d'inductance $L = 0,25~\mathrm{H}$ et d'un condensateur de capacité $C = 10~\mathrm{mF}$. Montrer que ces deux oscillateurs ont la même pulsation propre et la même période propre.
\tcblower
\textbf{Oscillateur mécanique :}
\[
\omega_0 = \sqrt{\frac{k}{m}} = \sqrt{\frac{100}{0,25}} = \sqrt{400} = 20~\mathrm{rad\,s^{-1}}.
\]
\textbf{Oscillateur électrique :}
\[
\omega_0 = \frac{1}{\sqrt{LC}} = \frac{1}{\sqrt{0,25 \times 10\times 10^{-3}}} = \frac{1}{\sqrt{2,5\times 10^{-3}}} = \frac{1}{0,05} = 20~\mathrm{rad\,s^{-1}}.
\]
Les pulsations coïncident parfaitement. La période propre commune est :
\[
T_0 = \frac{2\pi}{\omega_0} = \frac{2\pi}{20} \approx 0,314~\mathrm{s}.
\]
\textbf{Conclusion.} La condition d'analogie est bien $\dfrac{k}{m} = \dfrac{1}{LC}$. Ici, $m=L=0,25$ et $k=100 = 1/C$ (avec $C=0,01~\mathrm{F}$). Le ressort et le circuit oscillent en parfaite correspondance.
\end{exoresolu}

\begin{savaistu}[Simuler un pont avec des composants électroniques]
Avant l'ère des supercalculateurs, les ingénieurs étudiaient la résistance des ponts et des bâtiments aux vibrations (vent, séismes, passage de convois) en construisant des \emph{circuits électriques analogues} : chaque masse devenait une bobine, chaque élasticité un condensateur, chaque amortisseur une résistance. Le réseau électrique, facile et peu coûteux à modifier, reproduisait fidèlement les oscillations de la structure. Ces « calculateurs analogiques » ont permis de dimensionner de grands ouvrages dans les années 1950--1970. Aujourd'hui encore, l'analogie électromécanique est utilisée pour concevoir les microphones, les haut-parleurs et les suspensions automobiles : un même formalisme mathématique gouverne des mondes très différents.
\end{savaistu}

\courssub{Synthèse de la leçon}

Rassemblons les acquis de toute la leçon en une carte mentale.

\begin{popfigure}
\begin{center}
\begin{tikzpicture}[font=\small, grow cyclic, mindmap,
  level 1 concept/.append style={level distance=3.4cm, sibling angle=90},
  level 2 concept/.append style={level distance=2.2cm, sibling angle=45}]
\node[concept, concept color=popblueL, font=\bfseries] {Oscillations électriques libres}
  child[concept color=poporangeL] { node[concept] {La bobine}
    child { node[concept, font=\scriptsize] {$u = ri + L\dfrac{di}{dt}$} }
    child { node[concept, font=\scriptsize] {$E_m = \tfrac12 Li^2$} }
  }
  child[concept color=popgreenL] { node[concept] {Circuit LC}
    child { node[concept, font=\scriptsize] {$\ddot{q}+\dfrac{q}{LC}=0$} }
    child { node[concept, font=\scriptsize] {$T_0 = 2\pi\sqrt{LC}$} }
    child { node[concept, font=\scriptsize] {$E = \tfrac12 Li^2 + \dfrac{q^2}{2C}$ = cte} }
  }
  child[concept color=poppinkL] { node[concept] {Circuit RLC}
    child { node[concept, font=\scriptsize] {$L\ddot{q}+R\dot{q}+\dfrac{q}{C}=0$} }
    child { node[concept, font=\scriptsize] {3 régimes selon $R$} }
  }
  child[concept color=popgoldL] { node[concept] {Analogies}
    child { node[concept, font=\scriptsize] {$q\leftrightarrow x$, $L\leftrightarrow m$} }
    child { node[concept, font=\scriptsize] {$\dfrac{1}{C}\leftrightarrow k$, $R\leftrightarrow f$} }
  };
\end{tikzpicture}
\end{center}
\legende{Carte mentale de synthèse : la bobine fournit l'élément d'inertie, le circuit LC l'oscillateur idéal, le RLC la réalité amortie, et les analogies relient le tout à la mécanique.}
\end{popfigure}

\begin{aretenir}[L'essentiel de la leçon]
\begin{itemize}
\item La bobine ($L$, $r$) : $u = ri + L\dfrac{di}{dt}$ ; elle stocke l'énergie magnétique $E_m = \dfrac{1}{2}Li^2$ et s'oppose aux variations du courant.
\item Circuit $LC$ idéal : $\ddot{q} + \dfrac{q}{LC} = 0$, solution $q(t) = Q_m\cos(\omega_0 t + \varphi)$ avec $\omega_0 = \dfrac{1}{\sqrt{LC}}$ et $T_0 = 2\pi\sqrt{LC}$ (formule de Thomson). L'énergie électromagnétique $E = \dfrac{1}{2}Li^2 + \dfrac{q^2}{2C}$ se conserve.
\item Circuit $RLC$ : $L\ddot{q} + R\dot{q} + \dfrac{q}{C} = 0$ ; l'énergie décroît à cause de l'effet Joule ; trois régimes : pseudo-périodique, critique ($R_c = 2\sqrt{L/C}$), apériodique.
\item Analogies : $q \leftrightarrow x$, $i \leftrightarrow v$, $L \leftrightarrow m$, $R \leftrightarrow f$, $\dfrac{1}{C} \leftrightarrow k$ — même équation, même physique, même énergie.
\end{itemize}
\end{aretenir}

\courssub{Ouverture : entretenir les oscillations (complément)}

Tout circuit réel possède une résistance : les oscillations libres finissent toujours par mourir, l'énergie électromagnétique étant irrémédiablement convertie en chaleur. Comment alors obtenir des oscillations électriques \emph{permanentes}, indispensables à la radio, aux téléphones mobiles (réseaux MTN, Orange) ou aux horloges électroniques ?

L'idée est simple et élégante : introduire dans le circuit un \cle{dispositif actif} (un amplificateur à transistor, par exemple) qui, à chaque période, restitue au circuit exactement l'énergie $Ri^2$ dissipée par effet Joule. Vu des bornes du dipôle $RLC$, ce dispositif se comporte comme une « résistance négative » qui compense $R$ : l'amortissement disparaît et les oscillations s'entretiennent d'elles-mêmes. C'est le principe des \cle{oscillateurs électroniques}, au cœur de tout émetteur radio.

Mais il existe une seconde manière, plus subtile, d'entretenir des oscillations : au lieu de laisser le circuit osciller à sa pulsation propre, on l'excite par un générateur extérieur dont on fait varier la fréquence. Lorsque la fréquence imposée s'approche de la fréquence propre du circuit, l'amplitude des oscillations devient considérable : c'est le phénomène de \cle{résonance}, qui permet à ton poste radio de sélectionner une station parmi des centaines. Ce sera l'objet de la prochaine leçon : les oscillations électriques \emph{forcées}.

\newpage

\coursec{Activité d'intégration}

\courssub{Objectifs de l'activité}

À l'issue de cette activité d'intégration, tu seras capable de :
\begin{itemize}
\item modéliser un circuit réel (antenne, bobine, condensateur) par un circuit $LC$ puis $RLC$ série et \cle{établir l'équation différentielle} régissant la charge $q(t)$ ;
\item \cle{dimensionner} un composant (ici la capacité $C$) pour accorder un circuit oscillant sur une fréquence imposée, en exploitant la formule de la période propre $T_0 = 2\pi\sqrt{LC}$ ;
\item calculer l'\cle{énergie électromagnétique} stockée dans le circuit et analyser sa décroissance en présence d'une résistance ;
\item déterminer la \cle{nature du régime} d'oscillation (pseudo-périodique, critique ou apériodique) par comparaison de la résistance totale à la résistance critique ;
\item argumenter, dans un langage technique correct, le rôle d'un dispositif d'entretien des oscillations.
\end{itemize}

La restitution se fera sous la forme d'une \textbf{fiche technique} rédigée, comme le ferait un technicien d'ENEO ou un installateur d'équipements de télécommunication.

\courssub{Situation}

\textbf{Conception du circuit d'accord d'un poste radio FM artisanal.}

Dans le quartier Nkolndongo à Yaoundé, Amina, élève de Terminale C passionnée d'électronique, participe au club scientifique de son lycée. Le club veut fabriquer un récepteur radio FM minimaliste capable de capter une station locale émettant à la fréquence $f_0 = \num{100,0}~\text{MHz}$. Le principe est connu : l'antenne capte toutes les ondes radio environnantes, mais seul un \cle{circuit oscillant} correctement \emph{accordé} sélectionne la station désirée, car il n'oscille efficacement qu'à sa fréquence propre.

Amina récupère dans les stocks du laboratoire :
\begin{itemize}
\item une \textbf{bobine} d'inductance $L = \num{0,22}~\mu\text{H}$, dont la résistance interne, mesurée à l'ohmmètre, vaut $r = \num{0,5}~\Omega$ ;
\item un \textbf{condensateur variable} (type « condensateur d'accord » des anciens postes à lampes) dont la capacité est réglable entre $C_{\min} = \num{5}~\text{pF}$ et $C_{\max} = \num{30}~\text{pF}$ ;
\item une antenne filaire qui, sous l'effet de l'onde radio, fournit au circuit une charge initiale estimée à $Q_0 = \num{1,0}~\text{nC}$ sur les armatures du condensateur.
\end{itemize}

Le professeur de physique pose le défi suivant : « Votre circuit oscillant est-il capable de sélectionner la station ? Les oscillations captées survivront-elles assez longtemps pour être exploitées ? Rédigez la fiche technique du montage, avec tous les calculs justificatifs. »

On rappelle que la bande FM camerounaise s'étend de $87{,}5$ à $108~\text{MHz}$, et que la sélectivité d'un récepteur (sa capacité à isoler une station parmi d'autres) repose précisément sur la finesse de la résonance du circuit d'accord.

\begin{popfigure}
\begin{center}
\begin{circuitikz}[european, scale=1.0, transform shape]
\draw (0,0) to[C=$C$, v={$u_C$}] (3,0) to[L=$L$, i={$i$}] (3,-2.5) to[R=$r$, v={$u_r$}] (0,-2.5) -- (0,0);
\end{circuitikz}
\end{center}
\legende{Circuit oscillant série du récepteur : condensateur d'accord $C$, bobine d'inductance $L$ et de résistance interne $r$. Dans le modèle idéal, on néglige $r$ ; la flèche $i$ indique le sens positif choisi, avec $i = \dfrac{dq}{dt}$.}
\end{popfigure}

\courssub{Consignes}

Ta fiche technique devra répondre, dans l'ordre, aux questions suivantes. Chaque résultat numérique sera donné avec son unité SI et un nombre raisonnable de chiffres significatifs.

\begin{enumerate}
\item \textbf{Modélisation.} Schématiser le circuit formé par l'antenne (assimilée à la source de la charge initiale), la bobine idéale d'inductance $L$ et le condensateur de capacité $C$. En appliquant la loi des mailles, établir l'équation différentielle vérifiée par la charge $q(t)$ du condensateur dans le cas idéal ($r$ négligée). Préciser la nature des solutions.
\item \textbf{Accord du circuit.} Exprimer la fréquence propre $f_0$ en fonction de $L$ et $C$, puis calculer la valeur de $C$ qui accorde le circuit sur la station à $\num{100,0}~\text{MHz}$. Cette valeur est-elle accessible avec le condensateur variable d'Amina ? Conclure sur la faisabilité.
\item \textbf{Caractéristiques temporelles et énergétiques.} Calculer la période propre $T_0$ des oscillations. En supposant le circuit idéal, calculer l'énergie électromagnétique maximale $E$ stockée dans le circuit lorsque la charge initiale vaut $Q_0 = \num{1,0}~\text{nC}$. Où se trouve cette énergie à l'instant initial ?
\item \textbf{Prise en compte de la résistance.} La bobine réelle possède une résistance $r = \num{0,5}~\Omega$. Écrire l'équation différentielle du circuit $RLC$ série ainsi obtenu. Expliquer qualitativement pourquoi l'énergie électromagnétique décroît au cours du temps. Calculer la résistance critique $R_C = 2\sqrt{L/C}$ et la comparer à $r$. Quel est le régime d'oscillation ? Le circuit reste-t-il utilisable comme circuit d'accord ?
\item \textbf{Entretien des oscillations.} Proposer une justification physique du rôle d'un amplificateur (ou d'un montage à résistance négative) dans un récepteur réel : que doit-il compenser, et à quel rythme ?
\item \textbf{Rédaction.} Présenter l'ensemble sous forme d'une fiche technique : schéma, tableau des grandeurs calculées, conclusion sur la faisabilité et les limites du montage.
\end{enumerate}

\courssub{Ressources à mobiliser}

\begin{aretenir}[Boîte à outils de la leçon]
\begin{itemize}
\item Bobine : $u_L = ri + L\dfrac{di}{dt}$ ; énergie magnétique $E_L = \dfrac{1}{2}Li^2$.
\item Condensateur : $u_C = \dfrac{q}{C}$ ; énergie électrique $E_C = \dfrac{q^2}{2C}$ ; $i = \dfrac{dq}{dt}$.
\item Circuit $LC$ idéal : loi des mailles $\Rightarrow$ $\ddot{q} + \dfrac{1}{LC}q = 0$ ; solutions $q(t) = Q_m\cos(\omega_0 t + \varphi)$ avec $\omega_0 = \dfrac{1}{\sqrt{LC}}$, $T_0 = 2\pi\sqrt{LC}$, $f_0 = \dfrac{1}{2\pi\sqrt{LC}}$.
\item Conservation : $E = E_C + E_L = \dfrac{Q_m^2}{2C} = \text{constante}$ (circuit idéal).
\item Circuit $RLC$ série : $\ddot{q} + \dfrac{R}{L}\dot{q} + \dfrac{1}{LC}q = 0$ ; l'énergie décroît par effet Joule ($P_J = Ri^2$).
\item Régimes : $R < R_C$ pseudo-périodique ; $R = R_C$ critique ; $R > R_C$ apériodique, avec $R_C = 2\sqrt{\dfrac{L}{C}}$.
\end{itemize}
\end{aretenir}

\begin{attention}[Pièges fréquents dans ce type de problème]
\begin{itemize}
\item \textbf{Unités.} Les capacités des circuits d'accord sont de l'ordre du picofarad ($1~\text{pF} = 10^{-12}~\text{F}$) et les inductances du microhenry ($1~\mu\text{H} = 10^{-6}~\text{H}$) : convertis \emph{toujours} en unités SI avant tout calcul.
\item \textbf{Signe de la loi des mailles.} En convention récepteur avec $i = \dfrac{dq}{dt}$, la maille fermée impose $u_L + u_C = 0$ (circuit idéal) : le signe moins est absorbé par le choix d'orientation des flèches de tension. Une erreur de signe conduit à $\ddot{q} - \dfrac{q}{LC} = 0$, dont les solutions ne sont \emph{pas} sinusoïdales — c'est un bon auto-contrôle.
\item \textbf{Résistance critique.} C'est la résistance \emph{totale} du circuit (ici $r$ seule) qu'il faut comparer à $R_C = 2\sqrt{L/C}$, et non pas à $\sqrt{L/C}$.
\end{itemize}
\end{attention}

\courssub{Démarche et corrigé commenté}

\begin{exoresolu}[Fiche technique du circuit d'accord FM — corrigé complet]
Énoncé : dimensionner le circuit $LC$ d'un récepteur FM ($f_0 = \num{100,0}~\text{MHz}$, $L = \num{0,22}~\mu\text{H}$, $r = \num{0,5}~\Omega$, $C$ réglable de $5$ à $30~\text{pF}$, $Q_0 = \num{1,0}~\text{nC}$), étudier son régime réel et justifier l'entretien des oscillations.
\tcblower
\textbf{Étape 1 — Équation différentielle du circuit $LC$ idéal.}

Le circuit série comporte la bobine idéale ($u_L = L\dfrac{di}{dt}$) et le condensateur ($u_C = \dfrac{q}{C}$). La loi des mailles s'écrit $u_L + u_C = 0$. Avec $i = \dfrac{dq}{dt}$, donc $\dfrac{di}{dt} = \ddot{q}$ :
\[
L\ddot{q} + \frac{q}{C} = 0
\qquad\Longleftrightarrow\qquad
\ddot{q} + \frac{1}{LC}\,q = 0.
\]
C'est l'équation d'un \cle{oscillateur harmonique} de pulsation propre $\omega_0 = \dfrac{1}{\sqrt{LC}}$. Vérification dimensionnelle : $[LC] = \text{H}\cdot\text{F} = \text{s}^2$, donc $[1/\sqrt{LC}] = \text{s}^{-1}$, homogène à une pulsation. Les solutions sont sinusoïdales : $q(t) = Q_m\cos(\omega_0 t + \varphi)$, et $i(t) = \dot{q} = -\omega_0 Q_m\sin(\omega_0 t + \varphi)$.

\textbf{Étape 2 — Calcul de la capacité d'accord.}

De $f_0 = \dfrac{1}{2\pi\sqrt{LC}}$, on tire en élevant au carré :
\[
C = \frac{1}{4\pi^2 f_0^2 L}.
\]
Application numérique avec $f_0 = \num{1,00e8}~\text{Hz}$ et $L = \num{2,2e-7}~\text{H}$ :
\begin{align*}
C &= \frac{1}{4\pi^2 \times (\num{1,00e8})^2 \times \num{2,2e-7}} \\
&= \frac{1}{39{,}48 \times \num{1,00e16} \times \num{2,2e-7}} \\
&= \frac{1}{\num{8,68e10}} \approx \num{1,15e-11}~\text{F} \approx \num{11,5}~\text{pF}.
\end{align*}
Or $C_{\min} = \num{5}~\text{pF} < \num{11,5}~\text{pF} < \num{30}~\text{pF} = C_{\max}$ : la valeur est \textbf{accessible}. Le montage est faisable ; il suffit de régler le condensateur variable vers le milieu de sa course. Remarquons au passage que la plage $5$--$30~\text{pF}$ permet d'accorder le circuit de $f = \dfrac{1}{2\pi\sqrt{LC_{\max}}} \approx \num{72}~\text{MHz}$ à $f = \dfrac{1}{2\pi\sqrt{LC_{\min}}} \approx \num{152}~\text{MHz}$ : toute la bande FM camerounaise ($87{,}5$ à $108~\text{MHz}$) est couverte.

\textbf{Étape 3 — Période propre et énergie maximale.}

Par construction de l'accord :
\[
T_0 = \frac{1}{f_0} = \frac{1}{\num{1,00e8}} = \num{1,0e-8}~\text{s} = \num{10}~\text{ns}.
\]
(On vérifie : $T_0 = 2\pi\sqrt{LC} = 2\pi\sqrt{\num{2,2e-7}\times\num{1,15e-11}} \approx \num{1,0e-8}~\text{s}$.)

À $t = 0$, l'intensité est nulle et toute l'énergie est stockée sous forme électrique dans le condensateur. L'énergie électromagnétique totale, conservée dans le circuit idéal, vaut :
\[
E = \frac{Q_0^2}{2C} = \frac{(\num{1,0e-9})^2}{2\times\num{1,15e-11}} = \frac{\num{1,0e-18}}{\num{2,3e-11}} \approx \num{4,3e-8}~\text{J} \approx \num{43}~\text{nJ}.
\]
Cette énergie, certes faible, oscille intégralement entre le condensateur et la bobine : un quart de période plus tard ($t = T_0/4$), elle est entièrement magnétique, avec $i_{\max} = \omega_0 Q_0 = 2\pi f_0 Q_0 \approx \num{0,63}~\text{A}$. On vérifie la cohérence énergétique : $\dfrac{1}{2}Li_{\max}^2 = \dfrac{1}{2}\times\num{2,2e-7}\times(0{,}63)^2 \approx \num{4,3e-8}~\text{J} = E$.

\textbf{Étape 4 — Circuit réel : amortissement et régime.}

Avec la résistance $r$ de la bobine, la loi des mailles donne $u_L + u_r + u_C = 0$, soit :
\[
L\ddot{q} + r\dot{q} + \frac{q}{C} = 0
\qquad\Longleftrightarrow\qquad
\ddot{q} + \frac{r}{L}\dot{q} + \frac{1}{LC}q = 0.
\]
En dérivant l'énergie totale $E = \dfrac{q^2}{2C} + \dfrac{1}{2}Li^2$ par rapport au temps, on obtient :
\[
\frac{dE}{dt} = \frac{q}{C}\dot{q} + Li\,\frac{di}{dt} = i\left(\frac{q}{C} + L\ddot{q}\right) = i(-ri) = -ri^2 \leq 0 :
\]
l'énergie décroît par \cle{effet Joule} dans la résistance de la bobine. Les oscillations sont donc amorties.

Calculons la résistance critique :
\[
R_C = 2\sqrt{\frac{L}{C}} = 2\sqrt{\frac{\num{2,2e-7}}{\num{1,15e-11}}} = 2\sqrt{\num{1,91e4}} \approx 2\times 138 \approx \num{2,8e2}~\Omega.
\]
Comparaison : $r = \num{0,5}~\Omega \ll R_C \approx \num{280}~\Omega$. Le régime est donc \textbf{pseudo-périodique}, et même très faiblement amorti : la pseudo-période est pratiquement égale à $T_0 = \num{10}~\text{ns}$. Le circuit oscille de nombreuses périodes avant de s'éteindre (le facteur de qualité vaut $Q = \dfrac{1}{r}\sqrt{\dfrac{L}{C}} \approx \dfrac{138}{\num{0,5}} \approx 276$) : il est \textbf{utilisable comme circuit d'accord}, et sa sélectivité est même excellente.

\textbf{Étape 5 — Rôle de l'amplificateur.}

L'antenne reçoit en permanence l'onde de la station, mais l'énergie captée est minuscule et s'évanouit par effet Joule. Un amplificateur (transistor alimenté par une pile, par exemple) prélève de l'énergie à une source extérieure et la réinjecte dans le circuit \emph{en phase} avec les oscillations, compensant exactement, à chaque période, l'énergie dissipée dans $r$. Il se comporte comme une « résistance négative » $-r$ qui annule l'amortissement : les oscillations sont \cle{entretenues}, d'amplitude constante, et le signal devient exploitable par l'étage audio du poste.

\textbf{Étape 6 — Fiche technique (synthèse).}

\begin{center}
\begin{tabularx}{\linewidth}{|l|X|}\hline
\rowcolor{popblueL} \textbf{Grandeur} & \textbf{Valeur} \\ \hline
Fréquence d'accord $f_0$ & $\num{100,0}~\text{MHz}$ \\ \hline
Inductance $L$ & $\num{0,22}~\mu\text{H}$ \\ \hline
Capacité calculée $C$ & $\num{11,5}~\text{pF}$ (accessible : $5$ à $30~\text{pF}$) \\ \hline
Période propre $T_0$ & $\num{10}~\text{ns}$ \\ \hline
Énergie maximale $E$ & $\num{4,3e-8}~\text{J}$ \\ \hline
Intensité maximale $i_{\max}$ & $\approx \num{0,63}~\text{A}$ \\ \hline
Résistance critique $R_C$ & $\approx \num{280}~\Omega$ \\ \hline
Régime réel ($r = \num{0,5}~\Omega$) & pseudo-périodique, très faiblement amorti \\ \hline
Dispositif requis & amplificateur d'entretien \\ \hline
\end{tabularx}
\end{center}

Conclusion : le montage d'Amina est faisable et sélectif ; l'ajout d'un étage amplificateur est indispensable pour entretenir les oscillations et rendre la station audible.
\end{exoresolu}

\courssub{Bilan de l'activité}

Cette activité a permis de mobiliser, dans une situation technologique authentique, \textbf{l'ensemble des savoir-faire de la leçon} :
\begin{itemize}
\item la \textbf{modélisation} : passer d'un dispositif réel (antenne, bobine, condensateur) au schéma $LC$ puis $RLC$, et établir rigoureusement l'équation différentielle par la loi des mailles ;
\item l'\textbf{exploitation de la période propre} : la formule $T_0 = 2\pi\sqrt{LC}$ n'est pas une simple formule à réciter, c'est un outil de \emph{conception} qui permet de calculer un composant pour atteindre une fréquence cible ;
\item l'\textbf{étude énergétique} : l'énergie électromagnétique se conserve dans le circuit idéal, mais décroît par effet Joule dès qu'une résistance intervient ($\dfrac{dE}{dt} = -ri^2$), ce qui motive physiquement l'entretien des oscillations ;
\item le \textbf{diagnostic de régime} : la comparaison de $r$ à $R_C = 2\sqrt{L/C}$ tranche quantitativement entre pseudo-périodique, critique et apériodique — ici, un circuit d'accord performant exige $r \ll R_C$.
\end{itemize}

Tu as également découvert une idée essentielle en télécommunications : la \cle{sélectivité} d'un récepteur repose sur la résonance d'un circuit oscillant faiblement amorti. Le même principe se retrouve dans les téléphones mobiles, les téléviseurs et les liaisons satellite de CAMTEL. Enfin, la rédaction d'une fiche technique t'a entraîné à communiquer des résultats scientifiques de manière structurée, une compétence précieuse bien au-delà du baccalauréat.

\newpage

\coursec{Méthodes et exercices résolus}

\coursec{Oscillations électriques libres (circuits LC et RLC)}

\courssub{La bobine : inductance, tension et énergie magnétique}

\begin{definition}[Bobine réelle et bobine idéale]
Une \textbf{bobine réelle} est modélisée par une inductance $L$ (en henry, H) en série avec une résistance $r$ (en ohm, $\Omega$). La tension $u$ à ses bornes (convention récepteur) s'écrit :
\[ u = r\,i + L\,\frac{di}{dt}. \]
Une \textbf{bobine idéale} a une résistance nulle ($r=0$) : $u_L = L\,\dfrac{di}{dt}$.
\end{definition}

\begin{propriete}[Énergie magnétique stockée dans une bobine]
L'énergie magnétique emmagasinée dans le champ magnétique d'une bobine d'inductance $L$ parcourue par un courant $i$ est :
\[ E_L = \frac{1}{2}\,L\,i^2. \]
\emph{Démonstration exigible au Bac :} Puissance reçue $P = u_L i = L\,i\,\dfrac{di}{dt} = \dfrac{d}{dt}\left(\frac{1}{2}Li^2\right)$. Intégration entre $0$ et $t$ (avec $i(0)=0$) donne $E_L = \frac{1}{2}Li^2$.
\end{propriete}

\begin{aretenir}[Rappel : Condensateur]
Tension aux bornes (convention récepteur) : $u_C = \dfrac{q}{C}$.
Énergie électrique stockée : $E_C = \dfrac{1}{2}\dfrac{q^2}{C} = \dfrac{1}{2}C\,u_C^2$.
Relation courant-charge : $i = \dfrac{dq}{dt}$ (si $i$ arrive sur l'armature de charge $q$).
\end{aretenir}

\courssub{Le circuit LC idéal : équation différentielle des oscillations libres}

\begin{methode}[Établir l'équation différentielle des oscillations libres]
\begin{enumerate}
\item \textbf{Choisir les conventions} : orienter le circuit (sens positif de $i$) et flécher la tension $u_C$ aux bornes du condensateur en convention \cle{récepteur}, de sorte que $i = \dfrac{dq}{dt}$ avec $q$ la charge de l'armature vers laquelle arrive $i$.
\item \textbf{Écrire la loi d'additivité des tensions} (loi des mailles) : $u_L + u_C = 0$ pour le circuit LC idéal (résistance négligée).
\item \textbf{Exprimer chaque tension} : $u_C = \dfrac{q}{C}$ et, pour la bobine idéale, $u_L = L\dfrac{di}{dt} = L\dfrac{d^2q}{dt^2}$.
\item \textbf{Substituer et ordonner} : on obtient
\[ \ddot{q} + \frac{1}{LC}\,q = 0 \qquad \text{soit} \qquad \ddot{q} + \omega_0^2\, q = 0 \quad \text{avec} \quad \omega_0 = \frac{1}{\sqrt{LC}}. \]
\item \textbf{Vérifier l'homogénéité} : $[LC] = \text{s}^2$ (car $T_0 = 2\pi\sqrt{LC}$ est une durée), donc $\dfrac{1}{LC}$ est bien homogène à $\text{s}^{-2}$, comme $\ddot{q}/q$.
\end{enumerate}
\textbf{Réflexe} : si la bobine a une résistance $r$ non négligeable, ajouter le terme $r\,i = r\dot{q}$, ce qui donne l'équation du circuit RLC amorti.
\end{methode}

\begin{exoresolu}[Circuit LC : mise en équation]
On réalise un circuit comprenant un condensateur de capacité $C = \SI{10}{\micro F}$ initialement chargé sous la tension $U_0 = \SI{6,0}{V}$, et une bobine d'inductance $L = \SI{0,10}{H}$ de résistance négligeable. À $t = 0$, on ferme l'interrupteur reliant le condensateur à la bobine.
\begin{enumerate}
\item Établir l'équation différentielle vérifiée par la charge $q(t)$ du condensateur.
\item En déduire la valeur de la pulsation propre $\omega_0$ des oscillations.
\end{enumerate}
\tcblower
\textbf{1. Équation différentielle.}

On oriente le circuit dans le sens du courant $i$ et on flèche $u_C$ en convention récepteur : $i = \dfrac{dq}{dt}$.

La loi d'additivité des tensions dans la maille fermée s'écrit :
\[ u_C + u_L = 0. \]
Or $u_C = \dfrac{q}{C}$ (relation charge-tension du condensateur) et, la résistance de la bobine étant négligeable, $u_L = L\dfrac{di}{dt} = L\dfrac{d^2q}{dt^2} = L\ddot{q}$.

En substituant :
\[ \frac{q}{C} + L\ddot{q} = 0 \quad \Longleftrightarrow \quad \boxed{\ddot{q} + \frac{1}{LC}\,q = 0}. \]

C'est l'équation différentielle d'un \textbf{oscillateur harmonique} de pulsation propre $\omega_0 = \dfrac{1}{\sqrt{LC}}$.

\textbf{2. Pulsation propre.}
\[ \omega_0 = \frac{1}{\sqrt{LC}} = \frac{1}{\sqrt{\num{0,10}\times \num{1,0e-5}}} = \frac{1}{\sqrt{\num{1,0e-6}}} = \frac{1}{\num{1,0e-3}} = \SI{1,0e3}{rad.s^{-1}}. \]

\textbf{Conclusion} : la charge du condensateur oscille sinusoïdalement avec la pulsation propre $\omega_0 = \SI{1,0e3}{rad.s^{-1}}$, soit une période propre $T_0 = \dfrac{2\pi}{\omega_0} \approx \SI{6,3}{ms}$.
\end{exoresolu}

\courssub{Expressions de q(t) et i(t) — Conditions initiales}

\begin{methode}[Déterminer les expressions de q(t) et i(t)]
\begin{enumerate}
\item \textbf{Écrire la solution générale} de $\ddot{q} + \omega_0^2 q = 0$ :
\[ q(t) = Q_m\cos(\omega_0 t + \varphi) \qquad \text{avec} \quad \omega_0 = \frac{1}{\sqrt{LC}}. \]
\item \textbf{Traduire les conditions initiales} : à $t = 0$, $q(0) = Q_0$ (charge initiale du condensateur, souvent $Q_0 = CU_0$) et $i(0) = 0$ (l'intensité dans une bobine ne peut pas subir de discontinuité).
\item \textbf{Exploiter $q(0)$} : $Q_m\cos\varphi = Q_0$.
\item \textbf{Exploiter $i(0) = 0$} : $i(t) = \dot{q}(t) = -\omega_0 Q_m\sin(\omega_0 t + \varphi)$, donc $i(0) = -\omega_0 Q_m \sin\varphi = 0$, d'où $\varphi = 0$ (modulo $\pi$ ; le signe de $Q_0$ lève l'ambiguïté).
\item \textbf{Conclure} : si le condensateur est chargé à $Q_0 > 0$ puis isolé de son chargeur,
\[ q(t) = Q_0\cos(\omega_0 t), \qquad i(t) = -\omega_0 Q_0\sin(\omega_0 t) = -I_m\sin(\omega_0 t), \]
avec $I_m = \omega_0 Q_0 = \dfrac{Q_0}{\sqrt{LC}}$. L'intensité est en \cle{quadrature} (déphasage de $\dfrac{\pi}{2}$) avec la charge.
\end{enumerate}
\textbf{Réflexe} : vérifier toujours que $q(0)$ et $i(0)$ retrouvés à partir des expressions finales correspondent bien aux conditions initiales de l'énoncé.
\end{methode}

\begin{exoresolu}[Expressions de q(t) et i(t)]
Un condensateur de capacité $C = \SI{4,7}{\micro F}$ est chargé sous une tension $U_0 = \SI{12}{V}$, puis connecté à l'instant $t = 0$ à une bobine idéale d'inductance $L = \SI{50}{mH}$.
\begin{enumerate}
\item Calculer la charge initiale $Q_0$ du condensateur et la pulsation propre $\omega_0$.
\item Donner les expressions numériques de $q(t)$ et $i(t)$.
\item Calculer l'intensité maximale $I_m$.
\end{enumerate}
\tcblower
\textbf{1. Charge initiale et pulsation propre.}
\[ Q_0 = CU_0 = \num{4,7e-6}\times 12 = \SI{5,64e-5}{C} \approx \SI{56}{\micro C}. \]
\[ \omega_0 = \frac{1}{\sqrt{LC}} = \frac{1}{\sqrt{\num{5,0e-2}\times \num{4,7e-6}}} = \frac{1}{\sqrt{\num{2,35e-7}}} \approx \SI{2,06e3}{rad.s^{-1}}. \]

\textbf{2. Expressions de q(t) et i(t).}

La solution générale est $q(t) = Q_m\cos(\omega_0 t + \varphi)$.

Conditions initiales : $q(0) = Q_0$ et $i(0) = 0$ (continuité du courant dans la bobine).

$i(0) = -\omega_0 Q_m\sin\varphi = 0 \Rightarrow \varphi = 0$ (ou $\pi$). Comme $q(0) = Q_m\cos\varphi = Q_0 > 0$, on retient $\varphi = 0$ et $Q_m = Q_0$. D'où :
\[ q(t) = \num{5,64e-5}\cos(\num{2,06e3}\,t) \ \text{(en C)}, \]
\[ i(t) = \dot{q}(t) = -\omega_0 Q_0\sin(\omega_0 t) = -\num{0,116}\sin(\num{2,06e3}\,t) \ \text{(en A)}. \]

\textbf{3. Intensité maximale.}
\[ I_m = \omega_0 Q_0 = \num{2,06e3}\times \num{5,64e-5} \approx \SI{0,116}{A} = \SI{116}{mA}. \]

\textbf{Conclusion} : charge et intensité oscillent sinusoïdalement en quadrature ; l'intensité est maximale ($\SI{116}{mA}$) lorsque le condensateur est complètement déchargé.
\end{exoresolu}

\courssub{Période propre, fréquence et applications (Formule de Thomson)}

\begin{aretenir}[Formule de Thomson — Période et fréquence propres du circuit LC]
\[ T_0 = 2\pi\sqrt{LC}, \qquad f_0 = \frac{1}{T_0} = \frac{1}{2\pi\sqrt{LC}}, \qquad \omega_0 = \frac{2\pi}{T_0} = \frac{1}{\sqrt{LC}}. \]
La période propre $T_0$ ne dépend \emph{que} de $L$ et $C$, jamais de la tension de charge initiale $U_0$ ni de l'amplitude des oscillations.
\end{aretenir}

\begin{methode}[Calculer et exploiter $T_0$, $f_0$ et les paramètres du circuit]
\begin{enumerate}
\item \textbf{Convertir toutes les données en unités SI} avant tout calcul : $L$ en henry (H), $C$ en farad (F), $T_0$ en seconde (s).
\item \textbf{Appliquer la formule de Thomson} : $T_0 = 2\pi\sqrt{LC}$.
\item \textbf{Isoler la grandeur cherchée} si l'énoncé demande $L$ ou $C$ : en élevant au carré, $L = \dfrac{T_0^2}{4\pi^2 C}$ ou $C = \dfrac{T_0^2}{4\pi^2 L}$.
\item \textbf{Vérifier par analyse dimensionnelle} : $[L] = \text{V}\cdot\text{s}\cdot\text{A}^{-1}$ et $[C] = \text{A}\cdot\text{s}\cdot\text{V}^{-1}$, donc $[LC] = \text{s}^2$ et $\sqrt{LC}$ est bien une durée.
\item \textbf{Exploiter un oscillogramme} : mesurer la durée de $n$ périodes sur l'écran, diviser par $n$ pour réduire l'incertitude, puis identifier $T_0$.
\end{enumerate}
\textbf{Réflexe Bac} : la période propre $T_0$ ne dépend \emph{que} de $L$ et $C$, jamais de la tension de charge initiale $U_0$ : c'est une question piège classique.
\end{methode}

\begin{exoresolu}[Accord d'un poste radio]
La partie « réception » d'un poste radio utilisé dans un village de la région de l'Est comporte un circuit LC dont la bobine a pour inductance $L = \SI{2,0}{\micro H}$. Pour capter une station émettant à la fréquence $f_0 = \SI{100}{MHz}$ (bande FM), on ajuste un condensateur de capacité $C$ variable.
\begin{enumerate}
\item Rappeler la relation entre $f_0$, $L$ et $C$.
\item Calculer la valeur de $C$ permettant l'accord sur cette station.
\end{enumerate}
\tcblower
\textbf{1. Relation.} La fréquence propre du circuit LC est donnée par la formule de Thomson :
\[ f_0 = \frac{1}{2\pi\sqrt{LC}}. \]

\textbf{2. Calcul de C.} On élève au carré et on isole $C$ :
\[ f_0^2 = \frac{1}{4\pi^2 LC} \quad \Longrightarrow \quad C = \frac{1}{4\pi^2 L f_0^2}. \]

Application numérique avec $L = \num{2,0e-6}\ \text{H}$ et $f_0 = \num{1,0e8}\ \text{Hz}$ :
\[ C = \frac{1}{4\pi^2 \times \num{2,0e-6}\times (\num{1,0e8})^2} = \frac{1}{4\pi^2 \times \num{2,0e10}} \approx \frac{1}{\num{7,90e11}} \approx \SI{1,27e-12}{F}. \]

\textbf{Conclusion} : il faut régler le condensateur sur $C \approx \SI{1,3}{pF}$. Cette très faible capacité explique pourquoi, en FM, on utilise des condensateurs variables de quelques picofarads : c'est le principe du bouton de recherche de stations.
\end{exoresolu}

\courssub{Étude énergétique du circuit LC : conservation de l'énergie électromagnétique}

\begin{methode}[Exploiter la conservation de l'énergie électromagnétique]
\begin{enumerate}
\item \textbf{Écrire les trois énergies} :
\[ E_C = \frac{1}{2}\frac{q^2}{C} = \frac{1}{2}Cu_C^2 \quad (\text{électrique}), \qquad E_L = \frac{1}{2}Li^2 \quad (\text{magnétique}), \qquad E = E_C + E_L. \]
\item \textbf{Montrer la conservation} dans le LC idéal : dériver $E$ par rapport au temps et utiliser l'équation différentielle :
\[ \frac{dE}{dt} = \frac{q\dot{q}}{C} + Li\,\frac{di}{dt} = i\left(\frac{q}{C} + L\ddot{q}\right) = 0. \]
\item \textbf{Exploiter les situations extrêmes} : quand $q = \pm Q_m$, $i = 0$ et $E = \dfrac{Q_m^2}{2C}$ ; quand $i = \pm I_m$, $q = 0$ et $E = \dfrac{1}{2}LI_m^2$. La conservation donne la relation utile :
\[ \frac{Q_m^2}{2C} = \frac{1}{2}LI_m^2 \quad \Longrightarrow \quad I_m = \frac{Q_m}{\sqrt{LC}} = \omega_0 Q_m. \]
\item \textbf{Pour un état intermédiaire}, écrire $E = \dfrac{q^2}{2C} + \dfrac{1}{2}Li^2 = \dfrac{Q_m^2}{2C}$ et isoler la grandeur inconnue.
\item \textbf{Conclure sur les échanges} : l'énergie se transfère alternativement du condensateur vers la bobine et réciproquement, deux fois par période (les énergies oscillent à la pulsation $2\omega_0$).
\end{enumerate}
\end{methode}

\begin{exoresolu}[Bilan énergétique d'un circuit LC]
Un condensateur de capacité $C = \SI{20}{\micro F}$, chargé sous $U_0 = \SI{9,0}{V}$, est connecté à une bobine idéale d'inductance $L = \SI{0,20}{H}$.
\begin{enumerate}
\item Calculer l'énergie électromagnétique totale $E$ du circuit.
\item En déduire l'intensité maximale $I_m$ dans le circuit.
\item À un instant $t$, la tension aux bornes du condensateur vaut $u_C = \SI{5,0}{V}$. Calculer l'intensité $i$ à cet instant.
\end{enumerate}
\tcblower
\textbf{1. Énergie totale.} À $t = 0$, toute l'énergie est stockée dans le condensateur ($i(0) = 0$) :
\[ E = \frac{1}{2}CU_0^2 = \frac{1}{2}\times \num{2,0e-5}\times 9,0^2 = \SI{8,1e-4}{J} = \SI{0,81}{mJ}. \]
La résistance étant nulle, cette énergie se conserve au cours des oscillations.

\textbf{2. Intensité maximale.} Lorsque le condensateur est déchargé, toute l'énergie est magnétique :
\[ \frac{1}{2}LI_m^2 = E \quad \Longrightarrow \quad I_m = \sqrt{\frac{2E}{L}} = \sqrt{\frac{2\times \num{8,1e-4}}{\num{0,20}}} = \sqrt{\num{8,1e-3}} = \SI{9,0e-2}{A}. \]
Soit $I_m = \SI{90}{mA}$.

\textbf{3. Intensité lorsque $u_C = \SI{5,0}{V}$.} Conservation de l'énergie :
\[ \frac{1}{2}Cu_C^2 + \frac{1}{2}Li^2 = E \quad \Longrightarrow \quad i^2 = \frac{2E - Cu_C^2}{L}. \]
Calcul de l'énergie électrique restante :
\[ \frac{1}{2}Cu_C^2 = \frac{1}{2}\times \num{2,0e-5}\times 25 = \SI{2,5e-4}{J}. \]
D'où :
\[ i^2 = \frac{2(\num{8,1e-4} - \num{2,5e-4})}{\num{0,20}} = \frac{\num{1,12e-3}}{\num{0,20}} = \num{5,6e-3} \quad \Longrightarrow \quad |i| \approx \SI{7,5e-2}{A} = \SI{75}{mA}. \]

\textbf{Conclusion} : $E = \SI{0,81}{mJ}$, $I_m = \SI{90}{mA}$, et lorsque $u_C = \SI{5,0}{V}$ le courant vaut $\SI{75}{mA}$ en valeur absolue (son signe dépend du sens de l'échange à cet instant).
\end{exoresolu}

\courssub{Le circuit RLC série : oscillations amorties}

\begin{aretenir}[Équation différentielle et régimes du circuit RLC série]
Équation différentielle (loi des mailles : $u_L + u_R + u_C = 0$) :
\[ L\ddot{q} + R\dot{q} + \frac{q}{C} = 0 \qquad \Longleftrightarrow \qquad \ddot{q} + \frac{R}{L}\dot{q} + \frac{1}{LC}q = 0. \]
Résistance critique : $R_C = 2\sqrt{\dfrac{L}{C}}$.
\begin{itemize}
\item $R < R_C$ : régime \cle{pseudo-périodique} (oscillations d'amplitude décroissante), pseudo-pulsation $\omega = \sqrt{\omega_0^2 - \left(\dfrac{R}{2L}\right)^2}$, pseudo-période $T = \dfrac{2\pi}{\omega}$.
\item $R = R_C$ : régime \cle{critique} (retour à l'équilibre le plus rapide, sans oscillation).
\item $R > R_C$ : régime \cle{apériodique} (retour lent à l'équilibre, sans oscillation).
\end{itemize}
Bilan énergétique : $\dfrac{dE}{dt} = -R\,i^2 < 0$ (énergie dissipée par effet Joule).
\end{aretenir}

\begin{methode}[Étudier les oscillations amorties du circuit RLC]
\begin{enumerate}
\item \textbf{Établir l'équation différentielle} en ajoutant la tension $u_R = Ri$ aux bornes du résistor.
\item \textbf{Calculer la résistance critique} $R_C = 2\sqrt{L/C}$ et comparer $R$ à $R_C$ pour identifier le régime.
\item \textbf{En régime pseudo-périodique}, calculer la pseudo-période $T = \dfrac{2\pi}{\sqrt{\omega_0^2 - (R/2L)^2}}$ ; si l'amortissement est faible ($R \ll R_C$), alors $T \approx T_0 = 2\pi\sqrt{LC}$.
\item \textbf{Interpréter énergétiquement} : $\dfrac{dE}{dt} = -Ri^2 < 0$ : l'énergie électromagnétique décroît, dissipée par effet Joule dans le résistor.
\item \textbf{Pour entretenir les oscillations}, il faut un dispositif qui compense à chaque instant l'énergie dissipée (montage à résistance négative, amplificateur) : c'est le principe des oscillateurs entretenus.
\end{enumerate}
\end{methode}

\begin{exoresolu}[Régimes d'un circuit RLC série]
On associe en série un condensateur $C = \SI{1,0}{\micro F}$ chargé, une bobine $L = \SI{0,10}{H}$ de résistance négligeable et un résistor de résistance $R$ réglable.
\begin{enumerate}
\item Établir l'équation différentielle vérifiée par $q(t)$.
\item Calculer la résistance critique $R_C$.
\item On fixe $R = \SI{50}{\Omega}$. Quel est le régime observé ? Comparer la pseudo-période $T$ à la période propre $T_0$.
\end{enumerate}
\tcblower
\textbf{1. Équation différentielle.} Loi des mailles : $u_L + u_R + u_C = 0$, soit :
\[ L\ddot{q} + R\dot{q} + \frac{q}{C} = 0 \quad \Longleftrightarrow \quad \ddot{q} + \frac{R}{L}\dot{q} + \frac{1}{LC}q = 0. \]

\textbf{2. Résistance critique.}
\[ R_C = 2\sqrt{\frac{L}{C}} = 2\sqrt{\frac{\num{0,10}}{\num{1,0e-6}}} = 2\sqrt{\num{1,0e5}} = 2\times \num{3,16e2} \approx \SI{6,3e2}{\Omega}. \]
Soit $R_C \approx \SI{630}{\Omega}$.

\textbf{3. Régime pour $R = \SI{50}{\Omega}$.} Comme $R = \SI{50}{\Omega} < R_C \approx \SI{630}{\Omega}$, le régime est \textbf{pseudo-périodique} : on observe des oscillations dont l'amplitude décroît exponentiellement.

Période propre :
\[ T_0 = 2\pi\sqrt{LC} = 2\pi\sqrt{\num{0,10}\times\num{1,0e-6}} = 2\pi\times \num{3,16e-4} \approx \SI{2,0}{ms}. \]

Pseudo-période exacte :
\[ T = \frac{2\pi}{\sqrt{\dfrac{1}{LC} - \dfrac{R^2}{4L^2}}} = \frac{2\pi}{\sqrt{\num{1,0e7} - \dfrac{2500}{\num{0,04}}}} = \frac{2\pi}{\sqrt{\num{1,0e7} - \num{6,25e4}}} \approx \frac{2\pi}{\num{3,15e3}} \approx \SI{2,0}{ms}. \]

\textbf{Conclusion} : l'amortissement étant faible ($R \ll R_C$), la pseudo-période est pratiquement égale à la période propre : $T \approx T_0 \approx \SI{2,0}{ms}$. Seule l'amplitude des oscillations trahit la présence du résistor.
\end{exoresolu}

\courssub{Analogies électromécaniques et synthèse}

\begin{methode}[Utiliser le tableau des analogies électromécaniques]
\begin{enumerate}
\item \textbf{Identifier les grandeurs correspondantes} : la charge $q$ joue le rôle de l'élongation $x$, l'intensité $i = \dot{q}$ celui de la vitesse $v = \dot{x}$.
\item \textbf{Associer les éléments des dipôles aux éléments mécaniques} :
\begin{center}
\begin{tabularx}{\linewidth}{|l|X|X|}
\hline
\rowcolor{popblueL} \textbf{Grandeur} & \textbf{Oscillateur mécanique} & \textbf{Oscillateur électrique} \\
\hline
Position / charge & $x$ (m) & $q$ (C) \\
\hline
Vitesse / intensité & $v = \dot{x}$ (m/s) & $i = \dot{q}$ (A) \\
\hline
Masse / inductance (inertie) & $m$ (kg) & $L$ (H) \\
\hline
Raideur / inverse de capacité & $k$ (N/m) & $1/C$ (F$^{-1}$) \\
\hline
Frottement / résistance & $f$ (N.s/m) & $R$ ($\Omega$) \\
\hline
Énergie cinétique / magnétique & $\frac{1}{2}mv^2$ & $\frac{1}{2}Li^2$ \\
\hline
Énergie potentielle / électrique & $\frac{1}{2}kx^2$ & $\dfrac{q^2}{2C}$ \\
\hline
Période propre & $T_0 = 2\pi\sqrt{m/k}$ & $T_0 = 2\pi\sqrt{LC}$ \\
\hline
\end{tabularx}
\end{center}
\item \textbf{Traduire l'équation différentielle} : à $m\ddot{x} + f\dot{x} + kx = 0$ correspond $L\ddot{q} + R\dot{q} + \dfrac{q}{C} = 0$.
\item \textbf{Utiliser l'analogie pour raisonner vite} : la bobine « s'oppose aux variations de courant » comme la masse s'oppose aux variations de vitesse (inertie) ; le condensateur « emmagasine » comme le ressort.
\end{enumerate}
\end{methode}

\begin{exoresolu}[Du ressort au circuit LC]
Un solide de masse $m = \SI{100}{g}$ accroché à un ressort de raideur $k = \SI{40}{N.m^{-1}}$ oscille sans frottement sur un plan horizontal.
\begin{enumerate}
\item Calculer la période propre $T_0$ de cet oscillateur mécanique.
\item Déterminer l'inductance $L$ de la bobine à associer à un condensateur de capacité $C = \SI{10}{\micro F}$ pour obtenir un circuit LC de même période propre.
\item Justifier la correspondance entre les énergies des deux oscillateurs.
\end{enumerate}
\tcblower
\textbf{1. Période propre mécanique.}
\[ T_0 = 2\pi\sqrt{\frac{m}{k}} = 2\pi\sqrt{\frac{\num{0,100}}{40}} = 2\pi\sqrt{\num{2,5e-3}} = 2\pi\times \num{5,0e-2} \approx \SI{0,31}{s}. \]

\textbf{2. Inductance équivalente.} On impose $2\pi\sqrt{LC} = T_0$ :
\[ L = \frac{T_0^2}{4\pi^2 C} = \frac{(\num{0,314})^2}{4\pi^2\times \num{1,0e-5}} = \frac{\num{9,87e-2}}{\num{3,95e-4}} \approx \SI{2,5e2}{H}. \]
Soit $L \approx \SI{250}{H}$ (valeur très élevée, ce qui montre que les oscillations électriques « lentes » exigent de grosses inductances).

\textbf{3. Correspondance énergétique.} Par analogie $m \leftrightarrow L$, $k \leftrightarrow 1/C$, $x \leftrightarrow q$, $v \leftrightarrow i$ :
\begin{itemize}
\item énergie cinétique $\frac{1}{2}mv^2 \leftrightarrow$ énergie magnétique $\frac{1}{2}Li^2$ ;
\item énergie potentielle élastique $\frac{1}{2}kx^2 \leftrightarrow$ énergie électrique $\dfrac{q^2}{2C}$.
\end{itemize}
Dans les deux cas, l'énergie totale se conserve (pas de frottement, pas de résistance) et s'échange alternativement entre les deux réservoirs : ressort/masse d'un côté, condensateur/bobine de l'autre.

\textbf{Conclusion} : les deux oscillateurs sont régis par la même équation différentielle $\ddot{y} + \omega_0^2 y = 0$ ; l'analogie permet de transposer tout résultat de mécanique vers l'électricité, et réciproquement.
\end{exoresolu}

\begin{attention}[Les 5 erreurs qui coûtent des points au Bac]
\begin{enumerate}
\item \textbf{Oublier de convertir en unités SI.} Avec $C = \SI{10}{\micro F}$ non converti, $T_0 = 2\pi\sqrt{LC}$ donne un résultat faux d'un facteur $10^3$. Écris systématiquement $C = \num{1,0e-5}\ \text{F}$, $L = \num{5,0e-2}\ \text{H}$ avant l'application numérique.
\item \textbf{Confondre $i = \dfrac{dq}{dt}$ et $i = -\dfrac{dq}{dt}$.} Le signe dépend du sens choisi pour $i$ par rapport à l'armature portant $q$. Annonce clairement tes conventions sur un schéma fléché : sans cela, l'équation différentielle peut apparaître avec un signe faux et la copie est pénalisée.
\item \textbf{Croire que $T_0$ dépend de la tension de charge $U_0$.} La période propre $T_0 = 2\pi\sqrt{LC}$ ne dépend que de $L$ et $C$ : augmenter $U_0$ augmente l'amplitude et l'énergie, jamais la période. C'est le piège le plus fréquent des questions de cours.
\item \textbf{Écrire que l'énergie se conserve dans un circuit RLC.} Dans le RLC, $\dfrac{dE}{dt} = -Ri^2 < 0$ : l'énergie décroît par effet Joule. La conservation $E = \text{cte}$ n'est valable que pour le LC \emph{idéal} ($R = 0$). Précise toujours la condition d'application.
\item \textbf{Confondre pseudo-période et période propre, ou les régimes d'amortissement.} En régime pseudo-périodique, $T > T_0$ (mais $T \approx T_0$ si l'amortissement est faible). Retiens le critère sur $R_C = 2\sqrt{L/C}$ : $R < R_C$ pseudo-périodique, $R = R_C$ critique, $R > R_C$ apériodique — et non l'inverse.
\end{enumerate}
\end{attention}

\newpage

\coursec{Exercices}

\courssub{Je vérifie mes connaissances}

\begin{exercice}[QCM : la bobine et le circuit LC]
Pour chaque affirmation, choisir la (ou les) bonne(s) réponse(s) en justifiant brièvement.

\textbf{1.} La tension aux bornes d'une bobine d'inductance $L$ et de résistance interne $r$, parcourue par un courant $i$ en convention récepteur, s'écrit :
\begin{enumerate}[label=\alph*)]
\item $u = ri - L\dfrac{di}{dt}$
\item $u = ri + L\dfrac{di}{dt}$
\item $u = \dfrac{1}{L}\displaystyle\int i\,dt$
\end{enumerate}

\textbf{2.} Dans un circuit LC idéal, la période propre des oscillations est :
\begin{enumerate}[label=\alph*)]
\item $T_0 = 2\pi\sqrt{\dfrac{L}{C}}$
\item $T_0 = 2\pi\sqrt{LC}$
\item $T_0 = \dfrac{2\pi}{\sqrt{LC}}$
\end{enumerate}

\textbf{3.} L'énergie magnétique emmagasinée dans une bobine parcourue par un courant d'intensité $i$ vaut :
\begin{enumerate}[label=\alph*)]
\item $\dfrac{1}{2}Li^2$
\item $\dfrac{1}{2}L^2i$
\item $\dfrac{Li^2}{2C}$
\end{enumerate}

\textbf{4.} Dans un circuit LC idéal, quand la charge du condensateur est maximale en valeur absolue :
\begin{enumerate}[label=\alph*)]
\item l'intensité est maximale ;
\item l'intensité est nulle ;
\item l'énergie est entièrement sous forme magnétique.
\end{enumerate}
\end{exercice}

\begin{exercice}[Vrai ou faux, avec justification]
Répondre par \textbf{vrai} ou \textbf{faux} en justifiant chaque réponse par une phrase ou un calcul.

\textbf{1.} Dans un circuit LC idéal, l'énergie électromagnétique totale se conserve au cours du temps.

\textbf{2.} Si l'on double l'inductance $L$ d'un circuit LC, sa période propre double.

\textbf{3.} Dans un circuit RLC série, plus la résistance $R$ est grande, plus les oscillations sont amorties.

\textbf{4.} Le régime critique correspond à la plus grande valeur de $R$ pour laquelle on observe encore des oscillations.

\textbf{5.} Dans un circuit LC, $q(t)$ et $i(t)$ oscillent en quadrature de phase.
\end{exercice}

\begin{exercice}[Définitions et expressions]
\textbf{1.} Donner la définition de l'inductance $L$ d'une bobine et préciser son unité SI ainsi que le nom de cette unité.

\textbf{2.} Écrire l'équation différentielle vérifiée par la charge $q(t)$ du condensateur dans un circuit LC idéal, puis donner l'expression générale de sa solution $q(t)$.

\textbf{3.} Définir la \cle{pulsation propre} $\omega_0$ d'un circuit LC et donner les relations liant $\omega_0$, $T_0$ et la fréquence propre $f_0$.

\textbf{4.} Écrire l'équation différentielle des oscillations amorties de la charge $q(t)$ dans un circuit RLC série, en faisant apparaître $R$, $L$ et $C$.
\end{exercice}

\begin{exercice}[Calculs directs]
On réalise un circuit LC avec une bobine idéale d'inductance $L = \SI{0,20}{H}$ et un condensateur de capacité $C = \SI{10}{\micro F}$. Le condensateur est initialement chargé sous la tension $U_0 = \SI{12}{V}$.

\textbf{1.} Calculer la pulsation propre $\omega_0$, la période propre $T_0$ et la fréquence propre $f_0$ des oscillations.

\textbf{2.} Calculer la charge maximale $Q_m$ du condensateur et l'intensité maximale $I_m$ du courant.

\textbf{3.} Calculer l'énergie électromagnétique totale $E$ du circuit.
\end{exercice}

\courssub{Je m'entraîne}

\begin{exercice}[Oscillations libres dans un circuit LC]
Un condensateur de capacité $C = \SI{2,2}{\micro F}$ est chargé sous une tension $U = \SI{9,0}{V}$, puis connecté à l'instant $t = 0$ aux bornes d'une bobine idéale d'inductance $L = \SI{25}{mH}$. On note $q(t)$ la charge de l'armature reliée à la bobine.

\textbf{1.} Établir l'équation différentielle vérifiée par $q(t)$.

\textbf{2.} Déterminer les expressions numériques de $q(t)$ et de $i(t)$ en précisant les conditions initiales utilisées.

\textbf{3.} Calculer la période propre $T_0$ et l'intensité maximale $I_m$.

\textbf{4.} Calculer l'énergie électrique emmagasinée dans le condensateur à l'instant $t = \dfrac{T_0}{8}$. En déduire l'énergie magnétique de la bobine à cet instant.
\end{exercice}

\begin{exercice}[Dimensionnement d'un circuit oscillant]
Pour un projet de club scientifique au lycée, des élèves veulent réaliser un circuit oscillant de fréquence propre $f_0 = \SI{50}{kHz}$ à l'aide d'une bobine idéale d'inductance $L = \SI{2,0}{mH}$.

\textbf{1.} Déterminer la capacité $C$ du condensateur à associer à cette bobine.

\textbf{2.} Le condensateur est chargé initialement avec une charge $Q_0 = \SI{0,20}{\micro C}$. Calculer l'intensité maximale $I_m$ du courant dans le circuit.

\textbf{3.} Calculer l'énergie électromagnétique totale du circuit.

\textbf{4.} Les élèves souhaitent diviser la fréquence propre par deux. Proposer deux modifications indépendantes du circuit permettant d'atteindre cet objectif, en précisant les nouvelles valeurs.
\end{exercice}

\begin{exercice}[Les trois régimes du circuit RLC]
On réalise un circuit RLC série avec une bobine d'inductance $L = \SI{0,10}{H}$ (de résistance négligeable), un condensateur de capacité $C = \SI{1,0}{\micro F}$ et un conducteur ohmique de résistance $R$ réglable. On rappelle que la résistance critique du circuit vaut $R_C = 2\sqrt{\dfrac{L}{C}}$.

\textbf{1.} Calculer la valeur de $R_C$.

\textbf{2.} Préciser, en le justifiant, le régime observé pour chacune des valeurs suivantes : $R_1 = \SI{20}{\Omega}$ ; $R_2 = \SI{632}{\Omega}$ ; $R_3 = \SI{2,0}{k\Omega}$.

\textbf{3.} Esquisser, sur trois graphes $q(t)$ différents, l'allure de la charge du condensateur (initialement chargée) dans chacun des trois cas. On indiquera clairement les différences entre les allures.

\textbf{4.} Dans le cas $R_1 = \SI{20}{\Omega}$, calculer la pseudo-période $T$ en l'approchant par la période propre $T_0$. Justifier cette approximation.
\end{exercice}

\begin{exercice}[Analogie électromécanique]
On considère un oscillateur mécanique constitué d'un solide de masse $m$ accroché à un ressort de raideur $k$, soumis à une force de frottement fluide $\vec{f} = -\alpha\,\vec{v}$, et un circuit RLC série.

\textbf{1.} Recopier et compléter le tableau de correspondance suivant :

\begin{center}
\begin{tabularx}{\linewidth}{|l|X|X|}
\hline
\rowcolor{popblueL} \textbf{Grandeur} & \textbf{Oscillateur mécanique} & \textbf{Circuit RLC} \\
\hline
État du système & élongation $x$ & \ldots \\
\hline
Vitesse de variation & vitesse $v = \dot{x}$ & \ldots \\
\hline
Inertie & masse $m$ & \ldots \\
\hline
Rappel & raideur $k$ & \ldots \\
\hline
Amortissement & coefficient $\alpha$ & \ldots \\
\hline
Énergie potentielle & $\dfrac{1}{2}kx^2$ & \ldots \\
\hline
Énergie cinétique & $\dfrac{1}{2}mv^2$ & \ldots \\
\hline
\end{tabularx}
\end{center}

\textbf{2.} En utilisant la formule de Thomson $T_0 = 2\pi\sqrt{LC}$ et l'analogie du tableau, déduire l'expression de la période propre $T_{0m}$ de l'oscillateur mécanique non amorti en fonction de $m$ et $k$.

\textbf{3.} Vérifier par analyse dimensionnelle que l'expression obtenue est bien homogène à un temps.

\textbf{4.} À quelle grandeur électrique le frottement fluide est-il analogue ? Que devient l'énergie mécanique totale de l'oscillateur amorti au cours du temps, et quel est l'équivalent électrique de ce phénomène ?
\end{exercice}

\courssub{Je me prépare au Bac}

\begin{exercice}[Type Bac : oscillations libres dans un circuit LC]
\textit{(D'après Baccalauréat C, Cameroun)}

Au laboratoire d'un lycée de Douala, un professeur réalise le montage suivant : un condensateur de capacité $C = \SI{4,7}{\micro F}$ est chargé sous une tension continue $U = \SI{6,0}{V}$ à l'aide d'un générateur. À l'instant $t = 0$, on bascule l'interrupteur pour connecter le condensateur chargé aux bornes d'une bobine idéale d'inductance $L = \SI{50}{mH}$. On note $q(t)$ la charge de l'armature du condensateur reliée à la bobine et $i(t)$ l'intensité du courant, orientée en convention récepteur par rapport à $q$.

\textbf{1.} Faire le schéma du circuit après le basculement de l'interrupteur, en représentant les flèches de tension $u_C$ aux bornes du condensateur et $u_L$ aux bornes de la bobine.

\textbf{2.} En appliquant la loi d'additivité des tensions, établir l'équation différentielle vérifiée par la charge $q(t)$. La mettre sous la forme $\ddot{q} + \omega_0^2\,q = 0$ et exprimer $\omega_0$.

\textbf{3.} Calculer numériquement $\omega_0$, la période propre $T_0$ et la fréquence propre $f_0$ des oscillations.

\textbf{4.} La solution de l'équation différentielle s'écrit $q(t) = Q_m\cos(\omega_0 t + \varphi)$. Déterminer $Q_m$ et $\varphi$ à l'aide des conditions initiales, puis écrire les expressions numériques de $q(t)$ et de $i(t)$.

\textbf{5.} Calculer l'intensité maximale $I_m$ du courant dans le circuit.

\textbf{6.} Exprimer l'énergie électromagnétique totale $E$ du circuit en fonction de $q$, $i$, $L$ et $C$, puis calculer sa valeur numérique. Que peut-on dire de son évolution au cours du temps ?

\textbf{7.} Déterminer la valeur de $q$ pour laquelle l'énergie est répartie également entre le condensateur et la bobine.
\end{exercice}

\begin{exercice}[Type Bac : exploitation d'un oscillogramme en régime pseudo-périodique]
\textit{(D'après Baccalauréat C, Cameroun)}

Dans le cadre d'un TP, des élèves de Terminale C étudient la décharge d'un condensateur de capacité $C = \SI{1,0}{\micro F}$, initialement chargé sous $U_0 = \SI{5,0}{V}$, dans un circuit série comportant une bobine d'inductance $L$ inconnue et de résistance interne $r = \SI{15}{\Omega}$, ainsi qu'un conducteur ohmique de résistance $R = \SI{180}{\Omega}$. Ils visualisent la tension $u_C(t)$ aux bornes du condensateur sur un oscilloscope à mémoire et obtiennent l'oscillogramme ci-dessous.

\begin{popfigure}
\begin{center}
\begin{tikzpicture}
\begin{axis}[
  width=0.9\linewidth, height=5.2cm,
  axis lines=middle,
  xlabel={$t$ (ms)}, ylabel={$u_C$ (V)},
  xmin=0, xmax=13, ymin=-4, ymax=6,
  xtick={0,2,4,6,8,10,12},
  ytick={-4,-2,0,2,4,6},
  tick label style={font=\small},
  label style={font=\small},
  samples=300
]
\addplot[popcurve,domain=0:12.8] {5*exp(-0.22*x)*cos(deg(2*pi*x/4.19))};
\draw[dashed,popmuted] (axis cs:4.19,0) -- (axis cs:4.19,2.0);
\draw[dashed,popmuted] (axis cs:8.38,0) -- (axis cs:8.38,0.8);
\end{axis}
\end{tikzpicture}
\end{center}
\legende{Oscillogramme de la tension $u_C(t)$ en régime pseudo-périodique.}
\end{popfigure}

\textbf{1.} Justifier qualitativement que le régime observé est pseudo-périodique.

\textbf{2.} À partir de l'oscillogramme, mesurer la pseudo-période $T$ des oscillations. Expliquer la méthode utilisée et préciser pourquoi il est préférable de mesurer plusieurs pseudo-périodes.

\textbf{3.} En assimilant la pseudo-période $T$ à la période propre $T_0 = 2\pi\sqrt{LC}$, calculer l'inductance $L$ de la bobine. Justifier brièvement cette assimilation compte tenu des valeurs de $R$ et $r$.

\textbf{4.} Établir l'équation différentielle vérifiée par la charge $q(t)$ du condensateur dans ce circuit.

\textbf{5.} Expliquer, par un bilan énergétique détaillé, l'origine de la décroissance de l'amplitude des oscillations. Écrire la relation liant la variation d'énergie électromagnétique $dE$ à l'énergie dissipée pendant la durée $dt$.

\textbf{6.} Calculer l'énergie électromagnétique initiale $E_0$ du circuit, puis l'énergie dissipée par effet Joule entre $t = 0$ et l'instant où l'amplitude de $u_C$ vaut $\SI{2,0}{V}$ (on supposera qu'à cet instant l'intensité est nulle).
\end{exercice}

\begin{exercice}[Type Bac : démonstration énergétique et étude complète d'un circuit RLC]
\textit{(D'après Baccalauréat C, Cameroun)}

\textbf{Partie A : circuit LC idéal.}

Un condensateur de capacité $C$ et une bobine idéale d'inductance $L$ sont associés en série. Le condensateur porte la charge $q(t)$ et le circuit est parcouru par le courant $i(t) = \dfrac{dq}{dt}$.

\textbf{1.} Rappeler les expressions de l'énergie électrique $E_e$ emmagasinée dans le condensateur et de l'énergie magnétique $E_m$ emmagasinée dans la bobine.

\textbf{2.} On définit l'énergie électromagnétique totale $E = \dfrac{q^2}{2C} + \dfrac{1}{2}Li^2$. En admettant que $E$ se conserve, montrer que la condition $\dfrac{dE}{dt} = 0$ conduit à l'équation différentielle $\ddot{q} + \dfrac{q}{LC} = 0$.

\textbf{3.} Réciproquement, montrer que l'équation différentielle du circuit LC entraîne la conservation de $E$.

\textbf{Partie B : circuit RLC réel.}

On associe maintenant en série le condensateur de capacité $C = \SI{0,47}{\micro F}$, la bobine d'inductance $L = \SI{40}{mH}$ et de résistance interne $r = \SI{8,0}{\Omega}$, et un conducteur ohmique de résistance $R$. Le condensateur est initialement chargé sous $U_0 = \SI{10}{V}$.

\textbf{4.} Établir l'équation différentielle vérifiée par $q(t)$.

\textbf{5.} Calculer la résistance critique $R_C$ de ce circuit (on prendra en compte la résistance totale $R + r$). En déduire la valeur maximale de $R$ permettant d'observer un régime pseudo-périodique.

\textbf{6.} On règle $R = \SI{12}{\Omega}$. Calculer la période propre $T_0$ et justifier que la pseudo-période mesurée à l'oscilloscope en est très proche.

\textbf{7.} Calculer l'énergie totale dissipée par effet Joule dans le circuit entre $t = 0$ et l'instant où les oscillations sont totalement amorties.

\textbf{8.} Citer une application technologique des circuits oscillants dans le domaine des télécommunications au Cameroun, et expliquer en une phrase le rôle du circuit LC dans cette application.
\end{exercice}

\newpage

\coursec{Corrigés des exercices}

\coursec{Exercices corrigés — Oscillations électriques libres (circuits LC et RLC)}

\begin{corrige}[Exercice 1 — QCM : la bobine et le circuit LC]
\textbf{1.} Réponse \textbf{b)}. En \cle{convention récepteur}, la tension aux bornes d'une bobine réelle s'écrit $u = ri + L\,\dfrac{di}{dt}$ : le terme $ri$ traduit l'effet Joule dans le fil (résistance interne $r$), le terme $L\,di/dt$ est la tension d'induction (loi de Lenz) qui s'oppose aux variations du courant. La réponse a) correspondrait à une convention générateur ; la réponse c) est la loi du condensateur ($u_C = q/C$), pas de la bobine.

\textbf{2.} Réponse \textbf{b)}. $T_0 = \dfrac{2\pi}{\omega_0} = 2\pi\sqrt{LC}$. \textbf{Vérification dimensionnelle} : $[L] = \text{V}\cdot\text{s}\cdot\text{A}^{-1} = \text{H}$, $[C] = \text{A}\cdot\text{s}\cdot\text{V}^{-1} = \text{F}$, donc $[LC] = \text{s}^2$ et $[\sqrt{LC}] = \text{s}$ : homogène à une période. La réponse a) donnerait $[L/C] = \text{V}^2\cdot\text{A}^{-2}$ (sans dimension de temps) ; la réponse c) donne des $\text{s}^{-1}$ (une pulsation, à $2\pi$ près).

\textbf{3.} Réponse \textbf{a)}. $E_m = \dfrac{1}{2}Li^2$. Unité : $\text{H}\cdot\text{A}^2 = (\text{V}\cdot\text{s}\cdot\text{A}^{-1})\cdot\text{A}^2 = \text{V}\cdot\text{A}\cdot\text{s} = \text{W}\cdot\text{s} = \text{J}$. La réponse b) n'est pas homogène ($\text{H}^2\cdot\text{A}$) ; la réponse c) mélange les formules de la bobine et du condensateur.

\textbf{4.} Réponse \textbf{b)}. Si $|q| = Q_m$, alors $\cos(\omega_0 t + \varphi) = \pm 1$, donc $\sin(\omega_0 t + \varphi) = 0$ et $i = -\omega_0 Q_m \sin(\omega_0 t + \varphi) = 0$. À cet instant $E_m = \frac{1}{2}Li^2 = 0$ : l'énergie est entièrement sous forme \emph{électrique} ($E_e = E_{\text{tot}}$), ce qui rend la réponse c) fausse.
\end{corrige}

\begin{corrige}[Exercice 2 — Vrai ou faux]
\textbf{1.} \textbf{Vrai.} Un circuit LC idéal ne contient aucune résistance : il n'y a pas de dissipation par effet Joule. On montre que $\dfrac{dE}{dt} = i\left(\dfrac{q}{C} + L\ddot{q}\right) = 0$ d'après l'équation différentielle, donc $E = \text{constante}$.

\textbf{2.} \textbf{Faux.} $T_0 = 2\pi\sqrt{LC}$ est proportionnelle à $\sqrt{L}$, pas à $L$. Doubler $L$ multiplie $T_0$ par $\sqrt{2} \approx 1,41$ (augmentation d'environ $41\,\%$, pas $100\,\%$).

\textbf{3.} \textbf{Vrai.} Le coefficient d'amortissement $\sigma = \dfrac{R_{\text{tot}}}{2L}$ croît avec $R$ : l'enveloppe $Q_m e^{-\sigma t}$ décroît d'autant plus vite, et la puissance dissipée $R_{\text{tot}}i^2$ augmente. \textbf{Attention cependant} : au-delà de $R_C$, il n'y a plus d'oscillations du tout (régime apériodique).

\textbf{4.} \textbf{Faux.} C'est l'inverse : le régime critique ($R_{\text{tot}} = R_C = 2\sqrt{L/C}$) est la \emph{plus petite} valeur de la résistance pour laquelle les oscillations \emph{disparaissent}. Pour $R_{\text{tot}} < R_C$ on observe encore des oscillations (pseudo-périodique) ; pour $R_{\text{tot}} \geq R_C$, plus aucune oscillation.

\textbf{5.} \textbf{Vrai.} $i = \dfrac{dq}{dt}$. Si $q = Q_m\cos(\omega_0 t + \varphi)$, alors $i = -\omega_0 Q_m\sin(\omega_0 t + \varphi) = \omega_0 Q_m\cos\!\left(\omega_0 t + \varphi + \dfrac{\pi}{2}\right)$ : le déphasage de $i$ sur $q$ vaut exactement $+\dfrac{\pi}{2}$ (quadrature de phase), quelles que soient les conditions initiales.
\end{corrige}

\begin{corrige}[Exercice 3 — Définitions et expressions]
\textbf{1.} L'\cle{inductance propre} $L$ d'une bobine est le rapport du flux magnétique propre $\Phi$ traversant la bobine à l'intensité $i$ du courant qui le crée : $L = \dfrac{\Phi}{i}$. Unité SI : le \textbf{henry} (symbole H), avec $1\,\text{H} = 1\,\text{Wb}\cdot\text{A}^{-1} = 1\,\text{V}\cdot\text{s}\cdot\text{A}^{-1}$.

\textbf{2.} Équation différentielle : $\ddot{q} + \dfrac{1}{LC}\,q = 0$. Solution générale : $q(t) = Q_m\cos(\omega_0 t + \varphi)$, avec $\omega_0 = \dfrac{1}{\sqrt{LC}}$, $Q_m$ et $\varphi$ fixés par les conditions initiales.

\textbf{3.} La \cle{pulsation propre} est $\omega_0 = \dfrac{1}{\sqrt{LC}}$ (en $\text{rad}\cdot\text{s}^{-1}$). Relations : $T_0 = \dfrac{2\pi}{\omega_0} = 2\pi\sqrt{LC}$ (période propre, en s) et $f_0 = \dfrac{1}{T_0} = \dfrac{\omega_0}{2\pi}$ (fréquence propre, en Hz).

\textbf{4.} Pour un circuit RLC série (bobine de résistance interne $r$ incluse dans $R_{\text{tot}}$) :
\[
L\,\ddot{q} + R_{\text{tot}}\,\dot{q} + \frac{q}{C} = 0 \qquad\text{avec } R_{\text{tot}} = R + r.
\]
Forme réduite : $\ddot{q} + 2\sigma\,\dot{q} + \omega_0^2\,q = 0$ avec $\sigma = \dfrac{R_{\text{tot}}}{2L}$ (coefficient d'amortissement, en $\text{s}^{-1}$).
\end{corrige}

\begin{corrige}[Exercice 4 — Calculs directs]
Données : $L = \SI{0,20}{H}$, $C = \SI{10}{\micro\farad} = \SI{1,0e-5}{F}$, $U_0 = \SI{12}{V}$.

\textbf{1.} Pulsation propre :
\[
\omega_0 = \frac{1}{\sqrt{LC}} = \frac{1}{\sqrt{0,20 \times 1,0\times10^{-5}}} = \frac{1}{\sqrt{2,0\times10^{-6}}} = \frac{1}{1,414\times10^{-3}} \approx \SI{707}{rad/s}.
\]
Période propre : $T_0 = \dfrac{2\pi}{\omega_0} = \dfrac{6,283}{707} \approx \SI{8,89e-3}{s} = \SI{8,89}{ms}$.
Fréquence propre : $f_0 = \dfrac{1}{T_0} \approx \dfrac{1}{8,89\times10^{-3}} \approx \SI{112}{Hz}$.

\textbf{2.} Charge maximale : $Q_m = C\,U_0 = \SI{1,0e-5}{F} \times \SI{12}{V} = \SI{1,2e-4}{C} = \SI{120}{\micro\coulomb}$.
Intensité maximale : $I_m = \omega_0 Q_m = 707 \times 1,2\times10^{-4} \approx \SI{8,48e-2}{A} = \SI{84,8}{mA}$.

\textbf{3.} Énergie totale (conservée) :
\[
E = \frac{1}{2}C\,U_0^2 = 0,5 \times 1,0\times10^{-5} \times 144 = \SI{7,2e-4}{J} = \SI{720}{\micro\joule}.
\]
Vérification croisée : $E = \dfrac{1}{2}L I_m^2 = 0,5 \times 0,20 \times (8,48\times10^{-2})^2 \approx 0,10 \times 7,19\times10^{-3} \approx \SI{7,19e-4}{J}$. Cohérent (arrondis).
\end{corrige}

\begin{corrige}[Exercice 5 — Oscillations libres dans un circuit LC]
Données : $L = \SI{25}{mH} = \SI{2,5e-2}{H}$, $C = \SI{2,2}{\micro\farad} = \SI{2,2e-6}{F}$, $U = \SI{9,0}{V}$.

\textbf{1.} Après fermeture de l'interrupteur, la maille contient le condensateur et la bobine idéale. Loi d'additivité des tensions (convention récepteur pour les deux dipôles) : $u_C + u_L = 0$ avec $u_C = \dfrac{q}{C}$ et $u_L = L\,\dfrac{di}{dt} = L\,\ddot{q}$. D'où :
\[
\frac{q}{C} + L\,\ddot{q} = 0 \quad\Longleftrightarrow\quad \ddot{q} + \frac{1}{LC}\,q = 0 .
\]
Pulsation propre : $\omega_0 = \dfrac{1}{\sqrt{LC}} = \dfrac{1}{\sqrt{2,5\times10^{-2} \times 2,2\times10^{-6}}} = \dfrac{1}{\sqrt{5,5\times10^{-8}}} = \dfrac{1}{2,345\times10^{-4}} \approx \SI{4,26e3}{rad/s}$.

\textbf{2.} Conditions initiales : à $t = 0$, le condensateur est chargé sous $U = \SI{9,0}{V}$, donc $q(0) = Q_m = C\,U = 2,2\times10^{-6} \times 9,0 = \SI{19,8}{\micro\coulomb}$ ; le courant était nul avant la fermeture et ne peut pas subir de discontinuité dans la bobine (continuité du flux), donc $i(0) = 0$.
$q(0) = Q_m\cos\varphi = Q_m \Rightarrow \varphi = 0$. D'où :
\[
q(t) = \SI{19,8}{\micro\coulomb}\,\cos(4,26\times10^{3}\,t), \qquad
i(t) = \dot{q}(t) = -\SI{84,4}{mA}\,\sin(4,26\times10^{3}\,t).
\]
($\omega_0 Q_m = 4,26\times10^{3} \times 19,8\times10^{-6} \approx 8,44\times10^{-2}\,\text{A}$.)

\textbf{3.} $T_0 = \dfrac{2\pi}{\omega_0} = \dfrac{6,283}{4,26\times10^{3}} \approx \SI{1,47e-3}{s} = \SI{1,47}{ms}$. \quad $I_m = \omega_0 Q_m \approx \SI{84,4}{mA}$.

\textbf{4.} À $t = \dfrac{T_0}{8}$ : $\omega_0 t = \dfrac{2\pi}{T_0}\times\dfrac{T_0}{8} = \dfrac{\pi}{4}$, donc $q = Q_m\cos\dfrac{\pi}{4} = \dfrac{Q_m}{\sqrt{2}}$.
\[
E_e = \frac{q^2}{2C} = \frac{Q_m^2}{4C} = \frac{1}{2}\times\frac{Q_m^2}{2C} = \frac{E}{2}.
\]
Or $E = \dfrac{1}{2}C\,U^2 = 0,5 \times 2,2\times10^{-6} \times 81 = \SI{89,1}{\micro\joule}$.
Donc $E_e = \SI{44,6}{\micro\joule}$ et, par conservation de l'énergie, $E_m = E - E_e = \SI{44,6}{\micro\joule}$ : l'énergie est équitablement répartie entre le condensateur et la bobine.
\end{corrige}

\begin{corrige}[Exercice 6 — Dimensionnement d'un circuit oscillant]
Données : $f_0 = \SI{50}{kHz} = \SI{5,0e4}{Hz}$, $L = \SI{2,0}{mH} = \SI{2,0e-3}{H}$, $Q_0 = \SI{0,20}{\micro\coulomb} = \SI{2,0e-7}{C}$.

\textbf{1.} De $f_0 = \dfrac{1}{2\pi\sqrt{LC}}$, on tire $\sqrt{LC} = \dfrac{1}{2\pi f_0}$, puis :
\[
C = \frac{1}{(2\pi f_0)^2 L} = \frac{1}{(2\pi \times 5,0\times10^{4})^2 \times 2,0\times10^{-3}} = \frac{1}{(3,142\times10^{5})^2 \times 2,0\times10^{-3}} = \frac{1}{1,974\times10^{8}} \approx \SI{5,07e-9}{F} = \SI{5,07}{nF}.
\]

\textbf{2.} $\omega_0 = 2\pi f_0 = 2\pi \times 5,0\times10^{4} \approx \SI{3,14e5}{rad/s}$.
\[
I_m = \omega_0 Q_0 = 3,14\times10^{5} \times 2,0\times10^{-7} = \SI{6,28e-2}{A} = \SI{62,8}{mA}.
\]

\textbf{3.} $E = \dfrac{Q_0^2}{2C} = \dfrac{(2,0\times10^{-7})^2}{2 \times 5,07\times10^{-9}} = \dfrac{4,0\times10^{-14}}{1,014\times10^{-8}} \approx \SI{3,94e-6}{J} = \SI{3,94}{\micro\joule}$.
Vérification : $E = \frac{1}{2}LI_m^2 = 0,5 \times 2,0\times10^{-3} \times (6,28\times10^{-2})^2 = 1,0\times10^{-3} \times 3,94\times10^{-3} = \SI{3,94e-6}{J}$. Les deux calculs concordent parfaitement.

\textbf{4.} Diviser $f_0$ par 2 revient à multiplier $T_0 = 2\pi\sqrt{LC}$ par 2, donc à multiplier le produit $LC$ par 4. Deux solutions indépendantes :
\begin{itemize}
\item \textbf{Conserver $L = \SI{2,0}{mH}$} et prendre $C' = 4C \approx \SI{20,3}{nF}$ ;
\item \textbf{Conserver $C \approx \SI{5,07}{nF}$} et prendre $L' = 4L = \SI{8,0}{mH}$.
\end{itemize}
\end{corrige}

\begin{corrige}[Exercice 7 — Les trois régimes du circuit RLC]
Données : $L = \SI{0,10}{H}$, $C = \SI{1,0}{\micro\farad} = \SI{1,0e-6}{F}$.

\textbf{1.} Résistance critique :
\[
R_C = 2\sqrt{\dfrac{L}{C}} = 2\sqrt{\dfrac{0,10}{1,0\times10^{-6}}} = 2\sqrt{1,0\times10^{5}} = 2 \times 316,2 \approx \SI{632}{\Omega}.
\]

\textbf{2.}
\begin{itemize}
\item $R_1 = \SI{20}{\Omega} < R_C$ : \textbf{régime pseudo-périodique} ($\sigma < \omega_0$, racines complexes conjuguées) : oscillations dont l'amplitude décroît exponentiellement.
\item $R_2 = \SI{632}{\Omega} \approx R_C$ : \textbf{régime critique} ($\sigma = \omega_0$, racine double) : retour à l'équilibre le plus rapide possible sans oscillation.
\item $R_3 = \SI{2,0}{k\Omega} > R_C$ : \textbf{régime apériodique} ($\sigma > \omega_0$, deux racines réelles négatives) : retour lent vers l'équilibre, sans oscillation.
\end{itemize}

\textbf{3.} Allures attendues pour $q(t)$ (condensateur initialement chargé, $q(0) = Q_0 > 0$, $i(0) = 0$) :
\begin{itemize}
\item \textbf{Pseudo-périodique :} sinusoïde amortie partant de $Q_0$ avec tangente horizontale, coupant l'axe des temps de nombreuses fois, comprise entre les enveloppes $\pm Q_0 e^{-\sigma t}$.
\item \textbf{Critique :} décroissance rapide vers 0, sans changement de signe (ou un seul passage par zéro), tangente horizontale à l'origine.
\item \textbf{Apériodique :} décroissance monotone vers 0, plus lente que le régime critique aux temps longs, jamais de passage par zéro.
\end{itemize}

\textbf{4.} $T_0 = 2\pi\sqrt{LC} = 2\pi\sqrt{0,10 \times 1,0\times10^{-6}} = 2\pi \times 3,162\times10^{-4} \approx \SI{1,99e-3}{s} = \SI{1,99}{ms}$.
Vérifions l'approximation $T \approx T_0$ pour $R_1 = \SI{20}{\Omega}$ :
$\sigma = \dfrac{R_1}{2L} = \dfrac{20}{0,20} = \SI{100}{s^{-1}}$ et $\omega_0 = \dfrac{1}{\sqrt{LC}} \approx \SI{3162}{rad/s}$.
\[
\omega = \sqrt{\omega_0^2 - \sigma^2} = \sqrt{3162^2 - 100^2} = \sqrt{9,998\times10^{6} - 1,0\times10^{4}} \approx \SI{3161}{rad/s},
\]
d'où $T = \dfrac{2\pi}{\omega} \approx \SI{1,988}{ms} \approx T_0$ (écart relatif $< 0,1\,\%$).
\textbf{Justification :} $\sigma = 100 \ll \omega_0 = 3162$ (équivalent à $R_1 \ll R_C$) : l'amortissement est faible, donc $\omega \approx \omega_0$ et $T \approx T_0$.
\end{corrige}

\begin{corrige}[Exercice 8 — Analogie électromécanique]
\textbf{1.} Tableau complété :
\begin{center}
\begin{tabularx}{\linewidth}{|l|X|X|}
\hline
\rowcolor{popblueL} \textbf{Grandeur} & \textbf{Oscillateur mécanique} & \textbf{Circuit RLC série} \\
\hline
État du système & élongation $x$ & charge $q$ \\
\hline
Vitesse de variation & vitesse $v = \dot{x}$ & intensité $i = \dot{q}$ \\
\hline
Inertie & masse $m$ & inductance $L$ \\
\hline
Rappel & raideur $k$ & inverse de la capacité $1/C$ \\
\hline
Amortissement & coefficient $\alpha$ & résistance totale $R_{\text{tot}}$ \\
\hline
Énergie potentielle & $\frac{1}{2}kx^2$ & $\dfrac{q^2}{2C}$ \\
\hline
Énergie cinétique & $\frac{1}{2}mv^2$ & $\frac{1}{2}Li^2$ \\
\hline
\end{tabularx}
\end{center}

\textbf{2.} Les correspondances $L \leftrightarrow m$ et $\dfrac{1}{C} \leftrightarrow k$ donnent $LC \leftrightarrow \dfrac{m}{k}$. En transposant $T_0 = 2\pi\sqrt{LC}$ :
\[
T_{0m} = 2\pi\sqrt{\frac{m}{k}}.
\]

\textbf{3.} Analyse dimensionnelle : $[m] = \text{kg}$ ; $[k] = \text{N}\cdot\text{m}^{-1} = \text{kg}\cdot\text{m}\cdot\text{s}^{-2}\cdot\text{m}^{-1} = \text{kg}\cdot\text{s}^{-2}$.
\[
\left[\frac{m}{k}\right] = \frac{\text{kg}}{\text{kg}\cdot\text{s}^{-2}} = \text{s}^2
\quad\Rightarrow\quad
\left[\sqrt{\frac{m}{k}}\right] = \text{s}.
\]
L'expression est bien homogène à un temps.

\textbf{4.} Le frottement fluide de coefficient $\alpha$ est analogue à la \textbf{résistance électrique} $R_{\text{tot}} = R + r$. L'énergie mécanique totale de l'oscillateur amorti \textbf{décroît au cours du temps}, dissipée en chaleur par les frottements (puissance $\alpha v^2$). L'équivalent électrique est la \textbf{dissipation de l'énergie électromagnétique par effet Joule} dans la résistance : $\dfrac{dE}{dt} = -R_{\text{tot}}\,i^2$, d'où la décroissance de l'amplitude des oscillations.
\end{corrige}

\begin{corrige}[Exercice 9 — Type Bac : oscillations libres dans un circuit LC]
Données : $L = \SI{50}{mH}$, $C = \SI{4,7}{\micro\farad}$, $U = \SI{6,0}{V}$.

\textbf{1.} Schéma : une maille unique contenant le condensateur et la bobine. La flèche $u_C$ est orientée de l'armature chargée positivement vers l'autre ; le courant $i$ entre par la borne $+$ du condensateur (convention récepteur) et la flèche $u_L$ aux bornes de la bobine est dans le même sens que $u_C$ le long de la maille (convention récepteur pour la bobine également).

\begin{popfigure}
\begin{center}
\begin{circuitikz}[european]
\draw (0,0) to[C, l_={$C$}, v^=$u_C$] (0,2.5)
      to[L, l_={$L$}, v^=$u_L$] (3.5,2.5)
      to[short, i={$i$}] (3.5,0)
      to[short] (0,0);
\end{circuitikz}
\end{center}
\legende{Circuit LC après basculement de l'interrupteur (convention récepteur pour les deux dipôles).}
\end{popfigure}

\textbf{2.} Loi d'additivité des tensions : $u_C + u_L = 0$, avec $u_C = \dfrac{q}{C}$ et $u_L = L\,\dfrac{di}{dt} = L\,\ddot{q}$ (car $i = \dot{q}$ en convention récepteur). D'où :
\[
\frac{q}{C} + L\,\ddot{q} = 0
\quad\Longleftrightarrow\quad
\ddot{q} + \frac{1}{LC}\,q = 0,
\]
soit $\ddot{q} + \omega_0^2 q = 0$ avec $\boxed{\omega_0 = \dfrac{1}{\sqrt{LC}}}$.

\textbf{3.} $\omega_0 = \dfrac{1}{\sqrt{50\times10^{-3} \times 4,7\times10^{-6}}} = \dfrac{1}{\sqrt{2,35\times10^{-7}}} = \dfrac{1}{4,847\times10^{-4}} \approx \SI{2,06e3}{rad/s}$.
$T_0 = \dfrac{2\pi}{\omega_0} \approx \dfrac{6,283}{2063} \approx \SI{3,05e-3}{s} = \SI{3,05}{ms}$.
$f_0 = \dfrac{1}{T_0} \approx \SI{328}{Hz}$.

\textbf{4.} Conditions initiales : $q(0) = C\,U = 4,7\times10^{-6} \times 6,0 = \SI{28,2}{\micro\coulomb}$ et $i(0) = 0$ (continuité du courant dans la bobine).
$q(0) = Q_m\cos\varphi = Q_m$ et $i(0) = -\omega_0 Q_m\sin\varphi = 0$ donnent $\varphi = 0$ et $Q_m = \SI{28,2}{\micro\coulomb}$.
\[
q(t) = \SI{28,2}{\micro\coulomb}\,\cos(2,06\times10^{3}\,t), \qquad
i(t) = -\SI{58,2}{mA}\,\sin(2,06\times10^{3}\,t).
\]

\textbf{5.} $I_m = \omega_0 Q_m = 2,06\times10^{3} \times 28,2\times10^{-6} \approx \SI{58,2}{mA}$.

\textbf{6.} $E = \dfrac{q^2}{2C} + \dfrac{1}{2}Li^2$. Numériquement :
\[
E = \frac{1}{2}C\,U^2 = 0,5 \times 4,7\times10^{-6} \times 36 = \SI{8,46e-5}{J} = \SI{84,6}{\micro\joule}.
\]
Le circuit est idéal (aucune résistance) : pas d'effet Joule, donc \textbf{l'énergie électromagnétique totale se conserve} au cours du temps ; elle se transfère alternativement du condensateur vers la bobine et réciproquement.

\textbf{7.} ÉQUIPARTITION : $E_e = E_m = \dfrac{E}{2}$, soit $\dfrac{q^2}{2C} = \dfrac{1}{2}\times\dfrac{Q_m^2}{2C}$, d'où $q^2 = \dfrac{Q_m^2}{2}$ et
\[
|q| = \frac{Q_m}{\sqrt{2}} = \frac{28,2}{\sqrt{2}} \approx \SI{20,0}{\micro\coulomb}.
\]

\textbf{Barème indicatif (sur 10) :} schéma avec conventions (1) ; équation différentielle établie (2) ; $\omega_0$, $T_0$, $f_0$ (1,5) ; conditions initiales et expressions de $q(t)$, $i(t)$ (2) ; $I_m$ (0,5) ; énergie et conservation (1,5) ; équipartition (1,5).

\textbf{Points de vigilance :} citer explicitement la \emph{convention récepteur} et la loi d'additivité des tensions ; justifier $i(0) = 0$ par la continuité du courant dans la bobine ; ne pas oublier le signe moins dans $i(t) = -\omega_0 Q_m\sin(\omega_0 t)$ ; accompagner chaque résultat de son unité SI.
\end{corrige}

\begin{corrige}[Exercice 10 — Type Bac : exploitation d'un oscillogramme en régime pseudo-périodique]
Données : $C = \SI{1,0}{\micro\farad}$, $R = \SI{180}{\Omega}$, $r = \SI{15}{\Omega}$, $U_0 = \SI{5,0}{V}$. Oscillogramme de $u_C(t)$ fourni (non reproduit ici).

\textbf{1.} La tension $u_C(t)$ change de signe périodiquement (oscillations) tout en voyant son amplitude décroître au cours du temps selon une enveloppe exponentielle. C'est la signature du \textbf{régime pseudo-périodique} : $R_{\text{tot}} = R + r < R_C$.

\textbf{2.} On mesure la durée séparant deux maxima consécutifs (ou deux passages par zéro dans le même sens) : le premier maximum après $t = 0$ se situe vers $t_1 \approx \SI{4,2}{ms}$ et le second vers $t_2 \approx \SI{8,4}{ms}$, d'où
\[
T = t_2 - t_1 \approx \SI{4,2}{ms}.
\]
Mesurer la durée de \emph{plusieurs} pseudo-périodes (par exemple $2T$ ou $3T$) puis diviser par leur nombre réduit l'incertitude relative de lecture sur l'axe des temps (l'erreur de pointé est partagée entre toutes les périodes).

\textbf{3.} En assimilant $T$ à $T_0 = 2\pi\sqrt{LC}$ :
\[
L = \frac{T^2}{4\pi^2 C} = \frac{(4,19\times10^{-3})^2}{4\pi^2 \times 1,0\times10^{-6}} = \frac{1,756\times10^{-5}}{3,948\times10^{-5}} \approx \SI{0,445}{H} = \SI{445}{mH}.
\]
\textit{Justification de l'assimilation :} $R_{\text{tot}} = R + r = 180 + 15 = \SI{195}{\Omega}$, et $R_C = 2\sqrt{L/C} \approx 2\sqrt{0,445/10^{-6}} \approx \SI{1,33}{k\Omega}$. Comme $R_{\text{tot}} < R_C$, le régime est bien pseudo-périodique ; de plus $\sigma = \dfrac{R_{\text{tot}}}{2L} \approx \SI{219}{s^{-1}}$ est petit devant $\omega_0 = \dfrac{2\pi}{T} \approx \SI{1500}{rad/s}$, donc $\omega = \sqrt{\omega_0^2 - \sigma^2} \approx 0,99\,\omega_0$ et $T \approx T_0$ à $1\,\%$ près : l'approximation est justifiée.

\textbf{4.} Loi des mailles en convention récepteur : $u_C + u_R + u_{\text{bobine}} = 0$, soit $\dfrac{q}{C} + (R + r)\,i + L\,\dfrac{di}{dt} = 0$. Avec $i = \dot{q}$ :
\[
\boxed{L\,\ddot{q} + (R + r)\,\dot{q} + \frac{q}{C} = 0}.
\]

\textbf{5.} L'énergie électromagnétique $E = \dfrac{q^2}{2C} + \dfrac{1}{2}Li^2$ n'est pas conservée : la résistance totale $R_{\text{tot}} = R + r$ dissipe de l'énergie par \cle{effet Joule}. En dérivant et en utilisant l'équation différentielle :
\[
\frac{dE}{dt} = i\left(\frac{q}{C} + L\ddot{q}\right) = -R_{\text{tot}}\,i^2 \leq 0,
\]
soit, pendant la durée $dt$ : $\boxed{dE = -R_{\text{tot}}\,i^2\,dt}$. L'énergie dissipée par effet Joule est l'opposée de la variation d'énergie électromagnétique : c'est cette perte continue qui provoque la décroissance de l'amplitude des oscillations.

\textbf{6.} Énergie initiale : $E_0 = \dfrac{1}{2}C\,U_0^2 = 0,5 \times 1,0\times10^{-6} \times 25 = \SI{12,5}{\micro\joule}$.
À l'instant où l'amplitude de $u_C$ vaut $\SI{2,0}{V}$ avec $i = 0$ (extremum de charge), toute l'énergie restante est électrique :
\[
E = \frac{1}{2}C\,u_C^2 = 0,5 \times 1,0\times10^{-6} \times 4,0 = \SI{2,0}{\micro\joule}.
\]
Énergie dissipée par effet Joule entre les deux instants :
\[
W_{\text{Joule}} = E_0 - E = 12,5 - 2,0 = \SI{10,5}{\micro\joule}.
\]

\textbf{Barème indicatif (sur 10) :} justification du régime (1) ; mesure de $T$ et méthode (1,5) ; calcul de $L$ avec justification de l'approximation (2) ; équation différentielle (1,5) ; bilan énergétique et relation $dE = -R_{\text{tot}}i^2 dt$ (2) ; énergies $E_0$, $E$ et $W_{\text{Joule}}$ (2).

\textbf{Points de vigilance :} ne pas oublier la résistance interne $r$ de la bobine dans $R_{\text{tot}}$ ; justifier l'approximation $T \approx T_0$ par la comparaison de $R_{\text{tot}}$ à $R_C$ (ou de $\sigma$ à $\omega_0$), c'est une condition d'application exigée ; à un extremum de $u_C$, l'intensité est nulle, ce qui autorise $E = \frac{1}{2}Cu_C^2$ ; exprimer les énergies en joules avec un nombre cohérent de chiffres significatifs.
\end{corrige}

\begin{corrige}[Exercice 11 — Type Bac : démonstration énergétique et étude complète d'un circuit RLC]
Données Partie B : $L = \SI{40}{mH}$, $C = \SI{0,47}{\micro\farad}$, $r = \SI{8,0}{\Omega}$, $U_0 = \SI{10}{V}$.

\textbf{Partie A — Circuit LC idéal.}

\textbf{1.} Énergie électrique du condensateur : $E_e = \dfrac{q^2}{2C}$. Énergie magnétique de la bobine : $E_m = \dfrac{1}{2}Li^2$.

\textbf{2.} $E = \dfrac{q^2}{2C} + \dfrac{1}{2}Li^2$. En dérivant par rapport au temps (avec $i = \dot{q}$ et $\dot{i} = \ddot{q}$) :
\[
\frac{dE}{dt} = \frac{q}{C}\,\dot{q} + L\,i\,\dot{i} = i\left(\frac{q}{C} + L\,\ddot{q}\right).
\]
Si $E$ se conserve, $\dfrac{dE}{dt} = 0$ pour tout $t$. Comme $i$ n'est pas identiquement nul, il vient $\dfrac{q}{C} + L\ddot{q} = 0$, soit $\boxed{\ddot{q} + \dfrac{q}{LC} = 0}$.

\textbf{3.} Réciproque : multiplions l'équation différentielle $\ddot{q} + \dfrac{q}{LC} = 0$ par $L\,i = L\,\dot{q}$ :
\[
L\,i\,\ddot{q} + \frac{q}{C}\,\dot{q} = 0
\quad\Longleftrightarrow\quad
\frac{d}{dt}\!\left(\frac{1}{2}Li^2\right) + \frac{d}{dt}\!\left(\frac{q^2}{2C}\right) = 0
\quad\Longleftrightarrow\quad
\frac{dE}{dt} = 0.
\]
Donc $E$ est constante : l'équation différentielle du circuit LC entraîne bien la conservation de l'énergie électromagnétique.

\textbf{Partie B — Circuit RLC réel.}

\textbf{4.} La résistance totale de la maille est $R_{\text{tot}} = R + r$. Loi d'additivité des tensions : $\dfrac{q}{C} + (R + r)\,i + L\,\dfrac{di}{dt} = 0$, soit :
\[
\boxed{L\,\ddot{q} + (R + r)\,\dot{q} + \frac{q}{C} = 0}.
\]

\textbf{5.} $R_C = 2\sqrt{\dfrac{L}{C}} = 2\sqrt{\dfrac{40\times10^{-3}}{0,47\times10^{-6}}} = 2\sqrt{8,51\times10^{4}} = 2 \times 291,7 \approx \SI{583}{\Omega}$.
Le régime pseudo-périodique exige $R_{\text{tot}} = R + r < R_C$, donc :
\[
R < R_C - r = 583 - 8,0 \approx \SI{575}{\Omega}.
\]
La valeur maximale de $R$ est donc d'environ $\SI{575}{\Omega}$ (valeur limite, non incluse : inégalité stricte).

\textbf{6.} $T_0 = 2\pi\sqrt{LC} = 2\pi\sqrt{40\times10^{-3} \times 0,47\times10^{-6}} = 2\pi\sqrt{1,88\times10^{-8}} = 2\pi \times 1,371\times10^{-4} \approx \SI{8,61e-4}{s} = \SI{0,861}{ms}$.
Avec $R = \SI{12}{\Omega}$ : $R_{\text{tot}} = 12 + 8,0 = \SI{20}{\Omega}$, donc $\sigma = \dfrac{R_{\text{tot}}}{2L} = \dfrac{20}{0,080} = \SI{250}{s^{-1}}$, alors que $\omega_0 = \dfrac{1}{\sqrt{LC}} \approx \dfrac{1}{1,371\times10^{-4}} \approx \SI{7,29e3}{rad/s}$.
Comme $\sigma \ll \omega_0$ ($250 \ll 7290$, soit $R_{\text{tot}} = \SI{20}{\Omega} \ll R_C \approx \SI{583}{\Omega}$), l'amortissement est très faible : $\omega = \sqrt{\omega_0^2 - \sigma^2} \approx \omega_0$ (écart relatif $\approx 6\times10^{-4}$), et la pseudo-période mesurée $T = \dfrac{2\pi}{\omega}$ est pratiquement égale à $T_0$.

\textbf{7.} À l'état final ($t \to +\infty$), les oscillations sont totalement amorties : $q = 0$ et $i = 0$, donc $E_{\text{final}} = 0$. Toute l'énergie initiale a été dissipée par effet Joule dans $R$ et $r$ :
\[
W_{\text{Joule}} = E_0 = \frac{1}{2}C\,U_0^2 = 0,5 \times 0,47\times10^{-6} \times 10^2 = \SI{2,35e-5}{J} = \SI{23,5}{\micro\joule}.
\]

\textbf{8.} \textbf{Application :} le circuit accordé des récepteurs radio FM utilisés dans les téléphones et autoradios (stations comme CRTV, Radio Balafon, RFI reçues à Douala ou Yaoundé). \textbf{Rôle du circuit LC :} par résonance à sa fréquence propre $f_0 = \dfrac{1}{2\pi\sqrt{LC}}$, il \emph{sélectionne} la fréquence porteuse de la station désirée parmi toutes les ondes captées par l'antenne, en répondant avec une amplitude maximale à cette seule fréquence (accord réalisé en faisant varier $C$ ou $L$).

\textbf{Barème indicatif (sur 12) :} expressions de $E_e$ et $E_m$ (1) ; dérivation de $E$ et obtention de l'équation différentielle (2) ; réciproque (1,5) ; équation différentielle RLC avec $R + r$ (1,5) ; $R_C$ et condition sur $R$ (2) ; $T_0$ et justification $T \approx T_0$ (2) ; énergie dissipée totale (1) ; application et rôle du circuit LC (1).

\textbf{Points de vigilance :} dans la démonstration, expliciter $i = \dot{q}$ et $\dot{i} = \ddot{q}$ avant de dériver ; ne pas oublier la résistance interne $r$ dans $R_{\text{tot}}$ (piège classique du sujet) ; la condition du régime pseudo-périodique est une inégalité \emph{stricte} $R + r < R_C$ ; pour la question 7, justifier que l'énergie finale est nulle avant d'écrire $W_{\text{Joule}} = E_0$ ; citer la formule $R_C = 2\sqrt{L/C}$ avant de l'appliquer.
\end{corrige}

\newpage

\coursec{Fiche bilan}

\begin{popfigure}
\begin{tikzpicture}[
  mindmap,
  grow cyclic,
  every node/.style={concept, font=\small},
  root concept/.append style={concept color=popblue, font=\small\bfseries, minimum size=2.6cm},
  level 1 concept/.append style={level distance=3.4cm, sibling angle=60, font=\scriptsize},
  text=white
]
\node[root concept]{Oscillations électriques libres LC et RLC}
  child[concept color=poporange]{ node {Bobine : $u=ri+L\dfrac{di}{dt}$ ; $\mathcal{E}_m=\tfrac12 Li^2$} }
  child[concept color=popgreen]{ node {Circuit LC : $q''+\dfrac{q}{LC}=0$} }
  child[concept color=poppurple]{ node {$T_0=2\pi\sqrt{LC}$ ; $\omega_0=\dfrac{1}{\sqrt{LC}}$} }
  child[concept color=poppink]{ node {Énergie : $\mathcal{E}=\dfrac{q^2}{2C}+\tfrac12 Li^2=$ cte} }
  child[concept color=popgold]{ node {RLC : $Lq''+Rq'+\dfrac{q}{C}=0$} }
  child[concept color=popdark]{ node {Régimes : pseudo-périodique, critique, apériodique} };
\end{tikzpicture}
\legende{Carte mentale de la leçon : les six idées forces des oscillations électriques libres.}
\end{popfigure}

\begin{bilan}
\textbf{1. Définitions clés}
\begin{itemize}
  \item \cle{Inductance} $L$ d'une bobine : coefficient (en henry, H) tel que le flux propre vaut $\Phi = Li$.
  \item \cle{Tension aux bornes d'une bobine} (convention récepteur) : $u = ri + L\dfrac{di}{dt}$ ; pour une bobine idéale ($r=0$) : $u = L\dfrac{di}{dt}$.
  \item \cle{Oscillations libres} : échanges d'énergie entre condensateur et bobine, sans apport extérieur.
  \item \cle{Amortissement} : dissipation de l'énergie par effet Joule dans la résistance totale $R$ du circuit.
\end{itemize}

\textbf{2. Formules à connaître par c\oe ur}
\begin{center}
\begin{tabularx}{\linewidth}{|l|X|l|}
\hline
\rowcolor{popblueL} \textbf{Grandeur} & \textbf{Expression} & \textbf{Unité SI} \\
\hline
Énergie magnétique (bobine) & $\mathcal{E}_m = \dfrac{1}{2}Li^2$ & J \\
\hline
Énergie électrique (condensateur) & $\mathcal{E}_e = \dfrac{q^2}{2C} = \dfrac{1}{2}Cu^2$ & J \\
\hline
Équation du circuit LC & $q'' + \dfrac{1}{LC}\,q = 0$ & --- \\
\hline
Solution (charge) & $q(t) = Q_m\cos(\omega_0 t + \varphi)$ & C \\
\hline
Intensité & $i = \dfrac{dq}{dt} = -\omega_0 Q_m\sin(\omega_0 t+\varphi)$ & A \\
\hline
Pulsation propre & $\omega_0 = \dfrac{1}{\sqrt{LC}}$ & rad/s \\
\hline
Période propre & $T_0 = 2\pi\sqrt{LC}$ & s \\
\hline
Fréquence propre & $N_0 = \dfrac{1}{T_0} = \dfrac{1}{2\pi\sqrt{LC}}$ & Hz \\
\hline
Énergie électromagnétique (LC) & $\mathcal{E} = \dfrac{q^2}{2C} + \dfrac{1}{2}Li^2 = \dfrac{Q_m^2}{2C} = \text{cte}$ & J \\
\hline
Équation du circuit RLC & $Lq'' + Rq' + \dfrac{q}{C} = 0$ & --- \\
\hline
\end{tabularx}
\end{center}

\textbf{3. Méthodes express}
\begin{itemize}
  \item \textbf{Établir l'équation différentielle} : flécher les tensions, écrire la loi des mailles $u_C + u_L = 0$ (LC) ou $u_C + u_R + u_L = 0$ (RLC), remplacer $u_C = q/C$, $u_R = Ri$, $u_L = L\,di/dt$ avec $i = dq/dt$, puis ordonner.
  \item \textbf{Identifier $\omega_0$} : mettre l'équation sous la forme $q'' + \omega_0^2 q = 0$ et lire $\omega_0^2$.
  \item \textbf{Déterminer $Q_m$ et $\varphi$} : utiliser les conditions initiales $q(0)$ et $i(0) = q'(0)$.
  \item \textbf{Exploiter la conservation de l'énergie} : quand $q = \pm Q_m$, toute l'énergie est dans $C$ ; quand $i = \pm I_m$, toute l'énergie est dans $L$, avec $I_m = \omega_0 Q_m$.
  \item \textbf{Régimes du RLC} : $R$ faible $\Rightarrow$ pseudo-périodique ($T \approx T_0$) ; $R = R_c = 2\sqrt{L/C}$ $\Rightarrow$ critique ; $R > R_c$ $\Rightarrow$ apériodique.
\end{itemize}

\textbf{4. Analogies électromécaniques}
\begin{itemize}
  \item $q \leftrightarrow x$ ; \quad $i \leftrightarrow v$ ; \quad $L \leftrightarrow m$ ; \quad $1/C \leftrightarrow k$ ; \quad $R \leftrightarrow$ frottement $\lambda$.
  \item $\dfrac{q^2}{2C} \leftrightarrow \dfrac{1}{2}kx^2$ (potentielle) ; \quad $\dfrac{1}{2}Li^2 \leftrightarrow \dfrac{1}{2}mv^2$ (cinétique).
\end{itemize}
\end{bilan}

\begin{aretenir}[Checklist avant le Bac]
\begin{itemize}
  \item $\square$ Je sais écrire la tension aux bornes d'une bobine réelle : $u = ri + L\dfrac{di}{dt}$.
  \item $\square$ Je sais établir l'équation différentielle du circuit LC par la loi des mailles.
  \item $\square$ Je sais identifier $\omega_0 = \dfrac{1}{\sqrt{LC}}$ et en déduire $T_0 = 2\pi\sqrt{LC}$ et $N_0 = \dfrac{1}{T_0}$.
  \item $\square$ Je sais écrire $q(t) = Q_m\cos(\omega_0 t + \varphi)$ et $i(t) = q'(t)$, en respectant le signe.
  \item $\square$ Je sais déterminer $Q_m$ et $\varphi$ à partir des conditions initiales.
  \item $\square$ Je sais que l'énergie électromagnétique du circuit LC idéal se conserve : $\mathcal{E} = \dfrac{Q_m^2}{2C}$.
  \item $\square$ Je sais que $q$ et $i$ oscillent en quadrature de phase ($i$ maximal quand $q$ est nul).
  \item $\square$ Je sais établir l'équation du circuit RLC : $Lq'' + Rq' + \dfrac{q}{C} = 0$.
  \item $\square$ Je sais nommer et reconnaître les trois régimes : pseudo-périodique, critique, apériodique.
  \item $\square$ Je sais que l'amortissement est dû à l'effet Joule et que $T \approx T_0$ si l'amortissement est faible.
  \item $\square$ Je vérifie toujours l'homogénéité : $2\pi\sqrt{LC}$ est bien en secondes.
  \item $\square$ Je connais les analogies électromécaniques ($L \leftrightarrow m$, $1/C \leftrightarrow k$).
\end{itemize}
\end{aretenir}

\findecours

\end{document}
